% Options for packages loaded elsewhere
\PassOptionsToPackage{unicode}{hyperref}
\PassOptionsToPackage{hyphens}{url}
\PassOptionsToPackage{dvipsnames,svgnames,x11names}{xcolor}
\documentclass[
  letterpaper,
  11pt]{article}
\usepackage{xcolor}
\usepackage[margin=0.78in]{geometry}
\usepackage{amsmath,amssymb}
\setcounter{secnumdepth}{-\maxdimen} % remove section numbering
\usepackage{iftex}
\ifPDFTeX
  \usepackage[T1]{fontenc}
  \usepackage[utf8]{inputenc}
  \usepackage{textcomp} % provide euro and other symbols
\else % if luatex or xetex
  \usepackage{unicode-math} % this also loads fontspec
  \defaultfontfeatures{Scale=MatchLowercase}
  \defaultfontfeatures[\rmfamily]{Ligatures=TeX,Scale=1}
\fi
\usepackage{lmodern}
\ifPDFTeX\else
  % xetex/luatex font selection
\fi
% Use upquote if available, for straight quotes in verbatim environments
\IfFileExists{upquote.sty}{\usepackage{upquote}}{}
\IfFileExists{microtype.sty}{% use microtype if available
  \usepackage[]{microtype}
  \UseMicrotypeSet[protrusion]{basicmath} % disable protrusion for tt fonts
}{}
\makeatletter
\@ifundefined{KOMAClassName}{% if non-KOMA class
  \IfFileExists{parskip.sty}{%
    \usepackage{parskip}
  }{% else
    \setlength{\parindent}{0pt}
    \setlength{\parskip}{6pt plus 2pt minus 1pt}}
}{% if KOMA class
  \KOMAoptions{parskip=half}}
\makeatother
\usepackage{longtable,booktabs,array}
\usepackage{calc} % for calculating minipage widths
% Correct order of tables after \paragraph or \subparagraph
\usepackage{etoolbox}
\makeatletter
\patchcmd\longtable{\par}{\if@noskipsec\mbox{}\fi\par}{}{}
\makeatother
% Allow footnotes in longtable head/foot
\IfFileExists{footnotehyper.sty}{\usepackage{footnotehyper}}{\usepackage{footnote}}
\makesavenoteenv{longtable}
\setlength{\emergencystretch}{3em} % prevent overfull lines
\providecommand{\tightlist}{%
  \setlength{\itemsep}{0pt}\setlength{\parskip}{0pt}}
% Layout for the mathematical review; no mathematical definitions are changed.
\usepackage{microtype}
\usepackage{xurl}
\usepackage{fancyhdr}
\usepackage{titlesec}
\usepackage{etoolbox}
\usepackage{hyperref}
\usepackage{ragged2e}
\usepackage{needspace}
\definecolor{reviewnavy}{HTML}{17364D}
\definecolor{reviewgray}{HTML}{52616B}
\hypersetup{colorlinks=true,linkcolor=reviewnavy,urlcolor=reviewnavy,
  pdftitle={Geometrization in Lean: mathematical reading of the skeleton},
  pdfauthor={Prepared for Bennett Chow and Peng Lu}}
\pagestyle{fancy}
\fancyhf{}
\fancyhead[L]{\footnotesize\sffamily\color{reviewgray}GEOMETRIZATION IN LEAN}
\fancyhead[R]{\footnotesize\sffamily\color{reviewgray}MATHEMATICAL REVIEW}
\fancyfoot[L]{\scriptsize\sffamily\color{reviewgray}Frozen skeleton ea0fae60 \quad / \quad Blueprint 207}
\fancyfoot[R]{\small\thepage}
\setlength{\headheight}{15pt}
\setlength{\footskip}{23pt}
\renewcommand{\headrulewidth}{0.25pt}
\setlength{\emergencystretch}{3em}
\setlength{\parskip}{5pt plus 1pt}
\setlength{\parindent}{0pt}
\renewcommand{\arraystretch}{1.16}
\titleformat{\section}{\LARGE\bfseries\color{reviewnavy}}{}{0pt}{}
\titleformat{\subsection}{\large\bfseries\color{reviewnavy}}{}{0pt}{}
\titleformat{\subsubsection}{\normalsize\bfseries\color{reviewnavy}}{}{0pt}{}
\titlespacing*{\section}{0pt}{0pt}{16pt}
\titlespacing*{\subsection}{0pt}{17pt}{8pt}
\titlespacing*{\subsubsection}{0pt}{13pt}{6pt}
\newcommand{\sectionbreak}{\clearpage}
\AtBeginEnvironment{longtable}{\small}
\AtBeginEnvironment{Shaded}{\small}
\widowpenalty=8000
\clubpenalty=8000
\displaywidowpenalty=8000
\AtBeginDocument{\hypersetup{linkcolor=reviewnavy,urlcolor=reviewnavy}\sloppy}
\usepackage{bookmark}
\IfFileExists{xurl.sty}{\usepackage{xurl}}{} % add URL line breaks if available
\urlstyle{same}
\hypersetup{
  colorlinks=true,
  linkcolor={Maroon},
  filecolor={Maroon},
  citecolor={Blue},
  urlcolor={Blue},
  pdfcreator={LaTeX via pandoc}}

\author{}
\date{}

\begin{document}

\hypersetup{pageanchor=false}
\begin{titlepage}
\thispagestyle{empty}
\vspace*{0.45in}
{\sffamily\small\color{reviewgray}BLUEPRINT 207 / FROZEN LEAN SKELETON}\par
\vspace{0.7in}
{\Huge\bfseries\color{reviewnavy}Geometrization in Lean}\par
\vspace{0.25in}
{\LARGE Mathematical reading of the skeleton}\par
\vspace{0.45in}
{\Large Review copy for Bennett Chow and Peng Lu}\par
\vspace{0.7in}
{\large 121 authored declarations\quad\textbar\quad17 admitted theorems}\par
\vspace{0.14in}
{\large Six admissions occur in the main endpoint proof.}\par
\vspace{0.6in}
This report translates the actual definitions, statements and written\par
assembly proofs. It distinguishes mathematical obligations from proved\par
conditional compositions and identifies omitted stronger conclusions.\par
\vspace{0.35in}
No Lean mathematical source was changed for this review.\par
\vfill
{\small Source commit}\par
{\footnotesize\ttfamily ea0fae60ee01ef8c8d9b57a51794c6538f679c5d}\par
\vspace{0.12in}
{\small Private branch: codex/geometrization-blueprint-skeleton-207}\par
\vspace{0.35in}
{\small Statements awaiting mathematical review; not a proof-complete formalization.}
\end{titlepage}
\setcounter{page}{1}
\hypersetup{pageanchor=true}

{
\hypersetup{linkcolor=}
\setcounter{tocdepth}{2}
\tableofcontents
}
\section{Reading guide and actual proof map}\label{overview}

\subsection{Scope, authority, and how to read this
report}\label{reading-guide}

This report describes the Lean skeleton at commit
{\small\texttt{ea0fae6\allowbreak{}0ee01ef8\allowbreak{}c8d9b57a\allowbreak{}51794c65\allowbreak{}38f679c5\allowbreak{}d}},
on the private branch
{\small\texttt{codex/\allowbreak{}geometrization-\allowbreak{}blueprint-\allowbreak{}skeleton-\allowbreak{}207}}.
The associated mathematical blueprint is revision 207. The report
translates the code as it stands; it does not silently add hypotheses,
strengthen outputs, or supply missing geometric constructions. No Lean
mathematical source is changed by this report.

The immediate review question is: \textbf{Are these the mathematical
objects and claims that we intend to formalize, and are their stated
assumptions sufficient?} A second, separate question is whether their
current decomposition gives independent formalizers sufficiently precise
tasks. A true but very large existence theorem can be acceptable for the
first question while still being too coarse for the second.

The skeleton adds 22 mathematical files and 121 explicitly authored
declarations. Of its 55 named theorems, 17 have {\small\texttt{sorry}}
proofs and 38 have written proof terms. Some of those 38 proofs use
admitted theorems. The other declarations define objects, conditions,
structures, or abbreviations. The compiler also generates constructors,
projections, and auxiliary declarations; the corresponding audit covers
542 declarations.

We use four status descriptions throughout:

\begin{longtable}[]{@{}
  >{\raggedright\arraybackslash}p{(\linewidth - 2\tabcolsep) * \real{0.5000}}
  >{\raggedright\arraybackslash}p{(\linewidth - 2\tabcolsep) * \real{0.5000}}@{}}
\toprule\noalign{}
\begin{minipage}[b]{\linewidth}\raggedright
Status
\end{minipage} & \begin{minipage}[b]{\linewidth}\raggedright
Mathematical meaning
\end{minipage} \\
\midrule\noalign{}
\endhead
\bottomrule\noalign{}
\endlastfoot
Definition & The displayed data or condition is what the Lean name
means; no existence theorem follows merely from defining it. \\
Admitted & The statement is asserted using {\small\texttt{sorry}}. It
remains a mathematical and formalization obligation. \\
Proved, conditional & The code supplies a real argument from explicit
hypotheses or admitted upstream results. It does not prove those
inputs. \\
Proved, independent of the new admissions & The code supplies a proof
using the inherited library and none of the 17 new
{\small\texttt{sorry}} statements. This is not a fresh audit of the
entire inherited library. \\
\end{longtable}

In mathematical prose, ``there exists data'' translates Lean's
{\small\texttt{Nonempty}} or existential quantifier. It does not mean
that the data is computable or canonically selected. ``Smooth'' means
\(C^\infty\), unless a finite order is explicitly written. ``Connected''
includes nonempty in the relevant Lean manifold structures. A finite set
or family may be empty unless positivity of its size is explicitly
required.

For an efficient first reading, review the exact endpoint, the six main
admissions, the order of quantifiers in the selected-flow claim, and the
written contradiction/assembly proof. Then review the detailed collapse,
hyperbolic/area, topology, and finite-regularity sections. The
declaration register provides a route back from every one of the 121
names to its translation and original file. The source snapshot and its
hashes are the authority when checking a formula or a binder.

\subsection{The proof that is actually
assembled}\label{actual-proof-map}

The following is the mathematical shape of the present endpoint proof.
It is not a claim that the analytic inputs have been proved.

\begin{enumerate}
\def\labelenumi{\arabic{enumi}.}
\tightlist
\item
  Start with an arbitrary connected, compact, boundaryless, oriented
  smooth three-manifold \(M\). Obtain an arbitrary smooth Riemannian
  metric \(g\) from the inherited metric-existence theorem.
\item
  Invoke the admitted \textbf{selected-flow theorem}, at derivative
  order \(K=20\). It supplies one surgery flow with an additional
  sequence-level late decomposition property. A merely existing raw
  surgery flow does not suffice.
\item
  Suppose that arbitrarily late nonempty regular slices fail to have
  geometrized components. Choose one such bad slice at a time \(t_j>j\)
  for every integer \(j\geq0\).
\item
  For this fixed flow and this fixed sequence, the selected-flow
  property gives one function \(A(w)\) controlling the required
  finite-order derivative tests on the collapsed cut pieces. The static
  collapse theorem is then used to choose one tolerance \(w_0\). The
  flow property gives a sufficiently late index \(N\) valid for all
  later slices in the sequence and all their connected components.
\item
  Every cut piece of the selected component at index \(N\) is either an
  actual complete finite-volume hyperbolic interior, an intrinsically
  collapsed piece with its induced normalized flow metric, or a closed
  nonnegative-curvature piece. Static recognition turns the last two
  kinds into raw graph presentations.
\item
  If one of the actual cutting tori failed to inject on fundamental
  groups, the selected-flow property would supply a nonnegative physical
  disk-area function on a future half-line with an impossible
  differential upper barrier. The admitted scalar area lemma rules this
  out. Thus the actual cutting tori are incompressible in the component
  being decomposed.
\item
  The admitted \textbf{mixed graph/hyperbolic refinement theorem}
  supplies prime factors and geometric decompositions. The written
  assembly turns those outputs into the existing geometrization
  certificate. This contradicts the choice of the bad slice.
  Consequently all sufficiently late nonempty regular slices have
  geometrized components.
\item
  The inherited reconstruction theorem uses a finite observed surgery
  history, standard discarded-factor geometries, and its recorded smooth
  identifications to recover a certificate for the original \(M\). If
  the flow has an actual empty observation, the existing
  empty-observation branch supplies the reconstruction instead. No
  finite-extinction theorem is assumed for all initial manifolds.
\end{enumerate}

Only \textbf{six} of the new admitted statements occur in the
declaration dependency closure of this proof:

\begin{longtable}[]{@{}
  >{\raggedright\arraybackslash}p{(\linewidth - 2\tabcolsep) * \real{0.5000}}
  >{\raggedright\arraybackslash}p{(\linewidth - 2\tabcolsep) * \real{0.5000}}@{}}
\toprule\noalign{}
\begin{minipage}[b]{\linewidth}\raggedright
Main admitted input
\end{minipage} & \begin{minipage}[b]{\linewidth}\raggedright
Content
\end{minipage} \\
\midrule\noalign{}
\endhead
\bottomrule\noalign{}
\endlastfoot
Selected flow & One flow with the exact late-sequence tests and common
neck accuracy described below. \\
Closed static collapse & A compact connected closed carrier satisfying
the chosen intrinsic collapse tests has a raw graph presentation. \\
Boundary static collapse & The corresponding assertion with the explicit
nearly cuspidal boundary conditions. \\
Closed nonnegative classification & A compact connected closed oriented
three-manifold with sectional curvature at least zero has a raw graph
presentation. \\
Shifted scalar area contradiction & A continuous nonnegative function
cannot have the prescribed strict upper barriers forever. \\
Mixed graph/hyperbolic refinement & Actual incompressible torus gluing
of graph and hyperbolic pieces yields the chosen prime/geometric
endpoint data. \\
\end{longtable}

The other eleven admissions are independent prepared statements:
Mostow--Prasad rigidity; a local hyperbolic atlas theorem; cusp
curvature; positivity and the infinite case of the curvature radius;
common derivative bounds for a sequence; finite-order metric pullback;
finite-order coefficient compactness; continuity of least exterior disk
area from disk comparisons; and two raw-graph prime and geometric
refinement producers. Some have proved consumers of their own. They are
\textbf{not} already being applied inside the selected-flow or mixed
refinement {\small\texttt{sorry}}. A dependency edge suggested by the
blueprint is not a Lean application until the corresponding proof body
actually uses it.

\subsection{What the successful build does and does not
establish}\label{meaning-of-build}

The recorded full {\small\texttt{Differential\allowbreak{}Geometry}}
build passed with 26,440 Lake jobs. The new declaration audit checks the
542 elaborated declarations, permits exactly the 17 registered direct
admissions, and records their dependencies. The existing baseline audit
passed for 4,311 declarations. The 155 existing GC baseline Lean
modules, the accepted toolchain/dependency pins, and the endpoint
definitions were left unchanged by the skeleton.

These checks establish that the definitions are well-formed and the
written compositions type-check with the asserted admitted statements.
They do not establish that an admitted statement is true, that every
mathematical name matches its intended standard meaning, or that every
source theorem has already been translated correctly. Those are
precisely the matters for this review. A report with a faithful
translation is evidence for review, not a substitute for your
mathematical approval.

The skeleton is the principal argument, not a node-by-node translation
of all 649 selected DAG entries. Canonical JSJ uniqueness, all relative
marking exports, every local collapse/cloud construction, and the full
analytic decomposition of the selected-flow claim are not supplied as
separate Lean interfaces here. The detailed sections identify narrower
conclusions and unexported data explicitly.

\section{The exact endpoint and the actual flow-to-endpoint
argument}\label{endpoint-flow}

\subsection{E1. The initial manifold and the
conclusion}\label{e1-endpoint}

The public theorem concerns every \textbf{connected, compact, Hausdorff,
boundaryless smooth three-manifold with a specified orientation}.
Connected includes nonempty. No metric, simple connectedness, finite
fundamental group, curvature bound, irreducibility, or primeness is
assumed. It is a smooth oriented statement. The displayed second
formulation with explicit manifold instances is the same smooth
statement; the skeleton does not add a separate Moise-based theorem
about arbitrary topological three-manifolds or a theorem about
nonorientable manifolds.

Its conclusion is existence of the following data.

\textbf{Prime stage.} There is a finite \textbf{nonempty ordered list}
\(P_1,\ldots,P_m\) of connected closed oriented smooth three-manifolds,
each prime, and an orientation-preserving diffeomorphism
\(P_1\#\cdots\#P_m\longrightarrow M\). Here the actual definition of
prime is: for every pair of connected closed oriented smooth
three-manifolds \(A,B\), if \(P_i\) is smoothly diffeomorphic to
\(A\#B\), then \(A\) or \(B\) is smoothly diffeomorphic to the standard
\(S^3\). The diffeomorphisms in this test of primeness are not required
to preserve orientation. The reconstruction of \(M\) is required to
preserve orientation.

The definition does not exclude \(S^3\) from the list and does not
require the list to be irredundant or canonical. Although the inherited
finite connected sum operation defines the sum of the empty list to be
\(S^3\), the certificate itself requires a nonempty list. It does not
assert uniqueness of the prime decomposition.

\textbf{Torus-cut stage for each \(P_i\).} There is a compact oriented
smooth three-dimensional carrier \(C_i\), either boundaryless or modeled
on the three-dimensional half-space. It has a finite \textbf{positive}
number of pairwise disjoint connected clopen pieces covering it. Every
piece's intrinsic interior is also connected and nonempty. There are
finitely many pairs of boundary tori, with exact parametrizations,
smooth matching maps, and collars. Gluing these pairs gives an actual
quotient, a smooth oriented structure on that quotient, and an
orientation-preserving diffeomorphism from it to \(P_i\). The boundary
of \(C_i\) is exactly the union of the paired boundary tori; there is no
unaccounted boundary.

Each glued seam torus is required to induce an injective homomorphism
\(\pi_1(T^2,x)\to\pi_1(P_i,f(x))\) for \textbf{every} basepoint
\(x\in T^2\), using the actual seam map followed by the actual
reconstruction. The definition supplies a signed smooth collar around
each seam. Thus ``torus cut'' here is not just an abstract graph or a
collection of fundamental groups.

The number of seam pairs may be zero. Their left and right sides may
belong to the same connected cut piece, so nonseparating cuts are
allowed. The data record which piece owns each side. There is no
canonical-JSJ, minimality, nonparallelness, or uniqueness requirement.
In particular, do not read ``essential tori'' as silently including any
additional condition beyond the actual embedded collared seams and the
stated fundamental-group injections.

\textbf{Metric stage.} On the entire intrinsic interior of every
connected cut piece there is a complete smooth Riemannian metric locally
isometric to one of the eight fixed models described next. Hyperbolic
interiors must have finite total Riemannian volume. No finite-volume
condition is imposed on the other seven model types. The metric need not
extend over the boundary of the compact cut carrier. No metric matching
across the reglued tori is required; these are complete metrics on the
separate interiors.

These are the actual fields of
{\small\texttt{Geometrization\allowbreak{}Certificate}} and
{\small\texttt{Geometric\allowbreak{}Decomposition}}, unfolded through
{\small\texttt{Prime\allowbreak{}Decomposition}},
{\small\texttt{Compact\allowbreak{}Carrier.\allowbreak{}Components}},
{\small\texttt{Torus\allowbreak{}Gluing}}, and
{\small\texttt{Smooth\allowbreak{}Assembly}}. These definitions are
inherited and unchanged by the new skeleton.

\subsection{E2. What ``one of the eight geometries'' literally
means}\label{e2-models}

A geometric structure contains a model label, a smooth positive definite
metric \(h\), completeness for its intrinsic Riemannian extended
distance, and a metric-preserving local smooth atlas from the designated
model. On the connected interiors in the endpoint this is the usual
intrinsic metric completeness condition. It is \textbf{not} a bare
geometry label, a hypothesis that an arbitrary homogeneous metric
exists, or a metric on an unrelated abstract manifold.

For the spherical, Euclidean, and spherical-product labels, the fixed
metrics are respectively the round unit \(S^3\), Euclidean
\(\mathbb R^3\), and the product of the round unit \(S^2\) with
Euclidean \(\mathbb R\). For the other five labels, the code uses the
following coframes on \(\mathbb R^3\) with coordinates \((x,y,z)\). The
model inner product is the sum of the squares of the three displayed
one-forms.

\begin{longtable}[]{@{}
  >{\raggedright\arraybackslash}p{(\linewidth - 2\tabcolsep) * \real{0.5000}}
  >{\raggedright\arraybackslash}p{(\linewidth - 2\tabcolsep) * \real{0.5000}}@{}}
\toprule\noalign{}
\begin{minipage}[b]{\linewidth}\raggedright
Label
\end{minipage} & \begin{minipage}[b]{\linewidth}\raggedright
Fixed coframe
\end{minipage} \\
\midrule\noalign{}
\endhead
\bottomrule\noalign{}
\endlastfoot
Hyperbolic & \(e^{-z}dx,\ e^{-z}dy,\ dz\) \\
Hyperbolic product & \(e^{-y}dx,\ dy,\ dz\) \\
Universal \(\widetilde{\mathrm{SL}}_2\) &
\(e^{-y}dx,\ dy,\ dz+e^{-y}dx\) \\
Nil & \(dx,\ dy,\ dz-x\,dy\) \\
Sol & \(e^zdx,\ e^{-z}dy,\ dz\) \\
\end{longtable}

More explicitly, at every point of a geometric interior there must be a
smooth partial diffeomorphism from an open subset of the selected model
whose target contains that point and which pulls the interior metric
back to the displayed model metric at every source point and on all
tangent vectors. The atlas is not required to be orientation-preserving.
There is no explicit discrete group, developing map, holonomy
representation, or maximal-isometry-group action in this endpoint
record. The identification of these fixed models with the standard eight
geometries and the usual global quotient interpretation are matters of
the inherited model theory, not extra fields that can be assumed present
in the certificate.

The code also defines local and global homogeneity predicates elsewhere,
but neither is an additional field of
{\small\texttt{Geometric\allowbreak{}Structure}}; the fixed local-model
atlas is its actual requirement. The finite-volume field is the
implication ``if the label is hyperbolic, then the total volume is
finite,'' not a blanket finite-volume requirement.

\textbf{Source dictionary.} The exact inherited definitions are in
{\small\texttt{Geometrization/\allowbreak{}Prime.\allowbreak{}lean}},
{\small\texttt{Carrier.\allowbreak{}lean}},
{\small\texttt{Statement.\allowbreak{}lean}},
{\small\texttt{Torus\allowbreak{}Gluing.\allowbreak{}lean}},
{\small\texttt{Smooth\allowbreak{}Torus\allowbreak{}Reconstruction.\allowbreak{}lean}},
and
{\small\texttt{Geometry/\allowbreak{}Thurston/\allowbreak{}\{Atlas,Models,Model\allowbreak{}Atlas,Elementary\allowbreak{}Models\}.\allowbreak{}lean}}.
The review snapshot includes exact paths and source hashes.

\subsection{F1. The raw flow and its actual
observations}\label{f1-raw-flow}

The selected-flow claim uses an existing type
{\small\texttt{Raw\allowbreak{}Surgery(P,g)}}. Here \(P\) is a compact
Hausdorff oriented smooth three-dimensional stage, possibly disconnected
or empty, and \(g\) is a smooth Riemannian metric on that actual stage.
This type already contains considerably more than a sequence of
unrelated manifolds:

\begin{itemize}
\tightlist
\item
  For each integer \(n\geq0\), an actual finite retained-core surgery
  history through time \(n\), with a marked identification of its
  initial stage and initial metric with \((P,g)\).
\item
  Strictly increasing event times beginning after time zero, incoming
  Ricci-flow slabs, output metrics, actual cut-and-cap transitions and
  retained-core maps, and a final regular slab whenever its interval has
  positive length. The flow equation and the stipulated time/spatial
  smoothness on these slabs are fields of the inherited history type.
\item
  Compatibility under restriction from \([0,n+1]\) to \([0,n]\),
  including the initial identification. These are coherent histories of
  one flow, not fresh choices for unrelated finite intervals.
\item
  At every event: a singular incoming endpoint, the recorded
  boundary-frame reversal, and the inherited standard classification of
  discarded components. These are the precise three clauses of
  {\small\texttt{History\allowbreak{}Event\allowbreak{}Control}}.
\end{itemize}

Observation at a real time \(b\geq0\) means: take the history through
\(\lceil b\rceil\) and restrict it to \([0,b]\). The existing
compatibility theorems identify this with any longer observation
restricted to the same interval. Event times are locally finite. At an
event time the active stage is the post-event stage. Empty stages remain
empty in later observations.

The inherited library already constructs a
{\small\texttt{Raw\allowbreak{}Surgery(P,g)}}. The new admitted
selected-flow theorem asks for a flow with additional common accuracy
and late-geometric properties; it does \textbf{not} claim that every raw
flow, or the previously chosen arbitrary raw flow, has those properties.
No global canonical-neighborhood/noncollapse/parameter profile beyond
the fields actually stated below should be silently attached to the
words ``raw flow.''

\subsection{F2. Regular slices, normalizations, and finite common-time
selection}\label{f2-slices}

\textbf{Definition: {\small\texttt{Regular\allowbreak{}Slice}}.} For an
observation tower, a regular slice is a real time \(t>0\) which is not
an event time, together with the condition that the start time of its
last observed stage is strictly smaller than \(t\). Its history is
exactly the observation through \(t\); its stage is the last stage of
that history; and its metric is that stage's actual physical metric at
\(t\). A regular slice is allowed to be empty.

The accompanying definitions {\small\texttt{history}},
{\small\texttt{stage}}, {\small\texttt{metric}}, and
{\small\texttt{initial}} are exactly these projections and the inherited
initial identification. The two new metric definitions are \[
 \bar g(t)=t^{-1}g(t),\qquad \widehat g(t)=(4t)^{-1}g(t).
\] They are called {\small\texttt{normalized\allowbreak{}Metric}} and
{\small\texttt{curvature\allowbreak{}One\allowbreak{}Metric}}
respectively. The name of the latter is not itself an assertion that the
metric has curvature \(-1\). It is the normalization suited to a
curvature \(-1/4\) hyperbolic limit for \(\bar g\). For a connected
component \(C\) of the actual slice,
{\small\texttt{component\allowbreak{}Metric}} is precisely the
restriction of \(\bar g(t)\) to that component's open subtype, not a
newly chosen metric.

\textbf{Definition: arbitrarily late nonempty slices.} For every real
\(B\), there is a regular slice at a time \(t>B\) with nonempty stage.
There is no assertion that all towers have this property.

\textbf{Proved helper: regular times exist after every bound.} For every
tower and every real \(B\), there is a regular slice later than \(B\).
This does not assert that its stage is nonempty.

\textbf{Proved helper: empty slice or arbitrarily late nonempty slices.}
Every tower has an empty regular slice, or has arbitrarily late nonempty
regular slices. The displayed theorem is a disjunction, not an assertion
that a nonempty late component must survive. It is proved by negating
the late-nonempty property and then choosing a sufficiently late regular
time.

\textbf{Proved helper: a finite common late slice.} Suppose a tower has
arbitrarily late nonempty regular slices. Given a finite set \(I\) of
conditions \(q_i\) on slices, assume for every \(i\in I\) that some
bound \(B_i\) makes \(q_i\) true on \textbf{every} regular slice after
\(B_i\). Then, past any additionally prescribed bound \(B\), one can
choose a nonempty regular slice satisfying every \(q_i\). The proof
takes a finite maximum. The set \(I\) may be empty; the conclusion still
requires a nonempty slice later than \(B\). This does not intersect
arbitrary subsequences and does not take a maximum over infinitely many
derivative orders or thickness parameters.

\textbf{Proved reconstruction adapter.} If one particular
{\small\texttt{Raw\allowbreak{}Surgery}} from \((M,g)\) has a bound
\(B\) after which every nonempty regular slice has a geometrization
certificate on each connected component, then \(M\) geometrizes. This is
{\small\texttt{geometrizes\_\allowbreak{}of\_\allowbreak{}late\_\allowbreak{}slice\_\allowbreak{}supply}}.
Its proof passes the supplied certificates to the existing
reconstruction theorem for the \textbf{same} tower. The
empty-observation case is handled by that inherited theorem; nonempty
late slices need not be assumed separately.

The preceding selection and reconstruction adapters contain no new
admission of their own and do not need any of the 17 new admitted claims
when their stated hypotheses are supplied.

\subsection{F3. The actual cut carrier and its
metric}\label{f3-induced-cut-metric}

Let \(D\) be a torus decomposition of an actual connected closed
oriented component \(M\), in the precise sense spelled out in the
topology section: an actual compact cut carrier, its connected pieces,
paired boundary tori, a smooth quotient atlas, and a reconstruction to
\(M\). At this point neither incompressibility nor primeness is required
in \(D\).

For its \(i\)th cut piece \(C_i\), define \[
 q_i:C_i\longrightarrow M
\] to be inclusion into the cut carrier, followed by the torus quotient
map, followed by the reconstruction. This is
{\small\texttt{cut\allowbreak{}Piece\allowbreak{}Map}}.

For a given metric \(g\) on \(M\),
{\small\texttt{is\allowbreak{}Induced\allowbreak{}Cut\allowbreak{}Metric(g,D,i,h)}}
is exactly \[
 h_x(v,w)=g_{q_i(x)}\bigl(dq_i(v),dq_i(w)\bigr)
 \quad\text{for all }x\in C_i\text{ and all }v,w\in T_xC_i.
\] The equality includes boundary points and uses the manifold
derivative. The cut piece metric \(h\) is itself required to be a smooth
Riemannian metric. A favorable metric on an abstract diffeomorphic model
cannot be substituted without this equality.

\textbf{Definition:
{\small\texttt{Hyperbolic\allowbreak{}Or\allowbreak{}Collapsed}}.} For
given \(g,D,K,A,w_0\) and \(i\), the data have exactly one of the
following constructor forms. The type does not assert that the three
cases are mutually exclusive.

\begin{enumerate}
\def\labelenumi{\arabic{enumi}.}
\tightlist
\item
  A geometric structure on the \textbf{whole intrinsic interior} of
  \(C_i\) whose model label is hyperbolic. Completeness and finite
  volume are part of that structure. This hyperbolic metric is not
  required to equal the restriction of the finite-time flow metric
  \(g\).
\item
  A smooth metric \(h\) on \(C_i\), the equality \(h=q_i^*g\) above, and
  the selected static-collapse hypotheses with parameters \(K,A,w_0\).
\item
  A smooth metric \(h=q_i^*g\) on \(C_i\), empty boundary, and sectional
  curvature at least zero everywhere.
\end{enumerate}

Case 3 is explicitly separate: the assembly does not need to obtain a
finite curvature radius for a globally nonnegative closed component. In
case 2, balls, boundary distances and volume tests use the intrinsic
geometry of \(C_i\) with \(h\); they are not automatically ambient balls
in \(M\). No comparison with ambient balls is smuggled into the
definition.

\textbf{Proved conditional assembly.} Suppose every connected compact
carrier satisfying the static tests for the given \(K,A,w_0\) has a raw
graph presentation. Suppose also that the actual seam maps of \(D\) are
fundamental-group injective and that every cut piece has one of the
three forms just listed. Then \(M\) geometrizes. The proof invokes the
supplied static theorem in case 2, the admitted nonnegative
classification in case 3, and the admitted mixed graph/hyperbolic
refinement for the resulting collection. It does not prove any of those
admitted producers. The written topological assembly explicitly discards
the induced-metric equalities. Those equalities constrain what the
geometric producer must supply; this assembly does not use them or
require them of the final geometric metrics.

\subsection{F4. What common neck accuracy actually
requires}\label{f4-common-accuracy}

For a fixed flow \(F\) and function \(\delta:\mathbb R\to\mathbb R\),
{\small\texttt{has\allowbreak{}Common\allowbreak{}Neck\allowbreak{}Accuracy(F,delta)}}
says: \[
 \forall n\in\mathbb N\quad\exists p_n\quad
 \bigl(p_n.\delta=\delta\bigr)\ \text{and}\
 \bigl(\forall\text{ events }i\text{ in }F|_{[0,n]},\
       \exists\text{ a GeometricCutoffRecord for }i\text{ with }p_n\bigr).
\] The equality of the delta functions is on \textbf{all real
arguments}, not just the event times. The other fields of \(p_n\) may
depend on \(n\). In particular the definition does not provide one
global choice of the neck-radius and protected-radius functions, cap
scaffold, model radius, model order, model accuracy, or recentering
constant. Nor are records chosen for different prefixes required to
agree as data. The predicate is a precise common-delta condition, not
the full global analytic parameter register.

The parameter object contains positive neck-radius and protected-radius
functions on nonnegative times, \(0<\delta(t)<1\) there, a cap scaffold,
positive model radius and model accuracy, a natural-number model order,
and a recentering constant at least \(4\).

A cutoff record is an inherited, concrete geometric condition. To
prevent the abbreviation from hiding assumptions, its fields are
summarized here:

\begin{itemize}
\tightlist
\item
  The incoming event is singular. If the cut index set is nonempty there
  is a positive common nominal radius \(r\), satisfying
  \(r<\delta(t)^2\,r_{\rm neck}(t)\) and \(r^2\leq t\).
\item
  Each cut has \(0<\delta_\alpha\leq\delta(t)\) and a normalized neck of
  integer order at least
  \(\max\{m+6,\,2\lfloor\delta_\alpha^{-1}\rfloor_++4\}\). Its metric
  scaling factor is \(r^{-2}\). The neck buffers are pairwise disjoint,
  contain the actual tube domains, and their charts agree with the
  actual surgery tube maps. The incoming backward-neck condition is
  recorded for these necks and this radius.
\item
  The retained set lies in the incoming terminal regular region. Points
  with scalar curvature at most \(r_{\rm protected}(t)^{-2}\) lie in the
  interior of the actual retained set, and every retained connected core
  component meets that protected region.
\item
  Exactly one side of each cut is retained. If there are no cuts, some
  point of the pre-event core is discarded. Retained boundary components
  have presented static caps with the chosen scaffold and model
  parameters.
\item
  The recentered cap neck has scaling factor equal to the terminal
  scalar curvature at its center, the specified sphere marking, and
  accuracy \(c\delta_\alpha\). Its relative scaling error is at most
  \(c\delta_\alpha\). Its chart equals the original neck chart after the
  specified side-dependent change of axial coordinate
  \(s\mapsto\pm(1+s)\), and the transformed domain lies in the original
  buffer.
\item
  The event's old region is exactly the retained core. Every fixed
  Hamilton--Ivey region indexed by \(a>0\) that contains the entire
  terminal regular metric is preserved on the output. Every scalar lower
  bound \(L\leq0\) valid on that entire regular region is also preserved
  on the output.
\end{itemize}

The exact incoming-backward-neck and presented-static-cap notions are
inherited geometric structures, not newly admitted predicates defined as
{\small\texttt{True}}. Their detailed cap/neck fields remain part of the
inherited PC foundation boundary of this report. This field-by-field
summary is not a new mathematical source audit of those inherited
structures.

\subsection{F5. The physical exterior-area
obstruction}\label{f5-area-obstruction}

For a flow \(F\), any connected closed oriented manifold \(M\) with a
torus decomposition \(D\), and any real \(t_0\), the predicate
{\small\texttt{has\allowbreak{}Exterior\allowbreak{}Area\allowbreak{}Obstruction\allowbreak{}After(F,D,t0)}}
says the following for \textbf{each seam index \(i\) and each torus
basepoint \(x\)}. This predicate alone does not identify \(M\) as a flow
component or \(t_0\) as its observation time. In F6, the surrounding
late-sequence predicate supplies those exact identifications.

If the actual reconstructed seam map fails to inject on fundamental
groups at \(x\), then there exist numbers \(T,c\), subsets \(W(t)\) of
the actual post-event stage \(M_t\) of this \textbf{same} flow, and
parametrized continuous loops \(\gamma_t:S^1\to M_t\) for \(t\geq T\),
such that \[
 T\geq0,\qquad c>0,\qquad T\geq t_0.
\] Let \(a(t)\), for \(t\geq T\), be the infimum of the physical areas
of the admissible embedded exterior disks in \(W(t)\) with boundary
trace exactly \(\gamma_t\). The metric is the actual unnormalized
post-event metric \(g_F(t)\). The admissible disk class and the infimum
convention are unfolded in the hyperbolic/area section. The predicate
then requires:

\begin{enumerate}
\def\labelenumi{\arabic{enumi}.}
\tightlist
\item
  At every \(t\geq T\) an admissible disk exists and attains this
  infimum.
\item
  The real function \(a\) is continuous on \([T,\infty)\).
\item
  At every \(t\geq T\) there is a smooth upper test function touching
  \(a\) at \(t\), dominating it locally on the half-line, whose
  derivative at contact is strictly less than
  \[\frac{3a(t)}{4(t+c)}-\pi.\]
\end{enumerate}

These are conditions at every real time on the half-line, including
surgery times. All \(T,c,W,\gamma\) may depend on the seam and basepoint
under consideration. No uniform bound on these choices over the whole
sequence of slices is stated. The loops are not required before \(T\);
all-time definitions use a specified zero extension of the scalar area
before \(T\). The use of a post-event carrier at a surgery time is
explicit.

\textbf{What is not a field.} The predicate does not identify
\(\gamma_t\) with a primitive meridian or a transported kernel class of
the failed seam. It does not give a time-continuous family of loops, a
persistent cusp embedding, a specified topological exterior, a common
unchanged compact carrier across surgery, or maps transporting disks
through that carrier. The sets \(W(t)\) are arbitrary subsets subject to
the disk conditions. Stability, minimal surface equations and boundary
estimates are not separate exported fields. The coarse producer is
expected to construct the displayed consequence using that geometric
mathematics; it has not yet been split into those interfaces.

This limitation has an exact logical expression. Given the scalar area
contradiction and the proved nonnegativity of the infimum, the
obstruction predicate implies incompressibility. Conversely, if the
seams are already incompressible, every failed-injection premise is
false, so the obstruction predicate holds vacuously. Thus this is a
geometric-looking \emph{conditional obstruction consequence}, logically
equivalent to incompressibility once the scalar lemma is available. It
is not an independent construction of least-area disks merely because
the quantified data appear in its statement. This equivalence is
mathematical analysis of the definitions in this report; the skeleton
does not contain a separately named equivalence theorem.

\subsection{F6. The late-sequence tests: full quantifier
order}\label{f6-sequence-tests}

Fix one flow \(F\) and one natural number \(K\). Put
\(\omega_3=4\pi/3\). The predicate
{\small\texttt{has\allowbreak{}Late\allowbreak{}Sequence\allowbreak{}Tests(F,K)}}
has the following order: \[
 \begin{gathered}
 \forall (s_j)_{j\geq0}\ \text{of actual regular slices of }F,\\
 [\ t_j>j\text{ for every }j\ ]\ \land\
 [\ M_{t_j}\ne\varnothing\text{ for every }j\ ]\quad\Longrightarrow\\
 \exists A:\mathbb R\to\mathbb R,
 \quad A(w)>0\quad(0<w<\omega_3),\\
 \forall w_0\in(0,\omega_3)\quad\exists N\in\mathbb N\quad
 \forall j\geq N\quad\forall\text{ connected components }C\subset M_{t_j},\\
 \exists\text{ a torus decomposition }D\text{ of that actual }C,\\
 [\ \text{every piece has }\mathrm{HyperbolicOrCollapsed}
       (t_j^{-1}g_F(t_j)|_C,D,K,A,w_0)\ ]\\
 \land\ [\ \mathrm{hasExteriorAreaObstructionAfter}(F,D,t_j)\ ].
 \end{gathered}
\]

The sequence of times need not be increasing; \(t_j>j\) suffices. The
choice of \(A\) may depend on the entire sequence and the already fixed
\(F,K\). It is one function for all components and all tested late
indices in that sequence. It has no continuity or monotonicity
hypothesis. The choice of \(N\) may depend on the sequence and \(w_0\),
but is uniform in \(j\geq N\) and in the component \(C\). The
decomposition may depend on \(w_0,j,C\); it is not required to be
nested, isotopic, or consistent between different slices or different
tolerances. The area-obstruction data may depend on this chosen
decomposition.

In particular, this is not a theorem of the form ``one function \(A\)
works for every late time on this flow.'' Nor does it select a new flow
after seeing \(w_0\). It supplies exactly the sequence-level uniformity
used in the contradiction proof. The universal quantifier over sequences
is vacuous if the flow has no sequence of arbitrarily late nonempty
slices; that is compatible with the separate actual-empty reconstruction
branch.

\subsection{F7. The admitted selected-flow
theorem}\label{f7-selected-flow}

\textbf{Admitted statement.} For every compact oriented smooth
three-stage \(P\), every smooth Riemannian metric \(g\) on \(P\), and
every integer \(K\geq20\), there exist a real function \(\delta\) and
one raw surgery flow \(F\) starting from that same \((P,g)\) such that:

\begin{enumerate}
\def\labelenumi{\arabic{enumi}.}
\tightlist
\item
  \(0<\delta(t)<1\) for every \(t\geq0\).
\item
  \(\delta\) is nonincreasing on \([0,\infty)\).
\item
  For every \(\varepsilon>0\) there is \(B\in\mathbb R\) such that
  \(\delta(t)<\varepsilon\) for every real \(t>B\).
\item
  \(F\) has common neck accuracy \(\delta\) in the precise prefix sense
  of F4.
\item
  \(F\) has the late-sequence tests of F6 at order \(K\).
\end{enumerate}

The positivity, monotonicity and eventual smallness are separate
clauses; monotonicity alone would not imply decay to zero. There is no
assumed normalization of \(g\), and no positivity or smallness
assumption on its curvature. The selected flow and \(\delta\) may depend
on \(P,g,K\). This does not assert a single flow working simultaneously
for every \(K\). There are no requirements on \(\delta\) at negative
times apart from the literal all-real function equalities in the prefix
condition and the eventual inequality just written.

The exact declaration is
{\small\texttt{GC.\allowbreak{}Long\allowbreak{}Time.\allowbreak{}exists\_\allowbreak{}surgery\_\allowbreak{}with\_\allowbreak{}late\_\allowbreak{}sequence\_\allowbreak{}tests}}.
This is the largest new admitted composite. The claim does not take
{\small\texttt{Geometrizes}} as a hypothesis, but it incorporates most
of the still unseparated long-time geometric work in its conclusion. The
extra common delta/profile clauses are \emph{not used} in the current
capstone after this theorem has been invoked; the proof projects only
the selected flow and its late-sequence tests. This is a legitimate use
of a stronger conclusion, not evidence that the unused analytic
conditions have been derived elsewhere.

\subsection{F8. The written bad-sequence
proof}\label{f8-bad-sequence-proof}

\textbf{Proved, conditional statement.} If a fixed raw flow \(F\)
satisfies the late-sequence tests for an integer \(K\geq10\), then some
real \(B\) has the following property: every nonempty regular slice of
this same flow at a time greater than \(B\) has a geometrization
certificate on each actual connected component.

Here is the actual proof, including the order in which its inputs are
used.

\begin{enumerate}
\def\labelenumi{\arabic{enumi}.}
\tightlist
\item
  Negate the conclusion. For each integer \(j\geq0\), choose a nonempty
  regular slice \(s_j\) at time \(t_j>j\) for which not every component
  geometrizes. The code uses classical choice. It does not assume this
  sequence is monotone in time.
\item
  Apply the given late-sequence tests to obtain \(A\). Apply the
  admitted closed and boundary static theorems to this \(K,A\). Their
  proved minimum-threshold combination gives one \(w_0\in(0,\omega_3)\)
  valid for all connected compact carriers satisfying the appropriate
  tests.
\item
  Apply the sequence tests at this \(w_0\) to obtain \(N\). Fix any
  actual connected component \(C\) of the slice \(s_N\), and obtain its
  decomposition \(D\) and its three-way piece data.
\item
  To prove incompressibility, suppose one actual seam map fails to
  inject. Its area obstruction supplies \(T,c,a\). The scalar area is
  nonnegative by the definition and proved nonnegativity of the
  disk-area infimum. It is continuous and has the displayed strict upper
  barriers. Since \(T\geq0\) and \(c>0\), the denominator has the
  required positivity. The admitted shifted area lemma contradicts these
  properties.
\item
  Apply the conditional assembly in F3, using the static threshold from
  step 2, the now established seam injections, and the actual piece
  data. Thus \(C\) geometrizes. Since \(C\) was arbitrary, all
  components of \(s_N\) geometrize, contradicting the choice of \(s_N\).
\end{enumerate}

The proof uses the obstruction's continuity and barriers but does not
use its separate attainment field or its bound \(T\geq t_N\). Those
stronger clauses remain part of the admitted geometric output. It also
does not use the separate common-derivative-bound lemma, Mostow--Prasad,
local compactness, or cusp-family transport inside this proof; such
inputs belong under the admitted producer when that proof is eventually
written.

This theorem is
{\small\texttt{components\_\allowbreak{}geometrize\_\allowbreak{}of\_\allowbreak{}late\_\allowbreak{}sequence\_\allowbreak{}tests}}.
The regular slice, component, metric and decomposition in every
application refer to the same chosen flow and the same selected index.
There is no unrecorded change of carrier in the composition.

\subsection{F9. From the late flow to the original
manifold}\label{f9-capstone}

\textbf{Proved, conditional theorem with a supplied initial metric.} For
every connected closed oriented smooth \(M\) and every smooth metric
\(g\) on its corresponding stage, \(M\) geometrizes. The actual proof of
{\small\texttt{geometrizes\_\allowbreak{}of\_\allowbreak{}metric}}
invokes F7 at \(K=20\), applies F8 using \(20\geq10\), and then applies
the already proved same-flow reconstruction adapter in F2. Its status is
conditional on the admitted producers, even though its statement has no
explicit analytic hypothesis: the proof references the globally asserted
{\small\texttt{sorry}} theorems.

The inherited reconstruction runs backward through one finite
observation history. At each event it uses the actual smooth cut-and-cap
completion, geometrization of the classified discarded components, and
the standard \(S^2\times S^1\) factor for graph cycles. It keeps track
of the initial oriented diffeomorphism. If the final observed manifold
is empty, the componentwise endpoint assertion there is vacuous, and the
same history reconstruction still supplies the original certificate.
Neither this argument nor F8 asserts that every flow becomes extinct.

\textbf{Public capstone and its two reformulations.}
{\small\texttt{GC.\allowbreak{}Endpoint.\allowbreak{}geometrization}}
obtains a smooth metric on the compact \(M\) and applies
{\small\texttt{geometrizes\_\allowbreak{}of\_\allowbreak{}metric}}.
{\small\texttt{geometrization\_\allowbreak{}conjecture}} is the same
theorem viewed as a universally quantified proposition.
{\small\texttt{smooth\_\allowbreak{}geometrization\_\allowbreak{}conjecture}}
is its equivalent formulation with the topology, smooth atlas,
compactness, connectedness and orientation supplied as explicit
arguments. These are three names for the same endpoint conclusion in the
smooth oriented category. They are not three separate proofs and do not
remove any admission.

\subsection{F10. Review decisions for the
integration}\label{f10-review-decisions}

Five decisions deserve explicit review before independent proof tasks
are assigned. The joint-review worksheet returns to these with source
references.

\begin{enumerate}
\def\labelenumi{\arabic{enumi}.}
\tightlist
\item
  \textbf{Endpoint:} Approve E1's noncanonical prime and torus data,
  E2's model normalizations, and hyperbolic-only finite volume.
  Topological and nonorientable statements need separate formulations.
\item
  \textbf{Uniformity:} Approve F6's sequence-level quantifiers. They
  suffice for the contradiction but supply neither an all-times
  derivative function nor persistent geometric data between slices.
\item
  \textbf{Flow parameters:} Decide whether the common-delta condition is
  a sufficient exported interface. Other cutoff parameters may vary with
  the finite prefix.
\item
  \textbf{Incompressibility:} Meridians, persistent cusp families and
  surgery transports remain inside the admitted producer. The scalar
  contradiction does not supply those independent interfaces.
\item
  \textbf{Relative reconstruction:} The mixed producer can choose new
  primes, seams and metrics. Its conclusion, expanded in T9, governs
  which incoming markings are retained.
\end{enumerate}

These scope distinctions do not assert that the statements or written
implications are false. They identify precisely what approval would
cover.

\section{Static collapse: faithful mathematical reading of the Lean
skeleton}\label{static-collapse-faithful-mathematical-reading-of-the-lean-skeleton}

\subsection{C0. Scope, conventions, and proof
status}\label{collapse-scope}

This section translates \textbf{all 35 authored declarations in the
seven collapse-related modules} at commit
{\small\texttt{ea0fae6\allowbreak{}0ee01ef8\allowbreak{}c8d9b57a\allowbreak{}51794c65\allowbreak{}38f679c5\allowbreak{}d}}.
It is a description of the actual Lean code, not a proposed improvement
or a claim that all collapse nodes in the blueprint have been
formalized. No Lean code was changed in preparing this report.

There are \textbf{seven admitted theorem bodies}, \textbf{four
independently proved elementary monotonicity theorems}, and \textbf{two
composition theorems with genuine Lean proofs that use admitted collapse
theorems}. The other 22 declarations are definitions, structures, or
abbreviations, including two explicitly constructed weakening
operations. ``Independently proved'' here means their axiom audit does
not contain {\small\texttt{sorry\allowbreak{}Ax}}; it does not mean
their entire inherited mathematical foundation was re-audited for this
report. A theorem with an actual proof can still depend on admissions
through the theorems it invokes.

In the three-dimensional statements, a \textbf{compact carrier} means an
oriented compact Hausdorff second-countable smooth 3-manifold, modeled
either on Euclidean 3-space or the standard half-space. Its orientation
is part of the supplied data. It need not be connected or nonempty
merely by being a compact carrier. Whenever the collapse theorems assume
``connected,'' the Lean {\small\texttt{Connected\allowbreak{}Space}}
assumption includes nonemptiness. There are smooth boundaries here, not
arbitrary manifolds with corners. The notation \(\partial W\) denotes
the boundary determined by the actual manifold model, not a selected
subset or a boundary imported from another carrier.

Every metric \(g\) in these seven modules is a \textbf{smooth
positive-definite Riemannian metric}. Finite regularity enters through
cusp embeddings, not through nonsmooth input metrics to the static
collapse theorem.

Write

\[
 d_g(x,y)\in[0,\infty],\qquad B_g(p,r)=\{q:d_g(p,q)<r\},\qquad
 V_g(p,r)=\operatorname{Vol}_g B_g(p,r),\qquad \omega_3=4\pi/3.
\]

The distance is the intrinsic Riemannian distance \textbf{within the
specified carrier}, defined using lengths of curves in that carrier.
Thus a ball in a cut piece is not automatically the intersection of that
piece with a ball in a larger flow manifold. Volume is the actual
Riemannian volume measure on the same carrier. The code uses extended
nonnegative reals for distance and volume. For real \(r\leq0\), its
open-ball definition gives the empty set; the substantive tests below
always impose \(r>0\).

For the dimension-independent auxiliary definitions, the setting is a
smooth manifold with a finite-dimensional real normed model, a model
with corners, and the indicated Hausdorff and sigma-compact assumptions.
They do not silently assume dimension three, compactness, orientation,
or connectedness. The constant \(\omega_3\) and the exponent three are
nevertheless explicitly three-dimensional in the collapse predicates.

The convention for sectional curvature is fixed by the identity

\[
 \sec_g\geq a\quad\Longleftrightarrow\quad
 a\bigl(g(v,v)g(w,w)-g(v,w)^2\bigr)\leq \mathrm{Rm}_g(v,w,w,v)
\]

for every point and every pair of tangent vectors. Dependent and zero
vectors are included; division by a possibly zero Gram determinant is
avoided.

\subsection{C1. What the curvature-derivative notation actually
means}\label{collapse-tensor}

\textbf{Definitions, with no admitted proofs.} For a raw covariant
\(s\)-tensor field \(T\), the code defines \(\nabla_g T\) by the usual
coordinate expression for the Levi--Civita covariant derivative:
differentiate the coordinate tensor and subtract the connection action
in each covariant slot. The differentiation slot is prepended to the
existing tensor slots. At a boundary point the derivative is taken
within the range of the manifold's model with corners. The result is
converted from coordinates back to a tensor in the actual tangent space.

The iteration is defined recursively:

\[
 \nabla_g^0T=T,\qquad \nabla_g^{k+1}T=\nabla_g(\nabla_g^kT).
\]

The curvature quantity used everywhere below is

\[
 D_k(g,x):=\left|\nabla_g^k\mathrm{Rm}_g(x)\right|_g,
\]

where \(\mathrm{Rm}_g\) is the actual covariant four-tensor of the
Levi--Civita connection and the norm is the induced tensor
Hilbert--Schmidt norm, not the operator norm or the coordinate sup norm.
The norm is the square root of the corresponding tensor inner product.

\textbf{Total-definition qualification.} Lean's derivative operation is
defined even for functions that are not differentiable. Consequently the
raw operation \(T\mapsto\nabla_gT\) accepts an arbitrary tensor field
without a regularity premise. For differentiable tensor fields its
displayed formula is the ordinary covariant derivative. For arbitrary
nonsmooth fields one must not read its value as asserting that a
classical covariant derivative exists. Here the curvature field comes
from a smooth metric; the cusp-error application below has the explicit
\(C^{K+1}\) map regularity that is intended to justify the derivatives
through order \(K\). The new modules do not separately prove a general
finite-regularity compatibility theorem for this raw operation. This is
an implementation obligation, not an extra regularity theorem already
obtained from the definition.

\textbf{Lean references:}
{\small\texttt{metric\allowbreak{}Covariant\allowbreak{}Derivative}}
(Metric.lean:17),
{\small\texttt{iterated\allowbreak{}Metric\allowbreak{}Covariant\allowbreak{}Derivative}}
(Metric.lean:30), and
{\small\texttt{curvature\allowbreak{}Derivative\allowbreak{}Norm}}
(DerivativeNorm.lean:18).

\subsection{C2. Curvature radius, volume collapse, and derivative
tests}\label{collapse-radius}

\textbf{Definitions.} The curvature radius is the extended-real supremum

\[
 R_g(p)=\sup\left\{r>0:\sec_g(q)\geq-r^{-2}\ \text{for every }q\in B_g(p,r)\right\}.
\]

This radius is \textbf{not capped at 1}. It is permitted to equal
\(+\infty\). It uses a sectional lower bound on the whole ball, not a
curvature condition only at its center. The code does not insert a
separately supplied scale function into this definition.

``Volume collapsed at curvature scale with parameter \(w\)'' means
exactly

\[
 \forall r>0,\quad R_g(p)=r\ \Longrightarrow\ V_g(p,r)\leq wr^3.
\]

The equality is equality in extended nonnegative reals. Therefore this
condition tests the one finite positive curvature radius when it exists.
\textbf{It is vacuous when \(R_g(p)=+\infty\)}, and also vacuous if a
radius were zero. The latter possibility is excluded by the first
admitted lemma below for the stated smooth-manifold setting. No
expression \(\operatorname{Vol}B(p,\infty)/\infty^3\) is used.

For a nonnegative integer \(K\), a real-valued function
\(A:\mathbb R\to\mathbb R\), and a real threshold \(w_0\), ``curvature
derivatives controlled'' means:

\[
\begin{split}
&\text{for every }p,\ w,r\in\mathbb R\text{ with }w_0\leq w<\omega_3,\quad 0<r<R_g(p),\\
&V_g(p,r)\geq wr^3
\quad\Longrightarrow\quad
 D_k(g,q)\leq A(w)r^{-k-2}
 \quad\text{for every }0\leq k\leq K\text{ and every }q\in B_g(p,r).
\end{split}
\]

The derivative bound includes \(k=0\) and \(k=K\). The radius inequality
is \textbf{strict}. The volume inequality is \textbf{nonstrict}. The
estimate is on \textbf{every point of the entire same ball}. The
definition by itself does not require \(A\) to be positive or \(w_0\) to
be positive; the main existence theorems supply these hypotheses. No
continuity or monotonicity of \(A\) is built into the predicate. The
metric and the function \(A\) are the same in all tests.

\textbf{Admitted lemma 1.} For every smooth metric in the auxiliary
smooth-manifold setting and every point \(p\),

\[
 R_g(p)>0.
\]

There is no compactness, completeness, connectedness, or boundaryless
assumption in this lemma.

\textbf{Admitted lemma 2.} If that manifold is compact and connected,
then, for every point \(p\),

\[
 R_g(p)=+\infty\quad\Longleftrightarrow\quad \sec_g\geq0\text{ everywhere}.
\]

This includes manifolds with boundary and other allowed models with
corners; the lemma is not restricted to closed 3-manifolds. Its forward
direction uses connectedness to reach every point from \(p\).
Compactness is an explicit hypothesis even if a stronger result might be
possible.

\textbf{Independently proved lemma.} If \(w_1\leq w_2\), derivative
control at threshold \(w_1\) implies derivative control at threshold
\(w_2\). The latter condition asks about fewer values of \(w\). This
implication requires no positivity assumption on the thresholds or on
\(A\).

\textbf{Lean references:} all eight declarations in CurvatureScale.lean:
{\small\texttt{curvature\allowbreak{}Radius}} (23),
{\small\texttt{ball\allowbreak{}Volume}} (28),
{\small\texttt{euclidean\allowbreak{}Unit\allowbreak{}Ball\allowbreak{}Volume}}
(31),
{\small\texttt{volume\allowbreak{}Collapsed\allowbreak{}At\allowbreak{}Curvature\allowbreak{}Scale}}
(33),
{\small\texttt{curvature\allowbreak{}Derivatives\allowbreak{}Controlled}}
(38), {\small\texttt{curvature\allowbreak{}Radius\_\allowbreak{}pos}}
(46),
{\small\texttt{curvature\allowbreak{}Radius\_\allowbreak{}eq\_\allowbreak{}top\_\allowbreak{}iff}}
(50), and
{\small\texttt{curvature\allowbreak{}Derivatives\allowbreak{}Controlled\_\allowbreak{}mono\_\allowbreak{}threshold}}
(55).

\subsection{C3. The exact model cusp and its curvature
normalization}\label{collapse-cuspmodel}

\textbf{Definitions.} The fixed topological torus is
\(T^2=S^1\times S^1\), with its product smooth structure. The
half-cylinder is

\[
 \mathcal H=T^2\times[0,\infty),
\]

with the product manifold-with-boundary model. The collar domain is the
relatively open subset

\[
 \mathcal U=T^2\times[0,100).
\]

In particular, \(z=0\) is included and \(z=100\) is excluded.

An exact cusp consists of a \textbf{supplied} smooth metric \(q\) on
this torus, the condition

\[
 \mathrm{Rm}_q(v,w,w,v)=0\quad\text{for every point and all tangent vectors},
\]

and a supplied smooth metric \(h\) on the whole half-cylinder satisfying
exactly

\[
 h=dz^2+e^{-z}q.
\]

These are data and equations, not an existence theorem producing \(q\)
or \(h\). No diameter, area, injectivity-radius, lattice-shape, or
aspect-ratio restriction is imposed on \(q\) at this stage. There is no
quotient group or chosen lattice basis in this definition.

\textbf{Admitted lemma 3.} For every such cusp and every point and pair
of tangent vectors,

\[
 \mathrm{Rm}_h(v,w,w,v)
 =-\frac14\bigl(h(v,v)h(w,w)-h(v,w)^2\bigr).
\]

Thus the normalization is sectional curvature \(-1/4\), not \(-1\). The
statement is a tensor identity including degenerate vector pairs. It is
asserted on the whole half-cylinder, including its boundary, not merely
on the depth-100 collar. It does not assert completeness of the interior
\(T^2\times(0,\infty)\), finite volume for any arbitrary metric on an
ambient manifold, or existence of a cusp embedding.

\textbf{Lean references:} Hyperbolic/Cusp.lean:
{\small\texttt{Cusp\allowbreak{}Half\allowbreak{}Space}} (10),
{\small\texttt{cusp\allowbreak{}Domain}} (12),
{\small\texttt{Hyperbolic\allowbreak{}Cusp}} (14),
{\small\texttt{cusp\_\allowbreak{}constant\_\allowbreak{}sectional\_\allowbreak{}curvature}}
(23).

\subsection{C4. Exactly what a nearly cuspidal boundary
supplies}\label{collapse-cuspembedding}

Let \(W\) be a compact carrier with smooth metric \(g\). Fix
\(K\in\mathbb N\) and \(\delta\in\mathbb R\).

\textbf{Definitions of the error.} For a chosen exact cusp
\((\mathcal H,h)\) and a map \(f:\mathcal H\to W\), set

\[
 E=f^*g-h,\qquad
 (f^*g)_p(v,w)=g_{f(p)}(df_pv,df_pw).
\]

The code forms this as an actual covariant two-tensor using the manifold
derivative of \(f\). The error bound means

\[
 |\nabla_h^k E(p)|_h\leq\delta
 \quad(0\leq k\leq K,\ p\in T^2\times[0,100)).
\]

Both the derivatives and their tensor norms use the \textbf{reference
cusp metric} \(h\). This is a pointwise uniform bound on each derivative
order, not a sum over orders. It is not a coordinate coefficient norm or
a bound on derivatives measured with the target metric \(g\).

\textbf{Definition of a cusp embedding onto a prescribed boundary set
\(X\subset W\).} The supplied data must satisfy all of the following:

\begin{enumerate}
\def\labelenumi{\arabic{enumi}.}
\tightlist
\item
  An exact cusp as in C3 and a map \(f:\mathcal H\to W\).
\item
  The map is \(C^{K+1}\) on \(\mathcal U\), including its boundary
  points.
\item
  Its restriction to \(\mathcal U\), with the subspace topology, is a
  topological embedding.
\item
  Its differential is injective at every point of \(\mathcal U\).
\item
  Its depth-zero torus has exactly the prescribed image:
  \(f(T^2\times\{0\})=X\).
\item
  For every \((t,z)\in\mathcal U\), \(f(t,z)\in\partial W\) if and only
  if \(z=0\). Thus positive-depth points do not meet any ambient
  boundary component.
\item
  The error bounds displayed above hold through order \(K\).
\end{enumerate}

There are no conditions on \(f\) outside \(\mathcal U\). The map is
defined there because Lean represents maps as total functions. There is
no explicit field asserting a smooth inverse, a proper map on the
noncompact half-open collar, or a prescribed orientation for the
embedding.

\textbf{Definition of a nearly cuspidal boundary.} This consists of a
positive integer \(n\), subsets \(X_i\subset W\),
\(i\in\{0,\ldots,n-1\}\), and cusp embeddings onto those sets, with:

\[
 X_i\text{ connected and closed},\qquad X_i\cap X_j=\varnothing\ (i\ne j),\qquad
 \bigcup_iX_i=\partial W,
\]

and

\[
 d_g(x,y)\leq\delta\quad\text{for all }x,y\in X_i.
\]

Here connected sets are nonempty. The finite disjoint closed cover by
connected sets gives the actual boundary components, not arbitrary
pieces of one component. The diameter is measured by curves in
\textbf{all of \(W\)}. It is not the intrinsic length distance
constrained to \(X_i\), and it is not measured with the model torus
metric \(q_i\). This precise convention matters for comparison with
sources.

\textbf{No disjointness of the positive-depth cusp collars is required.}
The sets \(X_i\) are pairwise disjoint, and each collar individually is
embedded, but there is no field saying the images \(f_i(\mathcal U)\)
and \(f_j(\mathcal U)\) are disjoint. Nor does the definition give
common torus coordinates or compatible fiber structures on different
collars.

\textbf{Constructed operations, with no admissions.} If
\(\delta_1\leq\delta_2\), either a cusp embedding or a nearly cuspidal
boundary with error parameter \(\delta_1\) can be viewed as one with
error parameter \(\delta_2\), keeping the exact same maps, component
sets, and indices. Only the inequalities are weakened. The operations
impose no positivity assumption; substantive applications choose
positive thresholds.

\textbf{Lean references:} CuspBoundary.lean:
{\small\texttt{cusp\allowbreak{}Metric\allowbreak{}Error}} (15),
{\small\texttt{cusp\allowbreak{}Metric\allowbreak{}Error\allowbreak{}Bound}}
(22), {\small\texttt{Cusp\allowbreak{}Embedding}} (29),
{\small\texttt{Nearly\allowbreak{}Cuspidal\allowbreak{}Boundary}} (43),
{\small\texttt{Cusp\allowbreak{}Embedding.\allowbreak{}weaken}} (64),
{\small\texttt{Nearly\allowbreak{}Cuspidal\allowbreak{}Boundary.\allowbreak{}weaken}}
(70).

\subsection{C5. Which points must satisfy volume collapse near
boundary}\label{collapse-boundarydistance}

\textbf{Definitions.} Distance to the actual boundary is

\[
 d_g(p,\partial W)=\inf_{q\in\partial W}d_g(p,q)\in[0,\infty].
\]

For empty boundary the infimum is \(+\infty\). The boundary
volume-collapse predicate requires exactly

\[
 d_g(p,\partial W)>10\quad\Longrightarrow\quad
 p\text{ is volume collapsed at curvature scale with parameter }w.
\]

The strict inequality is \(>10\), not \(\geq10\). No curvature-scale
volume-collapse condition is imposed by this predicate at a point whose
distance to boundary is at most 10. Whole-ball curvature-derivative
tests will still be imposed at every center in the boundary theorem
below.

Although the distance-to-boundary function has a mathematically defined
value for empty boundary, the principal boundary theorem requires a
positive number of cusp boundary components. The closed theorem uses its
own separate predicate.

\textbf{Lean references:} CuspBoundary.lean:
{\small\texttt{distance\allowbreak{}To\allowbreak{}Boundary}} (56),
{\small\texttt{boundary\allowbreak{}Volume\allowbreak{}Collapsed}} (60).

\subsection{C6. The three static hypothesis packages and their
monotonicity}\label{collapse-hypotheses}

\textbf{Definitions.} For fixed \((W,g,K,A,w_0)\), the closed package is
the conjunction of these three conditions:

\begin{enumerate}
\def\labelenumi{\arabic{enumi}.}
\tightlist
\item
  \(\partial W=\varnothing\).
\item
  Volume collapse at curvature scale with \(w_0\) holds at every point.
\item
  Derivative control as in C2 holds with \(K,A,w_0\).
\end{enumerate}

The boundary package is the conjunction of these three conditions:

\begin{enumerate}
\def\labelenumi{\arabic{enumi}.}
\tightlist
\item
  There exists a nearly cuspidal boundary with \(K\) and \(\delta=w_0\).
\item
  The volume-collapse condition of C5 holds with \(w=w_0\).
\item
  Derivative control as in C2 holds with \(K,A,w_0\).
\end{enumerate}

The combined static package is the \textbf{disjunction} of these two
packages. None of the three definitions itself contains a graph-manifold
predicate or the desired graph presentation as a hypothesis.
Connectedness, the lower bound on \(K\), and positivity of \(A\) are
supplied separately by the theorems, not hidden inside the definitions.

\textbf{Three independently proved monotonicity lemmas.} If
\(w_1\leq w_2\), then:

\begin{itemize}
\tightlist
\item
  volume collapse at curvature scale with \(w_1\) implies it with
  \(w_2\);
\item
  the closed package with \(w_1\) implies the closed package with
  \(w_2\);
\item
  the boundary package with \(w_1\) implies the boundary package with
  \(w_2\).
\end{itemize}

The second and third claims combine the easier volume and collar-error
inequalities with the fact that the derivative-trigger interval
\([w_2,\omega_3)\) is a subset of \([w_1,\omega_3)\). These proofs do
not appeal to any of the seven admissions. They are what later permits
taking the minimum of two independently supplied thresholds.

\textbf{Lean references:} GraphManifold.lean:
{\small\texttt{closed\allowbreak{}Collapse\allowbreak{}Hypotheses}}
(11),
{\small\texttt{boundary\allowbreak{}Collapse\allowbreak{}Hypotheses}}
(17),
{\small\texttt{static\allowbreak{}Collapse\allowbreak{}Hypotheses}}
(22),
{\small\texttt{volume\allowbreak{}Collapsed\allowbreak{}At\allowbreak{}Curvature\allowbreak{}Scale\_\allowbreak{}mono}}
(26),
{\small\texttt{closed\allowbreak{}Collapse\allowbreak{}Hypotheses\_\allowbreak{}mono}}
(35),
{\small\texttt{boundary\allowbreak{}Collapse\allowbreak{}Hypotheses\_\allowbreak{}mono}}
(43).

\subsection{C7. What ``a raw graph presentation'' means in the
conclusions}\label{collapse-output}

This output is a concrete presentation on the \textbf{same supplied
carrier \(W\)}. It is not an abstract assertion that some homeomorphic
manifold is a graph manifold. Its complete definition is translated in
the topology section of this report; the parts needed for interpreting
the collapse claims are these.

The output supplies a compact oriented cut carrier with finitely many
nonempty connected components. Each component is a smooth, locally
trivial \textbf{ordinary circle bundle} over a compact connected smooth
surface, whose base may have boundary and need not be orientable.
Exceptional Seifert fibers are not encoded as exceptional fibers of
these particular bundles: a further decomposition, for example
separating appropriate solid-torus blocks, must account for them when
such a presentation is constructed.

There are finitely many paired boundary tori, specified smooth matching
maps, actual collars, a quotient reconstruction homeomorphism onto
\(W\), smoothness and orientation conditions on that reconstruction, a
diffeomorphism on cut interiors, and smooth signed collars across the
reassembled seams. Original external boundary tori have explicit collars
and labels. Their cut-side and reassembled collars agree under
reconstruction. Component ownership of the paired and external tori is
recorded.

The paired tori are \textbf{not required incompressible}. Neither
primeness, irreducibility, a JSJ uniqueness statement, geometric metrics
on the blocks, nor a geometrization certificate is part of this output.
Those belong to later refinements. The collapse theorems return
existence of this data; they do not construct it in their current
{\small\texttt{sorry}} proofs.

\textbf{Lean dependency:}
{\small\texttt{GC.\allowbreak{}Graph\allowbreak{}Manifold.\allowbreak{}Raw\allowbreak{}Graph\allowbreak{}Presentation}},
Presentation.lean:129, together with its
{\small\texttt{Circle\allowbreak{}Fibration}},
{\small\texttt{Torus\allowbreak{}Pairing}}, and
{\small\texttt{Boundary\allowbreak{}Tori}} fields. Its topology
translation remains the authority for the full field-by-field expansion.

\subsection{C8. The three admitted graph-presentation
producers}\label{collapse-admitted}

\textbf{Admitted lemma 4: the separate nonnegative-curvature case.} Let
\(W\) be a connected compact oriented smooth 3-manifold, with empty
boundary, and let \(g\) be any smooth metric with \(\sec_g\geq0\)
everywhere. Then there exists a raw graph presentation of \(W\), in the
exact sense of C7.

There is \textbf{no small-volume assumption}, no prescribed \(K\), no
derivative function \(A\), and no curvature-radius finiteness assumption
here. This is a topological recognition/classification input for closed
nonnegative-curvature 3-manifolds. It is separate because
curvature-scale volume collapse is vacuous for infinite curvature
radius, and the finite-radius argument in C10 cannot absorb such
components.

\textbf{Admitted theorem 5: closed static collapse.} In precisely this
order:

\[
 \forall K\in\mathbb N\ (K\geq10),\quad
 \forall A:\mathbb R\to\mathbb R\ \relax
 [\forall w\in(0,\omega_3),\ A(w)>0],\quad
 \exists w_0\in(0,\omega_3),
\]

such that \textbf{for every} connected compact oriented smooth 3-carrier
\(W\) and \textbf{every} smooth metric \(g\) on it, the closed package
of C6 implies existence of a raw graph presentation of \(W\).

The threshold depends on \(K,A\), not on \(W\), \(g\), or a finite
family chosen later. The function \(A\) is real-valued and hence finite;
no monotonicity, continuity, measurability, or local boundedness of it
is assumed. Its values outside \((0,\omega_3)\) are unrestricted and
irrelevant to these tests. The theorem does not give an explicit
numerical threshold. The output need not respect any previously chosen
circle fibration because no such fibration is an input.

\textbf{Admitted theorem 6: static collapse with nearly cuspidal
boundary.} The same quantifier order supplies a possibly different
\(w_0\in(0,\omega_3)\). For every connected compact oriented smooth
carrier \(W\), smooth metric \(g\), and \textbf{supplied specific nearly
cuspidal boundary} \(B\) with \(K\) and \(\delta=w_0\), assume the C5
volume-collapse condition and the C2 derivative control. Then there
exist:

\[
 G:\text{a raw graph presentation of }W,
 \qquad e:\{0,\ldots,B.n-1\}\ \xrightarrow{\ \cong\ }\ \relax
 \{0,\ldots,G.n_{\rm ext}-1\},
\]

with

\[
 \operatorname{image}(G\text{'s external torus map numbered }e(i))=B.X_i
 \quad\text{for every }i.
\]

Thus the number and identity of the original boundary components are
retained through an explicit index bijection and equality of actual
image sets. The theorem does \textbf{not} say that the output torus
parametrization equals the supplied cusp parametrization, that the
output external collar equals the finite-regularity cusp collar, that
prescribed longitude/meridian classes are retained, or that their circle
fibrations agree. Such stronger relative assertions cannot be read into
this statement merely from the word ``labels.'' The raw presentation
itself has its own smooth marked collars, whose compatibility is part of
C7.

The boundary theorem applies at nonempty boundary because \(B.n>0\). The
derivative condition is imposed at every center, including the boundary
and its depth-10 neighborhood; the volume-collapse condition alone has
the C5 exception there.

\textbf{Lean references:} GraphManifold.lean:
{\small\texttt{exists\_\allowbreak{}raw\allowbreak{}Graph\allowbreak{}Presentation\_\allowbreak{}of\_\allowbreak{}nonnegative}}
(52),
{\small\texttt{exists\_\allowbreak{}closed\_\allowbreak{}graph\_\allowbreak{}threshold}}
(60),
{\small\texttt{exists\_\allowbreak{}boundary\_\allowbreak{}graph\_\allowbreak{}threshold}}
(68).

\subsection{C9. The common threshold and finite-family
conclusion}\label{collapse-composition}

\textbf{Composition theorem with a genuine Lean proof.} With exactly the
\(K,A\) hypotheses of C8, there is one \(w_0\in(0,\omega_3)\) such that
the combined static package of C6 yields a raw graph presentation for
every connected compact carrier and metric.

The proof chooses the minimum of the admitted closed and boundary
thresholds, uses the proved monotonicity statements, and applies the
appropriate admitted theorem. Therefore it has no {\small\texttt{sorry}}
in its own body but \textbf{does depend on
{\small\texttt{sorry\allowbreak{}Ax}}}. In the boundary case this
combined theorem deliberately discards the additional boundary-index
bijection and image identities returned by the stronger producer; its
conclusion only says that a raw graph presentation exists.

\textbf{Finite-family composition theorem with a genuine Lean proof.}
The same choice of \(w_0\) works simultaneously for every natural number
\(n\), every family of connected compact carriers \(W_i\), \(i<n\), and
every family of smooth metrics \(g_i\), provided each satisfies the
combined static package with the same \(K,A,w_0\). The result is a
family of actual raw graph presentations, one on each \(W_i\).

The threshold is chosen \textbf{before} \(n\) and the family. The family
may be empty, \(n=0\), in which case the conclusion is the empty
function. Each individual connected carrier is still nonempty. The
conclusion retains the family index. It does not turn the disjoint union
into a single connected manifold, identify any flow carriers, or produce
one set of surgery/flow parameters.

\textbf{Lean references:} GraphManifold.lean:
{\small\texttt{exists\_\allowbreak{}graph\_\allowbreak{}threshold}}
(80),
{\small\texttt{exists\_\allowbreak{}componentwise\_\allowbreak{}graph\_\allowbreak{}threshold}}
(97).

\subsection{C10. Absorbing a sequence of eventual bounds into one
derivative function}\label{collapse-uniform}

\textbf{Admitted lemma 7.} Fix any \(K\in\mathbb N\); this lemma does
not require \(K\geq10\). Let \((W_j,g_j)\), \(j\in\mathbb N\), be a
sequence of compact carriers with smooth metrics. They are not assumed
connected, nonempty, or boundaryless. Supply real numbers \(R_j>0\) and
\(B_j\geq0\), with the following two bounds for each \(j\):

\[
 0<r<R_{g_j}(p)\quad\Longrightarrow\quad r\leq R_j
 \quad\text{for every }p\in W_j,
\]

and

\[
 D_k(g_j,p)\leq B_j\quad
 \text{for every }p\in W_j\text{ and every }0\leq k\leq K.
\]

The same \(B_j\) bounds all these derivative orders, although it may
vary arbitrarily with \(j\). The finite radius bound is a substantive
premise: if a nonempty component has infinite curvature radius, it
fails, since all positive real radii would have to be at most \(R_j\).
An empty carrier creates no difficulty because the pointwise premises
are vacuous.

Assume in addition that for each \(w\in(0,\omega_3)\) there are an
integer \(N(w)\) and a finite positive real \(C(w)\) such that, for all
\(j\geq N(w)\), all centers \(p\), and all \(0<r<R_{g_j}(p)\),

\[
 V_{g_j}(p,r)\geq wr^3
 \quad\Longrightarrow\quad
 D_k(g_j,q)\leq C(w)r^{-k-2}
 \quad(0\leq k\leq K,\ q\in B_{g_j}(p,r)).
\]

The conclusion is the existence of \textbf{one} function
\(A:\mathbb R\to\mathbb R\) satisfying \(A(w)>0\) for every \(w>0\),
such that for \textbf{every} sequence index \(j\) and \textbf{every}
\(w_0>0\), the derivative-control predicate of C2 holds for
\((g_j,K,A,w_0)\).

In expanded form, all the above volume-triggered tests hold for every
\(j\) and every \(0<w<\omega_3\), with \(C(w)\) replaced by this one
\(A(w)\). No common tail \(N\) for all \(w\) is assumed or concluded.
The output \(A\) may depend on the entire sequence and all its supplied
bounds. It is chosen before the static collapse theorem selects \(w_0\).
The conclusion for \(w_0\geq\omega_3\) is harmlessly vacuous, because
its derivative-trigger interval is empty.

The intended elementary proof, still admitted in Lean, takes a maximum
of the tail bound \(C(w)\) and finitely many quantities such as

\[
 B_jR_j^{k+2},\qquad j<N(w),\quad 0\leq k\leq K,
\]

then adds a positive constant. This explains why no uniform-in-\(w\)
tail is needed. It is not a proof that these bounds arise from a Ricci
flow: producing the sequence, intrinsic radius bounds, and eventual
whole-ball estimates belongs elsewhere.

\textbf{Lean reference:} UniformDerivativeBounds.lean:
{\small\texttt{exists\_\allowbreak{}common\_\allowbreak{}curvature\_\allowbreak{}derivative\_\allowbreak{}bound}}
(10).

\subsection{C11. Mathematical review findings and approval
points}\label{collapse-review}

No definite false assertion or explicit logical contradiction was found
in this read-through of these seven modules and their directly relevant
definitions. This is a qualified review conclusion: the seven admitted
claims remain mathematical obligations, and the stronger output notion
of a raw graph presentation must be checked against the topology
translation. A successful Lean build is not evidence that the admitted
mathematics is true.

The main distinctions that must survive subsequent implementation are:

\begin{enumerate}
\def\labelenumi{\arabic{enumi}.}
\tightlist
\item
  \textbf{Intrinsic balls and all-point estimates.} The static
  assumptions use the carrier's own metric distance and volume and
  estimate derivatives at all points of the same ball. Center estimates
  or estimates on ambient flow balls are not interchangeable with them.
\item
  \textbf{Infinite curvature radius is handled separately.} The
  collapse-at-radius predicate is vacuous at infinity. The closed
  theorem asserts a graph conclusion even in that situation; it must use
  nonnegative-curvature recognition. The sequence lemma cannot be
  applied to a nonempty infinite-radius component.
\item
  \textbf{Finite cusp regularity is explicit.} Maps are \(C^{K+1}\),
  error derivatives are measured through \(K\), and the reference metric
  is smooth. No assertion that an arbitrary finite-order metric can be
  smoothed while preserving curvature is present here.
\item
  \textbf{Boundary control has exact scope.} There is no requirement
  that different cusp collars be disjoint. Boundary diameter uses the
  distance in \(W\), not its induced boundary metric. The distance-10
  volume exception does not remove any derivative tests. These should be
  accepted deliberately, not supplied by readers from habit.
\item
  \textbf{Retained boundary labels are weaker than retained markings.}
  The stronger boundary theorem returns a bijection of indices with
  equality of image sets; it does not retain a prescribed
  parametrization, slope basis, fibration, or finite-regularity collar.
  The generic common-threshold theorem forgets even that extra
  bijection.
\item
  \textbf{The quantifier order is load-bearing.} Fix \(K,A\); then
  choose one \(w_0\); then allow all carriers and metrics. In the
  sequence lemma, fix the sequence first, then construct one \(A\), then
  invoke static collapse for \(w_0\). No uniform bound across all
  possible flows or sequences is asserted.
\item
  \textbf{This is an interface for the principal static theorem.} It
  does not separately state the local model, marker, cloud, one-sheet,
  or simultaneous cutoff constructions from the blueprint. Their proof
  work is currently inside the two large admitted threshold producers.
  It would be misleading to assign every Chapter 14 construction a
  separate Lean statement already present here.
\end{enumerate}

For this report, actual source definitions were reread in the seven
modules and in the dependent carrier, graph-presentation, Riemannian
distance/ball/volume, pullback, curvature, tensor metric, and coordinate
covariant-derivative files. Blueprint BBR03
({\small\texttt{master207B.\allowbreak{}tex}}, lines 10592--10651) and
LC89 ({\small\texttt{master207A.\allowbreak{}tex}}, lines 31165--31204)
were reread to check the selected boundary and sequence contracts. Their
strict inequalities and order of quantifiers agree with the translation
above. Existing exact Kleiner--Lott source and errata checks recorded in
{\small\texttt{collapse\_\allowbreak{}review.\allowbreak{}md}} were not
rerun and are not represented here as a new literature audit. No source
archive, Lean file, theorem contract, or blueprint text was changed.

\subsection{C12. File key and declaration
coverage}\label{collapse-files}

All paths below are relative to the checked code repository
{\small\texttt{GC\_\allowbreak{}BASELINE\_\allowbreak{}EXPORT}} and lie
under {\small\texttt{Differential\allowbreak{}Geometry/\allowbreak{}}}.

\begin{longtable}[]{@{}
  >{\raggedright\arraybackslash}p{(\linewidth - 2\tabcolsep) * \real{0.5000}}
  >{\raggedright\arraybackslash}p{(\linewidth - 2\tabcolsep) * \real{0.5000}}@{}}
\toprule\noalign{}
\begin{minipage}[b]{\linewidth}\raggedright
Short file reference
\end{minipage} & \begin{minipage}[b]{\linewidth}\raggedright
Actual file
\end{minipage} \\
\midrule\noalign{}
\endhead
\bottomrule\noalign{}
\endlastfoot
CurvatureScale.lean &
{\small\texttt{Geometry/\allowbreak{}Collapse/\allowbreak{}Curvature\allowbreak{}Scale.\allowbreak{}lean}} \\
CuspBoundary.lean &
{\small\texttt{Geometry/\allowbreak{}Collapse/\allowbreak{}Cusp\allowbreak{}Boundary.\allowbreak{}lean}} \\
GraphManifold.lean &
{\small\texttt{Geometry/\allowbreak{}Collapse/\allowbreak{}Graph\allowbreak{}Manifold.\allowbreak{}lean}} \\
UniformDerivativeBounds.lean &
{\small\texttt{Geometry/\allowbreak{}Collapse/\allowbreak{}Uniform\allowbreak{}Derivative\allowbreak{}Bounds.\allowbreak{}lean}} \\
Hyperbolic/Cusp.lean &
{\small\texttt{Geometry/\allowbreak{}Hyperbolic/\allowbreak{}Cusp.\allowbreak{}lean}} \\
Metric.lean &
{\small\texttt{Geometry/\allowbreak{}Connection/\allowbreak{}Tensor\allowbreak{}Nabla/\allowbreak{}Iterated/\allowbreak{}Metric.\allowbreak{}lean}} \\
DerivativeNorm.lean &
{\small\texttt{Geometry/\allowbreak{}Curvature/\allowbreak{}Metric/\allowbreak{}Derivative\allowbreak{}Norm.\allowbreak{}lean}} \\
Presentation.lean (dependent definition) &
{\small\texttt{Topology/\allowbreak{}Three\allowbreak{}Manifold/\allowbreak{}Graph\allowbreak{}Manifold/\allowbreak{}Presentation.\allowbreak{}lean}} \\
\end{longtable}

The accompanying
{\small\texttt{collapse\_\allowbreak{}coverage.\allowbreak{}json}} maps
every authored declaration to its exact section, original source line,
admitted/independent/composition/definition status, and inspected audit
axiom status. It covers all 35 authored declarations in the seven
modules; automatically generated structure projections and recursors are
represented through the complete field-by-field descriptions of their
parent structures.

\section{Hyperbolic geometry, exterior disk area, and the
incompressibility
consumers}\label{hyperbolic-geometry-exterior-disk-area-and-the-incompressibility-consumers}

This report translates the \textbf{actual Lean statements} at commit
{\small\texttt{ea0fae6\allowbreak{}0ee01ef8\allowbreak{}c8d9b57a\allowbreak{}51794c65\allowbreak{}38f679c5\allowbreak{}d}}.
It covers all \textbf{28 authored declarations in six files}: nine
definitions, four directly admitted theorems, and fifteen theorems with
written Lean proofs. A written proof can still depend on an admitted
theorem; that distinction is recorded below. The code has not been
changed in preparing this report.

The six reviewed source files are:

\begin{itemize}
\tightlist
\item
  \href{https://github.com/qinz1yang/differential-geometry-dev/blob/ea0fae60ee01ef8c8d9b57a51794c6538f679c5d/DifferentialGeometry/Analysis/ODE/AreaUpperBarrier.lean}{Analysis/ODE/AreaUpperBarrier.lean}
\item
  \href{https://github.com/qinz1yang/differential-geometry-dev/blob/ea0fae60ee01ef8c8d9b57a51794c6538f679c5d/DifferentialGeometry/Geometry/Hyperbolic/Rigidity.lean}{Geometry/Hyperbolic/Rigidity.lean}
\item
  \href{https://github.com/qinz1yang/differential-geometry-dev/blob/ea0fae60ee01ef8c8d9b57a51794c6538f679c5d/DifferentialGeometry/Geometry/Hyperbolic/ModelAtlas.lean}{Geometry/Hyperbolic/ModelAtlas.lean}
\item
  \href{https://github.com/qinz1yang/differential-geometry-dev/blob/ea0fae60ee01ef8c8d9b57a51794c6538f679c5d/DifferentialGeometry/Geometry/MinimalSurface/ExteriorDiskArea.lean}{Geometry/MinimalSurface/ExteriorDiskArea.lean}
\item
  \href{https://github.com/qinz1yang/differential-geometry-dev/blob/ea0fae60ee01ef8c8d9b57a51794c6538f679c5d/DifferentialGeometry/Geometry/Flow/RicciFlow/LongTime/CuspIncompressibility.lean}{Geometry/Flow/RicciFlow/LongTime/CuspIncompressibility.lean}
\item
  \href{https://github.com/qinz1yang/differential-geometry-dev/blob/ea0fae60ee01ef8c8d9b57a51794c6538f679c5d/DifferentialGeometry/Geometry/Flow/RicciFlow/LongTime/ExteriorDiskFlow.lean}{Geometry/Flow/RicciFlow/LongTime/ExteriorDiskFlow.lean}
\end{itemize}

The central distinction for mathematical approval is this: \textbf{the
area argument is encoded as a correct conditional deduction, but the
existence of the required geometric area data is not established in
these files.} The continuous nonnegative area function and all-time
upper barriers would be contradictory. Producing them from a
hypothetical compressing torus is the difficult geometric work; that
work remains inside the broader admitted flow producer reviewed
elsewhere.

\subsection{HA-0. Conventions and status labels}\label{ha-0}

\begin{itemize}
\tightlist
\item
  \textbf{Definition:} an explicit mathematical object or predicate, not
  a theorem asserting that the predicate holds.
\item
  \textbf{Admitted:} the theorem body is {\small\texttt{sorry}}. Lean
  has checked its type, not its mathematical truth.
\item
  \textbf{Written proof, admission-dependent:} a real Lean proof whose
  audited dependency closure contains
  {\small\texttt{sorry\allowbreak{}Ax}}.
\item
  \textbf{Written proof, no skeleton admission:} the recorded
  declaration audit found no {\small\texttt{sorry\allowbreak{}Ax}} in
  its closure. This is about the code's proof dependencies, not an
  independent reproof of all inherited mathematics.
\end{itemize}

Unless a subsection says otherwise, a smooth 3-manifold here is modeled
on all of \(\mathbb R^3\), hence has no boundary as a manifold. Subsets
\(W\) may have boundary. The Mostow and model-atlas statements
explicitly assume Hausdorffness and sigma compactness. The generic
disk-area definitions and consumers do not add these assumptions. Their
use on the actual compact flow carriers supplies the customary manifold
conditions.

The loop parameter is \(\mathbb R/\mathbb Z\). The closed parameter disk
is \(\overline{\mathbb D}\subset\mathbb C\), with boundary
parametrization \(\theta\mapsto e^{2\pi i\theta}\). An equality of
boundary traces below is an equality of \textbf{parametrized maps}, not
merely of their images or homotopy classes.

\subsection{HA-1. Constant-curvature geometry and rigidity}\label{ha-1}

\subsubsection{HA-1.1. Constant sectional curvature}\label{ha-1-1}

\textbf{Lean:}
{\small\texttt{Differential\allowbreak{}Geometry.\allowbreak{}Geometry.\allowbreak{}Hyperbolic.\allowbreak{}has\allowbreak{}Constant\allowbreak{}Sectional\allowbreak{}Curvature}}.
\textbf{Status:} definition. Source:
{\small\texttt{Geometry/\allowbreak{}Hyperbolic/\allowbreak{}Rigidity.\allowbreak{}lean}},
line 20.

For a smooth Riemannian metric \(g\) on a smooth 3-manifold \(M\), the
predicate \(\operatorname{ConstSec}(g,K)\) means \[
 \forall p\in M\ \forall v,w\in T_pM,
 \quad (v,w)\text{ linearly independent}
 \ \Longrightarrow\ \sec_g(p;v,w)=K.
\] There is no assertion about the library's value of sectional
curvature on linearly dependent vectors. The definition itself does not
require connectedness, completeness, finite volume, Hausdorffness, or
sigma compactness. Its elaborated signature retains only the smooth
manifold and metric data it uses.

\subsubsection{HA-1.2. Marked finite-volume Mostow--Prasad
rigidity}\label{ha-1-2}

\textbf{Lean:}
{\small\texttt{Differential\allowbreak{}Geometry.\allowbreak{}Geometry.\allowbreak{}Hyperbolic.\allowbreak{}mostow\_\allowbreak{}prasad}}.
\textbf{Status:} admitted. Source:
{\small\texttt{Geometry/\allowbreak{}Hyperbolic/\allowbreak{}Rigidity.\allowbreak{}lean}},
line 24.

Let \(M,N\) be connected, nonempty, Hausdorff, sigma compact smooth
3-manifolds without boundary. Let \(g,h\) be smooth Riemannian metrics,
and let \(K<0\) be \textbf{the same constant for both metrics}. Assume:

\begin{enumerate}
\def\labelenumi{\arabic{enumi}.}
\tightlist
\item
  \(\operatorname{ConstSec}(g,K)\) and \(\operatorname{ConstSec}(h,K)\).
\item
  Both Riemannian metrics are complete.
\item
  Both total Riemannian volumes are finite.
\item
  A specified continuous homotopy equivalence \(u:M\to N\) is given.
\end{enumerate}

Then there is a unique smooth diffeomorphism \(f:M\to N\) satisfying
both \[
 h_{f(p)}(df_pv,df_pw)=g_p(v,w)
 \quad\text{for every }p,v,w,
 \qquad f\simeq u.
\] Thus uniqueness is \textbf{within the specified unbased homotopy
class}, among smooth diffeomorphisms preserving the metric. It is not
uniqueness among all isometries. It is not an equation between
arbitrarily chosen based fundamental-group homomorphisms. The input
homotopy equivalence includes a continuous inverse up to homotopy; no
smoothness of \(u\) is required.

``Complete'' is implemented as completeness of the intrinsic Riemannian
extended distance. In the connected case here this is the customary
metric completeness notion. Finite volume is literally
\(\operatorname{vol}_g(M)<+\infty\), and similarly for \(h\).

No orientability or compactness assumption is imposed. Empty manifolds
are excluded by {\small\texttt{Connected\allowbreak{}Space}}; noncompact
finite-volume cusped manifolds are included. This declaration contains
no quantitative closeness conclusion, control of cusp coordinates, or
time-dependent family of isometries.

\subsubsection{HA-1.3. Local curvature gives a hyperbolic model
atlas}\label{ha-1-3}

\textbf{Lean:}
{\small\texttt{Differential\allowbreak{}Geometry.\allowbreak{}Geometry.\allowbreak{}Hyperbolic.\allowbreak{}has\_\allowbreak{}hyperbolic\_\allowbreak{}atlas\_\allowbreak{}of\_\allowbreak{}curvature\_\allowbreak{}neg\_\allowbreak{}one}}.
\textbf{Status:} admitted. Source:
{\small\texttt{Geometry/\allowbreak{}Hyperbolic/\allowbreak{}Model\allowbreak{}Atlas.\allowbreak{}lean}},
line 17.

Let \(M\) be a Hausdorff, sigma compact smooth 3-manifold, with smooth
metric \(g\) and \(\operatorname{ConstSec}(g,-1)\). Then every point has
a smooth local isometry from an open set in the fixed hyperbolic
coordinate model. That model is \(\mathbb R^3\), with coordinates
\((x,y,z)\) and metric \[
 h_{\mathbb H}=e^{-2z}(dx^2+dy^2)+dz^2.
\] More literally, for each \(x\in M\), there is a partial smooth
diffeomorphism \(e\) from this coordinate model into \(M\), with \(x\)
in its target, for which \(e^*g=h_{\mathbb H}\) everywhere in its
source. This is precisely the inherited predicate
{\small\texttt{Has\allowbreak{}Thurston\allowbreak{}Atlas g .\allowbreak{}hyperbolic}}.

No completeness, finite volume, connectedness, or nonemptiness is
assumed. On the empty manifold the pointwise assertion is vacuous. The
output is local model charts; it does not supply a global developing
map, a deck group, an explicit quotient presentation, or a selected cusp
decomposition.

\subsubsection{HA-1.4. Complete finite-volume geometric structure from
curvature -1/4}\label{ha-1-4}

\textbf{Lean:}
{\small\texttt{Differential\allowbreak{}Geometry.\allowbreak{}Geometry.\allowbreak{}Hyperbolic.\allowbreak{}finite\allowbreak{}Volume\allowbreak{}Geometric\allowbreak{}Structure}}.
\textbf{Status:} explicit construction, admission-dependent through
HA-1.3. Source:
{\small\texttt{Geometry/\allowbreak{}Hyperbolic/\allowbreak{}Model\allowbreak{}Atlas.\allowbreak{}lean}},
line 23.

On a Hausdorff, sigma compact smooth 3-manifold \(M\), assume \(g\) is
smooth, complete, finite volume, and has constant sectional curvature
\(-1/4\). The construction returns a
{\small\texttt{Geometric\allowbreak{}Structure}} on the \textbf{same
whole carrier \(M\)}, having:

\begin{itemize}
\tightlist
\item
  model label {\small\texttt{hyperbolic}};
\item
  metric \(\widehat g=\tfrac14g\);
\item
  completeness of \(\widehat g\);
\item
  a hyperbolic model atlas for \(\widehat g\);
\item
  finite total volume for \(\widehat g\).
\end{itemize}

The scaling is deliberate: \(\sec_{\frac14g}=4\sec_g=-1\), and in
dimension three volumes scale by \((1/4)^{3/2}=1/8\). Completeness and
volume scaling use existing proved lemmas; the atlas uses HA-1.3. There
is no connectedness assumption. The returned metric is \textbf{not
literally the input metric}, and this construction does not assert that
a truncated cusp core with its restricted original metric is complete.

\textbf{Companion declaration:}
{\small\texttt{Differential\allowbreak{}Geometry.\allowbreak{}Geometry.\allowbreak{}Hyperbolic.\allowbreak{}finite\allowbreak{}Volume\allowbreak{}Geometric\allowbreak{}Structure\_\allowbreak{}model}},
line 42, says exactly that the model label of this construction is
{\small\texttt{hyperbolic}}. Its proof is {\small\texttt{rfl}}; the
recorded full declaration closure still contains the admitted atlas
theorem through the construction. It adds no new geometric content.

\subsection{HA-2. The real-variable area contradiction}\label{ha-2}

\subsubsection{HA-2.1. A strict local smooth upper
barrier}\label{ha-2-1}

\textbf{Lean:}
{\small\texttt{Differential\allowbreak{}Geometry.\allowbreak{}Analysis.\allowbreak{}has\allowbreak{}Local\allowbreak{}Smooth\allowbreak{}Upper\allowbreak{}Barrier}}.
\textbf{Status:} definition. Source:
{\small\texttt{Analysis/\allowbreak{}ODE/\allowbreak{}Area\allowbreak{}Upper\allowbreak{}Barrier.\allowbreak{}lean}},
line 13.

Given \(A:\mathbb R\to\mathbb R\), a set \(S\subseteq\mathbb R\), and
real numbers \(t,b\), this means that there are an open set
\(U\subseteq\mathbb R\) and a function \(F:\mathbb R\to\mathbb R\) such
that: \[
 t\in U,\qquad F\in C^\infty(U),\qquad F(t)=A(t),
 \qquad A(s)\le F(s)\ (s\in U\cap S),\qquad F'(t)<b.
\] \(F\) is only required smooth on \(U\), despite being a function
defined on all of \(\mathbb R\). The bound on its derivative is
\textbf{strict}, and is required only at the contact time. The area
function \(A\) is not assumed differentiable. The definition alone does
not even require \(t\in S\); the consumers below impose it.

When \(S=[T,\infty)\), the neighborhood \(U\) is open in \(\mathbb R\).
Domination at \(t=T\) is required only to the right, whereas at any
\(t>T\), including a surgery time in an application, it is required on
both sides sufficiently near the contact.

\subsubsection{HA-2.2. The general shifted contradiction}\label{ha-2-2}

\textbf{Lean:}
{\small\texttt{Differential\allowbreak{}Geometry.\allowbreak{}Analysis.\allowbreak{}not\_\allowbreak{}nonnegative\_\allowbreak{}area\_\allowbreak{}upper\_\allowbreak{}barriers\_\allowbreak{}shift}}.
\textbf{Status:} admitted. Source:
{\small\texttt{Analysis/\allowbreak{}ODE/\allowbreak{}Area\allowbreak{}Upper\allowbreak{}Barrier.\allowbreak{}lean}},
line 18.

Let \(A:\mathbb R\to\mathbb R\), and let \(T,c,D\in\mathbb R\) satisfy
\[
 T+c>0,\qquad 0<D<2\pi.
\] Assume \(A\) is continuous on \([T,\infty)\), is nonnegative there,
and at \textbf{every} \(t\ge T\) has a local smooth upper barrier with
\[
 F'(t)<\frac{3A(t)}{4(t+c)}-2\pi+D.
\] Then these assumptions imply a contradiction. Neither \(T\) nor \(c\)
separately needs to be positive. The condition \(T+c>0\) makes every
denominator on the half-line positive. There is no boundedness
assumption on \(A\), no uniform contact-neighborhood radius, and no
differentiability assumption on \(A\).

The intended elementary argument uses \((t+c)^{-3/4}A(t)\) and the
divergent integral of \((t+c)^{-3/4}\). This is an explanation of the
admitted statement, \textbf{not a claim that this proof has been
formalized}. A finite time interval, barriers only at regular flow
times, or continuity only between surgeries would not satisfy the
statement.

\subsubsection{HA-2.3. The three written specializations}\label{ha-2-3}

Each row is a separate authored declaration in the same file, proved
from HA-2.2 and therefore admission-dependent.

\begin{longtable}[]{@{}
  >{\raggedright\arraybackslash}p{(\linewidth - 4\tabcolsep) * \real{0.3000}}
  >{\raggedright\arraybackslash}p{(\linewidth - 4\tabcolsep) * \real{0.3000}}
  >{\raggedleft\arraybackslash}p{(\linewidth - 4\tabcolsep) * \real{0.4000}}@{}}
\toprule\noalign{}
\begin{minipage}[b]{\linewidth}\raggedright
Lean name after
{\small\texttt{Differential\allowbreak{}Geometry.\allowbreak{}Analysis.\allowbreak{}}}
\end{minipage} & \begin{minipage}[b]{\linewidth}\raggedright
Exact specialization
\end{minipage} & \begin{minipage}[b]{\linewidth}\raggedleft
Line
\end{minipage} \\
\midrule\noalign{}
\endhead
\bottomrule\noalign{}
\endlastfoot
{\small\texttt{not\_\allowbreak{}nonnegative\_\allowbreak{}area\_\allowbreak{}upper\_\allowbreak{}barriers}}
& Set \(c=1/4\); require \(T>-1/4\), \(0<D<2\pi\), continuity and
nonnegativity on the whole half-line, and the bound
\(3A/[4(t+1/4)]-2\pi+D\). & 28 \\
{\small\texttt{not\_\allowbreak{}nonnegative\_\allowbreak{}area\_\allowbreak{}upper\_\allowbreak{}barriers\_\allowbreak{}pi\_\allowbreak{}shift}}
& Set \(D=\pi\); require \(T+c>0\), continuity and nonnegativity on the
half-line, and the bound \(3A/[4(t+c)]-\pi\). & 39 \\
{\small\texttt{not\_\allowbreak{}nonnegative\_\allowbreak{}area\_\allowbreak{}upper\_\allowbreak{}barriers\_\allowbreak{}pi}}
& Set \(c=1/4,D=\pi\); require the stronger endpoint condition
\(T\ge0\), continuity and nonnegativity, and the bound
\(3A/[4(t+1/4)]-\pi\). & 52 \\
\end{longtable}

All three conclude {\small\texttt{False}}, meaning their complete
collections of hypotheses cannot occur. None is a theorem producing
barriers for a Ricci flow.

\subsection{HA-3. What the exterior disk area actually
means}\label{ha-3}

\subsubsection{HA-3.1. The competitor class}\label{ha-3-1}

\textbf{Lean:}
{\small\texttt{Differential\allowbreak{}Geometry.\allowbreak{}Geometry.\allowbreak{}Minimal\allowbreak{}Surface.\allowbreak{}is\allowbreak{}Exterior\allowbreak{}Spanning\allowbreak{}Disk}}.
\textbf{Status:} definition. Source:
{\small\texttt{Geometry/\allowbreak{}Minimal\allowbreak{}Surface/\allowbreak{}Exterior\allowbreak{}Disk\allowbreak{}Area.\allowbreak{}lean}},
line 15.

Fix a subset \(W\subseteq M\), a continuous loop
\(\gamma:\mathbb R/\mathbb Z\to M\), and a continuous map
\(u:\overline{\mathbb D}\to M\). It is an admissible exterior spanning
disk precisely when all the following hold:

\begin{enumerate}
\def\labelenumi{\arabic{enumi}.}
\tightlist
\item
  \(u(e^{2\pi i\theta})=\gamma(\theta)\) for every \(\theta\).
\item
  \(\gamma(\mathbb R/\mathbb Z)\subseteq\operatorname{frontier}(W)\).
\item
  \(u\) is a topological embedding.
\item
  \(u(\overline{\mathbb D})\subseteq W\).
\item
  \(u(\mathbb D)\subseteq\operatorname{interior}(W)\).
\item
  There exists a function \(U:\mathbb C\to M\), equal to \(u\) on
  \(\overline{\mathbb D}\), that is smooth on an open neighborhood of
  the closed disk, and whose differential is injective at every point of
  the closed disk.
\end{enumerate}

Thus the competitors are smoothly immersed up to and across their
parameter boundary, and topologically embedded. No branch point is
allowed, including at the boundary. The extension need not be smooth
outside the specified neighborhood or have image in \(W\) outside the
closed disk.

There is \textbf{no metric in this predicate}. In particular it does not
mean minimal, stable, conformal, least area, or orthogonal to the
frontier. It does not require \(W\) to be closed, compact, connected, a
submanifold, or mean convex, nor identify \(W\) as the complement of any
particular hyperbolic core. The word ``exterior'' is a name for this
precisely specified side condition. For an open \(W\), the boundary
requirement and \(u(\overline{\mathbb D})\subseteq W\) usually make the
class empty; actual applications intend a compact exterior including its
boundary.

The elaborated predicate requires just topology and charts on \(M\); it
does not retain the surrounding file's unused smooth-manifold instance.
Its use with the area below does require the smooth metric/manifold
data.

\subsubsection{HA-3.2. The infimum and the meaning of
``area''}\label{ha-3-2}

\textbf{Lean:}
{\small\texttt{Differential\allowbreak{}Geometry.\allowbreak{}Geometry.\allowbreak{}Minimal\allowbreak{}Surface.\allowbreak{}least\allowbreak{}Exterior\allowbreak{}Disk\allowbreak{}Area}}.
\textbf{Status:} definition. Source:
{\small\texttt{Geometry/\allowbreak{}Minimal\allowbreak{}Surface/\allowbreak{}Exterior\allowbreak{}Disk\allowbreak{}Area.\allowbreak{}lean}},
line 22.

For a smooth metric \(g\), define \[
 a(g,W,\gamma)=\inf\{\operatorname{Area}_g(u):
       u\text{ satisfies HA-3.1}\}\in\mathbb R.
\] This is a \textbf{real-valued infimum}, not an extended-real infimum.
The inherited
{\small\texttt{riemannian\allowbreak{}Disk\allowbreak{}Area}} is the
integral over the Euclidean closed unit disk of the ordinary
two-dimensional metric Jacobian: \[
 \int_{\overline{\mathbb D}}
 \sqrt{g(U_x,U_x)g(U_y,U_y)-g(U_x,U_y)^2}\,dx\,dy.
\] Formally the library first extends a continuous disk map to
\(\mathbb C\) using the radial metric projection onto the closed disk.
Its derivative at the circle need not be the derivative of a smooth
extension. This causes no area discrepancy: the circle has measure zero,
and the inherited theorem
{\small\texttt{riemannian\allowbreak{}Disk\allowbreak{}Area\_\allowbreak{}eq\_\allowbreak{}of\_\allowbreak{}extension}}
proves equality with the displayed integral for any agreeing extension.
For the admitted competitor class, the smooth extension on a
neighborhood of the compact disk gives an integrable Jacobian by the
inherited compact integrability theorem. Thus these competitors have the
ordinary finite Riemannian area; the integral's totalized value on
nonintegrable arbitrary maps is not being used to manufacture competitor
areas.

The real infimum of the empty set is \textbf{zero in Lean}. Therefore
\(a(g,W,\gamma)\ge0\) by itself does not assert that a disk exists.
Attainment must be supplied separately. No strict positivity theorem,
minimality theorem, stability theorem, or Meeks--Yau theorem is asserted
by this definition.

Two authored helpers are proved with no skeleton admission:

\begin{itemize}
\tightlist
\item
  {\small\texttt{least\allowbreak{}Exterior\allowbreak{}Disk\allowbreak{}Area\_\allowbreak{}nonneg}}
  (line 27): for every \(g,W,\gamma\), \(0\le a(g,W,\gamma)\), including
  the empty-class case.
\item
  {\small\texttt{least\allowbreak{}Exterior\allowbreak{}Disk\allowbreak{}Area\_\allowbreak{}le}}
  (line 33): if a specified \(u\) is admissible,
  \(a(g,W,\gamma)\le\operatorname{Area}_g(u)\). The proof explicitly
  bounds all competitor areas below by zero before applying the infimum
  rule.
\end{itemize}

\subsubsection{HA-3.3. Time-varying carriers and late-only
loops}\label{ha-3-3}

\textbf{Lean:}
{\small\texttt{Differential\allowbreak{}Geometry.\allowbreak{}Geometry.\allowbreak{}Minimal\allowbreak{}Surface.\allowbreak{}exterior\allowbreak{}Disk\allowbreak{}Area\allowbreak{}On}}.
\textbf{Status:} definition. Source:
{\small\texttt{Geometry/\allowbreak{}Minimal\allowbreak{}Surface/\allowbreak{}Exterior\allowbreak{}Disk\allowbreak{}Area.\allowbreak{}lean}},
line 43.

One may specify a different smooth 3-manifold \(M_t\) for each real
\(t\), a smooth metric \(g_t\) on each \(M_t\), a subset
\(W_t\subseteq M_t\), a real starting time \(T\), and a continuous loop
\(\gamma_t\) \textbf{only for \(t\ge T\)}. Set \[
 A(t)=\begin{cases}
       a(g_t,W_t,\gamma_t),&t\ge T,\\
       0,&t<T.
      \end{cases}
\] There is no topology on the total family \(\coprod_tM_t\), no
continuity of \(g_t,W_t,\gamma_t\), no Ricci equation, and no surgery
assumption in this definition. The manifolds are allowed to change their
underlying types. The loop parameter depending on a proof of \(T\le t\)
is Lean bookkeeping; proof irrelevance makes it a single loop for each
permitted time.

\textbf{Companion theorem:}
{\small\texttt{exterior\allowbreak{}Disk\allowbreak{}Area\allowbreak{}On\_\allowbreak{}nonneg}}
(line 52), proved with no skeleton admission, says \(A(t)\ge0\) for
every real \(t\), including \(t<T\). No continuity from the left at
\(T\) is asserted. No loop is required on an earlier empty carrier.

\subsubsection{HA-3.4. Continuity from two-direction comparisons of
minimizers}\label{ha-3-4}

\textbf{Lean:}
{\small\texttt{Differential\allowbreak{}Geometry.\allowbreak{}Geometry.\allowbreak{}Minimal\allowbreak{}Surface.\allowbreak{}continuous\allowbreak{}On\_\allowbreak{}least\allowbreak{}Exterior\allowbreak{}Disk\allowbreak{}Area\_\allowbreak{}of\_\allowbreak{}local\_\allowbreak{}disk\_\allowbreak{}comparisons}}.
\textbf{Status:} admitted. Source:
{\small\texttt{Geometry/\allowbreak{}Minimal\allowbreak{}Surface/\allowbreak{}Exterior\allowbreak{}Disk\allowbreak{}Area.\allowbreak{}lean}},
line 65.

Use exactly the time-dependent data of HA-3.3. Assume:

\begin{enumerate}
\def\labelenumi{\arabic{enumi}.}
\tightlist
\item
  \textbf{Attainment at every late time.} For every \(t\ge T\), an
  admissible \(u_t\) exists with
  \(\operatorname{Area}_{g_t}(u_t)=A(t)\).
\item
  \textbf{Local two-direction comparison.} For every \(t_0\ge T\) and
  every \(\epsilon>0\), there is \(\delta>0\) such that, whenever
  \(t\ge T\) and \(|t-t_0|<\delta\), both statements hold:

  \begin{itemize}
  \tightlist
  \item
    For \textbf{every} attaining admissible disk \(u\) at \(t_0\), an
    admissible disk \(v\) at \(t\) exists with
    \(\operatorname{Area}_{g_t}(v)\le e^\epsilon
      \operatorname{Area}_{g_{t_0}}(u)\).
  \item
    For \textbf{every} attaining admissible disk \(v\) at \(t\), an
    admissible disk \(u\) at \(t_0\) exists with
    \(\operatorname{Area}_{g_{t_0}}(u)\le e^\epsilon
      \operatorname{Area}_{g_t}(v)\).
  \end{itemize}
\end{enumerate}

Then \(A\) is continuous on \([T,\infty)\).

The comparison disks need not themselves attain the infimum. The radius
\(\delta\) may depend on \(t_0,\epsilon\), but works for all nearby
times and all minimizing disks in the two quantifiers. The statement
does not require comparison maps between whole manifolds, or a
continuously chosen minimizer. It does not supply the two comparisons
from surgery geometry; it assumes them. The blueprint uses
\(e^{2\epsilon}\), whereas the Lean statement uses \(e^\epsilon\);
universal quantification over all positive \(\epsilon\) makes these
formulations equivalent by renaming \(\epsilon/2\). No uniformity is
lost in that renaming.

The direct numerical consequence is
\(e^{-\epsilon}A(t_0)\le A(t)\le e^\epsilon A(t_0)\). This handles
\(A(t_0)=0\) as well, without taking a logarithm of zero. The theorem's
assertion is mathematically elementary once its strong comparison
premises are available; its Lean proof is still admitted.

\subsection{HA-4. What is, and is not, proved about torus
injectivity}\label{ha-4}

\subsubsection{HA-4.1. A finite time-dependent family of continuous
torus maps}\label{ha-4-1}

\textbf{Lean:}
{\small\texttt{Differential\allowbreak{}Geometry.\allowbreak{}Geometry.\allowbreak{}Ricci\allowbreak{}Flow.\allowbreak{}cusp\_\allowbreak{}injective\_\allowbreak{}of\_\allowbreak{}exterior\_\allowbreak{}area\_\allowbreak{}barriers}}.
\textbf{Status:} written proof, admission-dependent through HA-2.2.
Source:
{\small\texttt{Geometry/\allowbreak{}Flow/\allowbreak{}Ricci\allowbreak{}Flow/\allowbreak{}Long\allowbreak{}Time/\allowbreak{}Cusp\allowbreak{}Incompressibility.\allowbreak{}lean}},
line 15.

Let \(M_t,g_t\) be any family as in HA-3.3, let \(n\in\mathbb N\), and
let \(T_0\in\mathbb R\). For each \(t\ge T_0\) and each
\(i\in\{0,\ldots,n-1\}\), suppose a continuous map
\(f_{t,i}:T^2\to M_t\) is given. Assume the following implication for
\textbf{each} \(t\ge T_0\), each \(i\), and each basepoint \(x\in T^2\):

\begin{quote}
If \((f_{t,i})_*:\pi_1(T^2,x)\to\pi_1(M_t,f_{t,i}(x))\) is not
injective, there exist \(T\ge0\) with \(T\ge t\), subsets \(W_s\), and
loops \(\gamma_s\) for all \(s\ge T\), such that the associated actual
disk-area function \(A\) of HA-3.3 has an admissible attaining disk at
every \(s\ge T\), is continuous on \([T,\infty)\), and has at every
\(s\ge T\) a strict local smooth upper barrier with derivative less than
\(3A(s)/[4(s+1/4)]-\pi\).
\end{quote}

Then all the displayed homomorphisms \((f_{t,i})_*\) are injective.

The proof takes a hypothetical noninjectivity, obtains the promised area
function, proves its nonnegativity from the definition, and applies
HA-2.3. It does not use the attainment conjunct; attainment is retained
as a geometric requirement of the premise. Nor does the contradiction
need \(t\le T\) once the premise has supplied a half-line with
\(T\ge0\).

\textbf{Exact scope:} the maps \(f_{t,i}\) need not be smooth, embedded,
time-continuous, or actual cusp maps. No Ricci flow is assumed. There is
no condition in this theorem relating \(\gamma_s\) to a particular
primitive kernel element of \(f_{t,i}\), or \(W_s\) to its exterior.
Such relations must be used inside the proof that produces the displayed
implication. If \(n=0\), both that implication and the conclusion are
vacuous.

\subsubsection{HA-4.2. Moving the torus along a product
collar}\label{ha-4-2}

\textbf{Lean:}
{\small\texttt{Differential\allowbreak{}Geometry.\allowbreak{}Geometry.\allowbreak{}Ricci\allowbreak{}Flow.\allowbreak{}cusp\_\allowbreak{}slice\_\allowbreak{}injective\_\allowbreak{}iff}}.
\textbf{Status:} written proof, no skeleton admission. Source:
{\small\texttt{Geometry/\allowbreak{}Flow/\allowbreak{}Ricci\allowbreak{}Flow/\allowbreak{}Long\allowbreak{}Time/\allowbreak{}Cusp\allowbreak{}Incompressibility.\allowbreak{}lean}},
line 43.

For any topological space \(M\), any continuous
\(c:T^2\times\mathbb R\to M\), any \(r,s\in\mathbb R\), and any
\(x\in T^2\), the map \(y\mapsto c(y,r)\) is injective on fundamental
groups at \(x\) if and only if \(y\mapsto c(y,s)\) is.

The proof uses the actual product homotopy and the existing theorem that
homotopic maps have equivalent injectivity properties, with the induced
basepoint track. No metric or smoothness is involved; \(c\) need not be
an embedding. It is not an inference from injection into an abstract
hyperbolic model to injection into a different ambient manifold.

\subsubsection{HA-4.3. Injectivity for the actual reconstruction
seams}\label{ha-4-3}

\textbf{Lean:}
{\small\texttt{Differential\allowbreak{}Geometry.\allowbreak{}Geometry.\allowbreak{}Ricci\allowbreak{}Flow.\allowbreak{}incompressible\_\allowbreak{}of\_\allowbreak{}area\_\allowbreak{}barriers}}.
\textbf{Status:} written proof, admission-dependent through HA-2.2.
Source:
{\small\texttt{Geometry/\allowbreak{}Flow/\allowbreak{}Ricci\allowbreak{}Flow/\allowbreak{}Long\allowbreak{}Time/\allowbreak{}Cusp\allowbreak{}Incompressibility.\allowbreak{}lean}},
line 51.

Let \(C\) be the existing compact carrier, \(G\) its actual paired-torus
gluing data, \(S\) a smooth assembly of that gluing, and
\(r:S.\mathrm{assembled}\to P\) an orientation-preserving smooth
reconstruction onto a connected closed oriented 3-manifold \(P\).
Despite the code name
{\small\texttt{torus\allowbreak{}In\allowbreak{}Prime}}, this type
\textbf{does not assume that \(P\) is prime}. For every gluing label
\(i\), the relevant map is the actual quotient seam followed by \(r\),
denoted \(j_i:T^2\to P\).

Assume for every \(i\) and every basepoint \(x\), failure of injectivity
of \((j_i)_*\) produces real data \(T\ge0\) and
\(A:\mathbb R\to\mathbb R\) which are continuous and nonnegative on
\([T,\infty)\), with a strict smooth local upper barrier at each
\(t\ge T\) having derivative less than \(3A(t)/[4(t+1/4)]-\pi\). Then
{\small\texttt{S.\allowbreak{}Incompressible r}} holds: every actual
seam map is injective on fundamental groups at every basepoint.

This adapter does not itself require that \(A\) be an area, or that it
be attached to a flow. It is a numerical contradiction consumer whose
conclusion is precisely the existing reconstruction's injectivity
predicate.

\textbf{Shifted companion:}
{\small\texttt{Differential\allowbreak{}Geometry.\allowbreak{}Geometry.\allowbreak{}Ricci\allowbreak{}Flow.\allowbreak{}incompressible\_\allowbreak{}of\_\allowbreak{}shifted\_\allowbreak{}area\_\allowbreak{}barriers}},
line 67, has exactly the same reconstruction data and conclusion. Its
noninjectivity premise produces \(T\ge0\), \(c>0\), and \(A\), with the
bound \(3A(t)/[4(t+c)]-\pi\). The proof derives \(T+c>0\) and uses
HA-2.3. The shift can depend on the failed seam and basepoint. This
theorem also has a written admission-dependent proof.

\subsubsection{HA-4.4. Approval issue: the consumers do not expose the
producer's substance}\label{ha-4-4}

After granting the area contradiction, each ``noninjectivity produces an
impossible area function'' premise above is logically equivalent to the
corresponding injectivity conclusion: the forward implication is the
written proof, while if injectivity already holds, the premise follows
vacuously. This does \textbf{not} make the conditional proof false. It
does mean these consumers by themselves provide almost no independent
check of the geometric proof of incompressibility.

The broader selected-flow producer in
{\small\texttt{Late\allowbreak{}Decomposition.\allowbreak{}lean}} uses
the actual same-flow disk area defined below. That rules out an
unrelated numerical function, but does not yet expose, as separately
checked fields, the fixed primitive meridian, its preservation through
surgery, actual persistent cusp embeddings, mean-convex exterior
geometry, or the common unchanged carrier and boundary-moving isotopy.
Mathematicians should review that producer as a large remaining theorem,
not treat the short
{\small\texttt{.\allowbreak{}.\allowbreak{}.\allowbreak{}of\_\allowbreak{}area\_\allowbreak{}barriers}}
proofs as having formalized Meeks--Yau or surgery avoidance.

\subsection{HA-5. The same-flow physical area and the event-time
convention}\label{ha-5}

All seven declarations in this section are in
{\small\texttt{Geometry/\allowbreak{}Flow/\allowbreak{}Ricci\allowbreak{}Flow/\allowbreak{}Long\allowbreak{}Time/\allowbreak{}Exterior\allowbreak{}Disk\allowbreak{}Flow.\allowbreak{}lean}}.
None has {\small\texttt{sorry\allowbreak{}Ax}} in its recorded
dependency closure. They take an existing
{\small\texttt{Observation\allowbreak{}Tower P g}}, a compatible
collection of actual finite surgery histories with specified initial
carrier \(P\) and initial metric \(g\). This section does not construct
such a tower or strengthen its analytic control.

\subsubsection{HA-5.1. The carrier at every real time}\label{ha-5-1}

\textbf{Lean:}
{\small\texttt{GC.\allowbreak{}Long\allowbreak{}Time.\allowbreak{}post\allowbreak{}Stage}}.
\textbf{Status:} definition. Line 18.

For \(t\in\mathbb R\), put \(\tau=\max\{t,0\}\), restrict the tower to
its observation through \(\tau\), and take that observation's last
stage. This defines \(M_t^{\mathrm{post}}\), including its existing
orientation and manifold structure. For \(t\ge0\), this is the actual
carrier at time \(t\). At an exact surgery time it is the
\textbf{post-surgery} stage. For \(t<0\), it is the time-zero stage;
this is a totalization convention, not a backward flow.

This event convention follows from the inherited restriction rule: the
active stage is the largest stage index whose start time is \(\le t\),
not \(<t\). The last stage of the observation is that active stage at
the observation horizon. A stage may be empty.

\subsubsection{HA-5.2. The physical metric on that
carrier}\label{ha-5-2}

\textbf{Lean:}
{\small\texttt{GC.\allowbreak{}Long\allowbreak{}Time.\allowbreak{}post\allowbreak{}Metric}}.
\textbf{Status:} definition. Line 22.

On {\small\texttt{post\allowbreak{}Stage O t}}, take the last stage's
actual {\small\texttt{stage\allowbreak{}Metric}} evaluated at the same
\(\tau=\max\{t,0\}\). At a surgery time this is the post-surgery output
metric. This is the \textbf{physical metric}. It is neither
\(t^{-1}g(t)\) nor \((4t)^{-1}g(t)\). No value is divided by zero at
\(t=0\).

\subsubsection{HA-5.3. Agreement with a regular slice}\label{ha-5-3}

Two separate authored theorems have written proofs without skeleton
admissions:

\begin{itemize}
\tightlist
\item
  {\small\texttt{GC.\allowbreak{}Long\allowbreak{}Time.\allowbreak{}post\allowbreak{}Stage\_\allowbreak{}regular\allowbreak{}Slice}}
  (line 26): for any {\small\texttt{Regular\allowbreak{}Slice O}} with
  time \(s.time>0\), the carrier
  {\small\texttt{post\allowbreak{}Stage O s.\allowbreak{}time}} equals
  the stage stored by the regular slice.
\item
  {\small\texttt{GC.\allowbreak{}Long\allowbreak{}Time.\allowbreak{}post\allowbreak{}Metric\_\allowbreak{}regular\allowbreak{}Slice}}
  (line 37): the physical metric
  {\small\texttt{post\allowbreak{}Metric O s.\allowbreak{}time}} is the
  same metric as {\small\texttt{s.\allowbreak{}metric}}, after
  accounting for that equality of the carriers. The Lean conclusion is
  heterogeneous equality because the metric types mention their
  respective carriers.
\end{itemize}

These are exact identifications of two descriptions of the same
observation, not a diffeomorphism between different surgery slices or a
metric-comparison estimate across surgery.

\subsubsection{HA-5.4. The actual same-flow area function}\label{ha-5-4}

\textbf{Lean:}
{\small\texttt{GC.\allowbreak{}Long\allowbreak{}Time.\allowbreak{}exterior\allowbreak{}Disk\allowbreak{}Area}}.
\textbf{Status:} definition. Line 57.

Choose, for every real time, a subset
\(W_t\subseteq M_t^{\mathrm{post}}\), a starting time \(T\in\mathbb R\),
and a loop \(\gamma_t\) in that carrier for every \(t\ge T\). Define \[
 A_O(t)=\begin{cases}
 a(\mathrm{postMetric}(O,t),W_t,\gamma_t),&t\ge T,\\
 0,&t<T.
 \end{cases}
\] This is HA-3.3 applied to the actual observation tower. The
definition itself still permits arbitrary \(W_t\) and \(\gamma_t\),
imposes no continuity or persistence on them, and does not require
\(T\ge0\). Its intended root application chooses \(T\) after a positive
regular slice, so only actual nonnegative flow times are used in the
area contradiction.

Two authored written helpers, without skeleton admissions, are:

\begin{itemize}
\tightlist
\item
  {\small\texttt{GC.\allowbreak{}Long\allowbreak{}Time.\allowbreak{}exterior\allowbreak{}Disk\allowbreak{}Area\_\allowbreak{}nonneg}}
  (line 62): \(A_O(t)\ge0\) for all real \(t\).
\item
  {\small\texttt{GC.\allowbreak{}Long\allowbreak{}Time.\allowbreak{}exterior\allowbreak{}Disk\allowbreak{}Area\_\allowbreak{}eq}}
  (line 68): if \(t\ge T\),
  \(A_O(t)=a(\mathrm{postMetric}(O,t),W_t,\gamma_t)\), exactly.
\end{itemize}

They provide no assertion of attainment, positive area, continuity,
smooth upper barriers, or identification of any \(W_t\) with a specified
thin complement. At a surgery contact, the future producer must use the
actual post-surgery infimum and prove the two-sided numerical barrier
using a common unchanged region. These helpers do not assert a global
identification of pre- and post-surgery manifolds.

\subsection{HA-6. Mathematical review findings and decisions
needed}\label{ha-6}

This reading found \textbf{no concrete counterexample to the four
admitted leaf statements} in these six files. Their principal
formulations are compatible with the corresponding natural-language
arguments, subject to the explicit limitations above. That conclusion is
a targeted mathematical review, not proof verification of admitted
claims.

The following distinctions should remain visible when approving the
skeleton:

\begin{longtable}[]{@{}
  >{\raggedright\arraybackslash}p{(\linewidth - 4\tabcolsep) * \real{0.3333}}
  >{\raggedright\arraybackslash}p{(\linewidth - 4\tabcolsep) * \real{0.3333}}
  >{\raggedright\arraybackslash}p{(\linewidth - 4\tabcolsep) * \real{0.3333}}@{}}
\toprule\noalign{}
\begin{minipage}[b]{\linewidth}\raggedright
Point for review
\end{minipage} & \begin{minipage}[b]{\linewidth}\raggedright
What the code actually does
\end{minipage} & \begin{minipage}[b]{\linewidth}\raggedright
What remains outside these six files
\end{minipage} \\
\midrule\noalign{}
\endhead
\bottomrule\noalign{}
\endlastfoot
Mostow--Prasad & Admits the finite-volume, same-curvature, marked
homotopy-class theorem in dimension three. & Its proof, quantitative
rigidity, controlled persistent families, and cusp coverage. \\
Hyperbolic geometric metric & Uses \(g/4\) on the same whole carrier
when \(\sec g=-1/4\). & A full-interior map from each actual cut piece
to the complete hyperbolic carrier. \\
Fixed-boundary disks & Uses embedded, unbranched smooth disks with
exactly parametrized boundary, on the specified side. & Actual existence
in the chosen exterior; source-to-class comparison; minimality and
stability. \\
Empty disk class & Gives real infimum zero, and proves nonnegativity. &
Attainment must rule out emptiness wherever geometry is claimed. \\
Continuity through surgery & Admits continuity from explicit forward and
backward minimizing-disk comparisons. & Common compact carrier, its
two-sided smooth metric, and actual boundary-preserving transports. \\
Area barrier & Requires a strict smooth upper barrier at every time of
one entire half-line. & Scalar bound, boundary estimates, first
variation, Gauss--Bonnet, and surgery-neck avoidance in the actual
flow. \\
Torus injection & Has genuine contradiction consumers for the actual
maps supplied to them. & Producer premises currently encapsulate the
essential geometric argument. \\
Empty families & Permits zero tori, with vacuous injection obligations;
permits empty stages in the tower helpers. & No fictitious loop is
required before its threshold; late loops themselves force nonempty
carriers where used. \\
\end{longtable}

In particular, there is no new standalone Lean theorem here saying
``Meeks--Yau holds for the actual exterior,'' ``every least-area disk
avoids the surgery necks,'' or ``the chosen meridian is preserved
through every event.'' Those would be natural future statements when the
large root producer is split into independently assignable tasks.
Approving the present statement layer does not approve a completed proof
of those facts.

\subsection{HA-7. Evidence for this translation}\label{ha-7}

Freshly inspected for this report:

\begin{itemize}
\tightlist
\item
  The complete six Lean source files, all 28 authored declarations, and
  their full elaborated signatures and recorded axiom closures in
  {\small\texttt{evidence/\allowbreak{}declarations.\allowbreak{}json}}.
\item
  The inherited definitions
  {\small\texttt{Smooth\allowbreak{}Disk\allowbreak{}Extension}},
  {\small\texttt{free\allowbreak{}Loop}},
  {\small\texttt{closed\allowbreak{}Disk}},
  {\small\texttt{disk\allowbreak{}Trace}},
  {\small\texttt{riemannian\allowbreak{}Disk\allowbreak{}Area}},
  {\small\texttt{riemannian\allowbreak{}Area\allowbreak{}Density}}, and
  {\small\texttt{tangent\allowbreak{}Two\allowbreak{}Jacobian}}; the
  existing extension-invariance and compact integrability statements for
  disk area.
\item
  {\small\texttt{Riemannian\allowbreak{}Metric\allowbreak{}Complete}},
  {\small\texttt{Has\allowbreak{}Thurston\allowbreak{}Atlas}},
  {\small\texttt{Coordinate\allowbreak{}Model\allowbreak{}Atlas}}, the
  hyperbolic {\small\texttt{coordinate\allowbreak{}Inner}}, and
  {\small\texttt{Geometric\allowbreak{}Structure}}.
\item
  {\small\texttt{Smooth\allowbreak{}Assembly.\allowbreak{}Reconstruction}},
  {\small\texttt{torus\allowbreak{}In\allowbreak{}Prime}}, and
  {\small\texttt{Incompressible}};
  {\small\texttt{Observation\allowbreak{}Tower.\allowbreak{}observe}},
  {\small\texttt{Observed\allowbreak{}History.\allowbreak{}active\allowbreak{}Stage}},
  and the exact-event and observation-horizon rules.
\item
  Blueprint {\small\texttt{master207A.\allowbreak{}tex}}, MPR79 lines
  15392--15424; IMS02--IMS03 around lines 18251--18338; IMS08--IMS09
  around lines 18533--18607; IAU02--IAU03 around lines 18724--18855.
  These were used to compare the actual predicates, not to substitute
  prose for code.
\item
  Archived
  {\small\texttt{Books\allowbreak{}Papers/\allowbreak{}MSM206.\allowbreak{}tex}},
  {\small\texttt{Mos\allowbreak{}Tow\allowbreak{}Section}}, lines
  17115--17166, for the Mostow statement and its homotopy-equivalence
  formulation. Its shorthand based-group equation is not copied into the
  Lean theorem.
\end{itemize}

The prior source review is
{\small\texttt{hyperbolic\_\allowbreak{}review.\allowbreak{}md}}, with
exact source identities in
{\small\texttt{hyperbolic\_\allowbreak{}crosswalk.\allowbreak{}json}}.
This report reuses unchanged source/errata checks only at their
documented scope; it does not claim a fresh literature or errata search.
It does not rerun the Lean build. Compilation and the axiom
classification above are read from the final recorded verification of
the pinned commit, supplemented by direct source inspection. The
machine-readable
{\small\texttt{hyperbolic\_\allowbreak{}area\_\allowbreak{}coverage.\allowbreak{}json}}
maps every declaration to the subsection that translates it.

\section{Mathematical reading of the topology and finite-regularity
skeleton}\label{topology-finite-regularity}

This report translates the actual Lean source at commit
{\small\texttt{ea0fae6\allowbreak{}0ee01ef8\allowbreak{}c8d9b57a\allowbreak{}51794c65\allowbreak{}38f679c5\allowbreak{}d}}.
It covers all 32 authored declarations in the five files listed below.
It is a reading of what the code says, including its weaker conclusions
and omitted data; it does not replace those statements by the stronger
blueprint assertions.

There are \textbf{five admitted theorems} in this part: three topology
producers, one finite-order compactness theorem, and one
finite-regularity pullback theorem. Seven further theorem bodies contain
actual proofs: three elementary collar facts independent of these
admissions, two corollaries of the admitted pullback theorem, and two
assemblies of admitted topology producers. The remaining 20 authored
declarations define mathematical objects or construct data. Anonymous
instances and compiler-generated structure projections are not separate
authored declarations in these counts.

The five source files are:

\begin{longtable}[]{@{}
  >{\raggedright\arraybackslash}p{(\linewidth - 2\tabcolsep) * \real{0.5000}}
  >{\raggedright\arraybackslash}p{(\linewidth - 2\tabcolsep) * \real{0.5000}}@{}}
\toprule\noalign{}
\begin{minipage}[b]{\linewidth}\raggedright
Short name in this report
\end{minipage} & \begin{minipage}[b]{\linewidth}\raggedright
Actual source
\end{minipage} \\
\midrule\noalign{}
\endhead
\bottomrule\noalign{}
\endlastfoot
Presentation &
\href{https://github.com/qinz1yang/differential-geometry-dev/blob/ea0fae60ee01ef8c8d9b57a51794c6538f679c5d/DifferentialGeometry/Topology/ThreeManifold/GraphManifold/Presentation.lean}{GraphManifold/Presentation.lean} \\
Refinement &
\href{https://github.com/qinz1yang/differential-geometry-dev/blob/ea0fae60ee01ef8c8d9b57a51794c6538f679c5d/DifferentialGeometry/Topology/ThreeManifold/GraphManifold/Refinement.lean}{GraphManifold/Refinement.lean} \\
TorusDecomposition &
\href{https://github.com/qinz1yang/differential-geometry-dev/blob/ea0fae60ee01ef8c8d9b57a51794c6538f679c5d/DifferentialGeometry/Topology/ThreeManifold/TorusCut/Decomposition.lean}{TorusCut/Decomposition.lean} \\
Pullback &
\href{https://github.com/qinz1yang/differential-geometry-dev/blob/ea0fae60ee01ef8c8d9b57a51794c6538f679c5d/DifferentialGeometry/Geometry/Metric/Pullback/FiniteRegularity.lean}{Metric/Pullback/FiniteRegularity.lean} \\
Compactness &
\href{https://github.com/qinz1yang/differential-geometry-dev/blob/ea0fae60ee01ef8c8d9b57a51794c6538f679c5d/DifferentialGeometry/Analysis/Calculus/Compactness/FiniteOrder.lean}{Calculus/Compactness/FiniteOrder.lean} \\
\end{longtable}

\subsection{T0. The mathematical types used throughout}\label{t0}

The carrier of a compact piece is an \textbf{actual compact, Hausdorff,
second-countable, oriented smooth 3-manifold}, possibly with boundary.
It is not a homeomorphism type, a fundamental group, or a model name.
Its chart model is either all of Euclidean 3-space or a Euclidean
half-space. A half-space chart model permits a manifold whose boundary
happens to be empty. The carrier type by itself need not be connected or
nonempty.

{\small\texttt{Compact\allowbreak{}Carrier.\allowbreak{}Components}}
adds a finite decomposition into actual subsets of that carrier. More
precisely, it supplies a positive integer \(n>0\), open-and-closed
subsets \(C_i\), \(0\leq i<n\), which are pairwise disjoint and cover
the carrier, with each \(C_i\) connected \textbf{and nonempty}. Each
\(\operatorname{int}(C_i)=C_i\cap\operatorname{int}(C)\) is also
required to be connected and nonempty. In Lean,
{\small\texttt{Connected\allowbreak{}Space}} includes nonemptiness. Thus
this component record cannot describe an empty carrier or an empty
component list. Individual torus-index lists below may nevertheless be
empty.

These inherited definitions are in
\href{https://github.com/qinz1yang/differential-geometry-dev/blob/ea0fae60ee01ef8c8d9b57a51794c6538f679c5d/DifferentialGeometry/Topology/ThreeManifold/Geometrization/Carrier.lean\#L32}{Carrier.lean:32}
and
\href{https://github.com/qinz1yang/differential-geometry-dev/blob/ea0fae60ee01ef8c8d9b57a51794c6538f679c5d/DifferentialGeometry/Topology/ThreeManifold/Geometrization/Carrier.lean\#L56}{Carrier.lean:56}.

An {\small\texttt{Interior\allowbreak{}Geometry}} on \(C_i\) means a
smooth Riemannian metric on \textbf{the whole actual interior}
\(C_i\cap\operatorname{int}(C)\), complete in its own Riemannian
distance, together with local isometric charts for one specified member
of the eight fixed Thurston models. If the model is hyperbolic, its
total volume is finite. For the other seven models, finite volume is
\textbf{not} required. This is the existing endpoint convention. It
permits, for example, complete infinite-volume flat product interiors.
It does not mean merely ``there is some locally homogeneous metric on
some abstract manifold related to this component.'' See
\href{https://github.com/qinz1yang/differential-geometry-dev/blob/ea0fae60ee01ef8c8d9b57a51794c6538f679c5d/DifferentialGeometry/Geometry/Thurston/Atlas.lean\#L56}{Atlas.lean:56}.

For a connected closed oriented smooth 3-manifold \(M\), a
{\small\texttt{Prime\allowbreak{}Decomposition M}} supplies a nonempty
finite list of actual connected closed oriented manifolds \(P_i\), each
prime, and an orientation-preserving smooth diffeomorphism

\[
       P_1\#\cdots\#P_m\;\longrightarrow\;M.
\]

Here prime means that every smooth connected-sum decomposition has a
summand diffeomorphic to \(S^3\). Primeness is formulated without an
orientation requirement on that summand-identifying diffeomorphism. In
particular, \(S^3\) is allowed as a recorded prime identity, and
\(S^2\times S^1\) is not excluded. The record does not itself supply an
embedded sphere system or cap balls in \(M\). See
\href{https://github.com/qinz1yang/differential-geometry-dev/blob/ea0fae60ee01ef8c8d9b57a51794c6538f679c5d/DifferentialGeometry/Topology/ThreeManifold/Geometrization/Prime.lean\#L11}{Prime.lean:11}.

A {\small\texttt{Geometric\allowbreak{}Decomposition P}} consists of
actual compact cut components, paired torus boundaries, their smooth
oriented reconstruction to \(P\), injection of each reconstructed torus
into \(\pi_1(P)\), and an {\small\texttt{Interior\allowbreak{}Geometry}}
for each cut component. The type can be formed for any connected closed
\(P\); it does not itself assume that \(P\) is prime. A
{\small\texttt{Geometrization\allowbreak{}Certificate M}} additionally
requires its chosen factors \(P_i\) to be prime. The proposition
{\small\texttt{Geometrizes M}} means that such a certificate exists. See
\href{https://github.com/qinz1yang/differential-geometry-dev/blob/ea0fae60ee01ef8c8d9b57a51794c6538f679c5d/DifferentialGeometry/Topology/ThreeManifold/Geometrization/Statement.lean\#L11}{Statement.lean:11}.

\subsection{T1. Surfaces and ordinary circle bundles}\label{t1}

\textbf{Declarations:} {\small\texttt{Surface\allowbreak{}Model}}
(Presentation:10),
{\small\texttt{Surface\allowbreak{}Model.\allowbreak{}Space}} (:13),
{\small\texttt{Surface\allowbreak{}Model.\allowbreak{}model}} (:22),
{\small\texttt{Compact\allowbreak{}Surface}} (:28), and
{\small\texttt{Circle\allowbreak{}Fibration}} (:43). All five define
data; none asserts an existence theorem.

{\small\texttt{Surface\allowbreak{}Model}} has two choices: the plane
and the closed half-plane. The {\small\texttt{Space}} and
{\small\texttt{model}} declarations associate their usual smooth
manifold models to those choices.
{\small\texttt{Compact\allowbreak{}Surface}} is a compact, connected,
nonempty, Hausdorff, second-countable smooth surface with one of these
models. \textbf{There is no assumption that the surface is orientable.}
The half-plane choice does not require nonempty boundary.

For an open subset \(U\subset C\), a
{\small\texttt{Circle\allowbreak{}Fibration C U}} supplies:

\begin{enumerate}
\def\labelenumi{\arabic{enumi}.}
\tightlist
\item
  A compact surface \(B\) as just defined.
\item
  A continuous, smooth, surjective map \(\pi:U\to B\).
\item
  For every \(b\in B\), an open neighborhood \(V_b\ni b\) and an actual
  smooth diffeomorphism \[
       \tau_b:\pi^{-1}(V_b)\longrightarrow V_b\times S^1
  \] whose first coordinate is exactly \(\pi\).
\end{enumerate}

The domain in item 3 is the inverse image \textbf{inside the actual
subset \(U\)}. This is an ordinary smooth circle bundle with local
product trivializations. Every fiber is a smooth circle. The record does
not require principal \(S^1\)-bundle transition functions, a global
circle action, a chosen fiber orientation, or an orientable base.
Arbitrary smooth fiberwise changes of circle coordinates are allowed. It
also has no exceptional fibers or orbifold base. Thus it represents the
regular circle-bundle blocks of the raw graph presentation, not a
definition of every Seifert fibration.

In the later
{\small\texttt{Raw\allowbreak{}Graph\allowbreak{}Presentation}}, \(U\)
is one of the compact clopen connected cut components. The generic
definition of {\small\texttt{Circle\allowbreak{}Fibration C U}} itself
does not separately assume that \(U\) is compact or connected.

\subsection{T2. Actual labelled boundary tori and their elementary
properties}\label{t2}

\textbf{Declaration:} {\small\texttt{Boundary\allowbreak{}Tori C n}}
(Presentation:55), a definition of data.

For each label \(i\in\{0,\ldots,n-1\}\), the record supplies a smooth
collar diffeomorphism from \(T^2\times[0,1)\) onto an open subset of
\(C\). Technically the object is a partial diffeomorphism defined on
ambient \(T^2\times[0,\infty)\), with designated source exactly
\(T^2\times[0,1)\). Only its restriction to the source carries the
asserted inverse and regularity properties. The zero slice lies in the
intrinsic boundary of \(C\). The target open sets of distinct collars
are disjoint.

This record \textbf{does not by itself say that these tori exhaust the
boundary}, or that they are incompressible. Boundary exhaustion is a
separate field of
{\small\texttt{Raw\allowbreak{}Graph\allowbreak{}Presentation}}. The
number \(n\) may be zero.

The following seven declarations make the convention precise:

\begin{longtable}[]{@{}
  >{\raggedright\arraybackslash}p{(\linewidth - 4\tabcolsep) * \real{0.3333}}
  >{\raggedright\arraybackslash}p{(\linewidth - 4\tabcolsep) * \real{0.3333}}
  >{\raggedright\arraybackslash}p{(\linewidth - 4\tabcolsep) * \real{0.3333}}@{}}
\toprule\noalign{}
\begin{minipage}[b]{\linewidth}\raggedright
Declaration and location
\end{minipage} & \begin{minipage}[b]{\linewidth}\raggedright
Exact mathematical content
\end{minipage} & \begin{minipage}[b]{\linewidth}\raggedright
Status
\end{minipage} \\
\midrule\noalign{}
\endhead
\bottomrule\noalign{}
\endlastfoot
{\small\texttt{Boundary\allowbreak{}Tori.\allowbreak{}torus\allowbreak{}Map}},
Presentation:64 & \(j_i(t)=\operatorname{collar}_i(t,0)\). &
Definition \\
{\small\texttt{Boundary\allowbreak{}Tori.\allowbreak{}image}}, :67 & The
union \(\bigcup_i j_i(T^2)\). & Definition \\
{\small\texttt{Boundary\allowbreak{}Tori.\allowbreak{}zero\_\allowbreak{}mem\_\allowbreak{}source}},
:70 & Every \((t,0)\) is in the designated collar source because
\(0<1\). This helper is private. & Proved, independent of admissions \\
{\small\texttt{Boundary\allowbreak{}Tori.\allowbreak{}torus\allowbreak{}Map\_\allowbreak{}smooth}},
:77 & Each \(j_i:T^2\to C\) is smooth. & Proved, independent of
admissions \\
{\small\texttt{Boundary\allowbreak{}Tori.\allowbreak{}boundary\allowbreak{}Map}},
:82 & The same \(j_i\), now packaged as a continuous map for use with
\(\pi_1\). & Definition using the proved smoothness \\
{\small\texttt{Boundary\allowbreak{}Tori.\allowbreak{}torus\allowbreak{}Map\_\allowbreak{}is\allowbreak{}Embedding}},
:86 & Each \(j_i\) is a topological embedding. The proof uses collar
injectivity and compactness of the torus into the Hausdorff carrier. &
Proved, independent of admissions \\
{\small\texttt{Boundary\allowbreak{}Tori.\allowbreak{}incompressible}},
:93 & For every label \(i\) and \textbf{every} basepoint \(x\in T^2\),
the map \(\pi_1(T^2,x)\to\pi_1(C,j_i(x))\) induced by this actual
\(j_i\) is injective. & Definition, not an asserted theorem \\
\end{longtable}

There is no implicit distinction here between ``injects into a block''
and ``injects into an ambient manifold'': the target is exactly the
carrier \(C\) passed to this record. At \(n=0\), the last predicate is
vacuously true, as expected for an empty torus family.

\subsection{T3. Paired torus ports and the quotient they actually
define}\label{t3}

\textbf{Declarations:} {\small\texttt{Torus\allowbreak{}Pairing}}
(Presentation:98),
{\small\texttt{Torus\allowbreak{}Pairing.\allowbreak{}Quotient\allowbreak{}Space}}
(:118), and
{\small\texttt{Torus\allowbreak{}Pairing.\allowbreak{}quotient\allowbreak{}Map}}
(:124). These define data, not new existence theorems.

For a compact carrier \(C\),
{\small\texttt{Torus\allowbreak{}Pairing C}} supplies finitely many
pairs \((L_i,R_i)\) of closed subsets of \(C\). Each is parametrized by
\(T^2\). All the subsets are disjoint, including
\(L_i\cap R_i=\varnothing\), and there is a homeomorphism
\(a_i:L_i\to R_i\). In the selected parametrizations, \(a_i\) is
represented by an actual smooth torus diffeomorphism \(m_i:T^2\to T^2\):

\[
       a_i(\ell_i(t))=r_i(m_i(t)).
\]

Both parametrizations are the zero slices of supplied smooth
half-collars, each with source \(T^2\times[0,1)\). The left and
matched-right collar derivatives induce opposite orientations from the
orientation on \(C\). Precisely, pull back the ambient three-dimensional
tangent orientation by the two collar derivatives at \((t,0)\); one
equals the negative of the other. This is the code's
boundary-orientation reversal condition. It is not an additional
assertion that \(m_i\) has a particular degree with respect to an
independently chosen torus orientation.

The generic pairing record does not itself contain boundary exhaustion.
In {\small\texttt{Raw\allowbreak{}Graph\allowbreak{}Presentation}},
boundary exhaustion forces every paired port to lie in the actual
boundary. It permits \(L_i\) and \(R_i\) to belong to the same connected
component of \(C\), although the subsets themselves must be distinct and
disjoint. Thus loop edges are permitted.

{\small\texttt{Quotient\allowbreak{}Space}} is the literal quotient that
identifies \(x\in L_i\) with \(a_i(x)\in R_i\), and identifies no other
distinct points. It is not a guessed assembled manifold. Disjointness of
all ports makes these pairwise identifications an equivalence relation.
The inherited compact Hausdorff quotient theorem supplies its Hausdorff
topology. {\small\texttt{quotient\allowbreak{}Map}} is the actual
continuous quotient map \(q:C\to C/{\sim}\).

The inherited gluing relation and its assumptions are in
\href{https://github.com/qinz1yang/differential-geometry-dev/blob/ea0fae60ee01ef8c8d9b57a51794c6538f679c5d/DifferentialGeometry/Topology/Attachment/BoundaryGluing.lean\#L154}{BoundaryGluing.lean:154}.
The exact orientation condition is in
\href{https://github.com/qinz1yang/differential-geometry-dev/blob/ea0fae60ee01ef8c8d9b57a51794c6538f679c5d/DifferentialGeometry/Topology/ThreeManifold/Geometrization/SmoothTorusReconstruction.lean\#L13}{SmoothTorusReconstruction.lean:13}.

\subsection{T4. What a raw graph presentation contains}\label{t4}

\textbf{Declaration:}
{\small\texttt{Raw\allowbreak{}Graph\allowbreak{}Presentation W}}
(Presentation:129). This is a definition of data, not a theorem that a
given \(W\) has such data.

It says that the actual compact oriented smooth carrier \(W\) is
obtained by gluing finitely many actual regular circle-bundle blocks
along actual paired torus boundaries, with all the following witnesses.

\begin{longtable}[]{@{}
  >{\raggedright\arraybackslash}p{(\linewidth - 2\tabcolsep) * \real{0.5000}}
  >{\raggedright\arraybackslash}p{(\linewidth - 2\tabcolsep) * \real{0.5000}}@{}}
\toprule\noalign{}
\begin{minipage}[b]{\linewidth}\raggedright
Lean field(s)
\end{minipage} & \begin{minipage}[b]{\linewidth}\raggedright
Mathematical data or requirement
\end{minipage} \\
\midrule\noalign{}
\endhead
\bottomrule\noalign{}
\endlastfoot
{\small\texttt{cut\allowbreak{}Carrier}}, {\small\texttt{components}} &
A compact smooth oriented cut carrier \(C\), with a positive finite
clopen decomposition \(C=\coprod_i C_i\) as in T0. \\
{\small\texttt{fibration}} & An ordinary circle-bundle structure on each
entire compact component \(C_i\), as in T1. \\
{\small\texttt{pairing}} & Paired internal ports and their smooth
matching maps and half-collars, as in T3. Write
\(q:C\to Q=C/{\sim}\). \\
{\small\texttt{external\allowbreak{}Count}}, {\small\texttt{external}},
{\small\texttt{cut\allowbreak{}External}} & A chosen finite label set
for external boundary tori, and disjoint smooth half-collars on both
\(W\) and \(C\), with those same labels. \\
{\small\texttt{external\_\allowbreak{}exhausted}} & The actual boundary
of \(W\) equals the union of its external zero-slice tori. \\
{\small\texttt{cut\_\allowbreak{}boundary\_\allowbreak{}exhausted}} &
The actual boundary of \(C\) equals the union of all paired internal
ports and all cut external zero-slice tori. \\
{\small\texttt{external\_\allowbreak{}disjoint}} & The paired internal
ports are disjoint from the cut external zero-slice tori. \\
{\small\texttt{reconstruction}} & An actual homeomorphism
\(h:Q\to W\). \\
{\small\texttt{quotient\_\allowbreak{}smooth}} & The map
\(F=h\circ q:C\to W\) is smooth, including in the manifold-with-boundary
sense at each boundary point. \\
{\small\texttt{quotient\_\allowbreak{}oriented}} & At every \(x\in C\),
\(dF_x\) is represented by a linear isomorphism carrying the chosen
tangent orientation of \(C\) to that of \(W\). Thus the half-side maps
have nonsingular orientation-preserving derivatives. \\
{\small\texttt{interior\allowbreak{}Image}},
{\small\texttt{interior\allowbreak{}Diffeomorph}},
{\small\texttt{interior\_\allowbreak{}map}} & An open subset
\(V\subset W\) and a smooth diffeomorphism \(C^\circ\to V\), whose
underlying map is exactly \(F\vert_{C^\circ}\). The interior here
removes \textbf{all} boundary of \(C\), paired and external. \\
{\small\texttt{seam}}, {\small\texttt{seam\_\allowbreak{}source}},
{\small\texttt{seam\_\allowbreak{}zero}} & For each pair, a smooth
two-sided seam chart \(s_i:T^2\times(-1,1)\to W\), onto an open set,
with \(s_i(t,0)=F(\ell_i(t))\). \\
{\small\texttt{seam\_\allowbreak{}positive}} & For \(0\leq u<1\),
\(s_i(t,u)=F(\text{rightCollar}_i(m_i(t),u))\). \\
{\small\texttt{seam\_\allowbreak{}negative}} & For \(-1<u\leq0\),
\(s_i(t,u)=F(\text{leftCollar}_i(t,-u))\). \\
{\small\texttt{seam\_\allowbreak{}interior}},
{\small\texttt{seam\_\allowbreak{}disjoint}} & Each seam chart has image
in \(W^\circ\), and these images are pairwise disjoint. \\
{\small\texttt{marked\_\allowbreak{}collar}} & On \textbf{every point
of} \(T^2\times[0,1)\), reconstruction of the cut external collar is
exactly the corresponding external collar of \(W\). This is equality of
parametrized collar maps, not just equality of homology classes. \\
{\small\texttt{external\_\allowbreak{}seam\_\allowbreak{}disjoint}} &
Every external collar image in \(W\) is disjoint from every internal
seam image. \\
{\small\texttt{left\allowbreak{}Piece}},
{\small\texttt{right\allowbreak{}Piece}},
{\small\texttt{left\_\allowbreak{}owned}},
{\small\texttt{right\_\allowbreak{}owned}} & Each paired port is
assigned to an actual connected cut component, and its subset is
contained in that component. The assignments may be equal for the two
sides of one pair. \\
{\small\texttt{external\allowbreak{}Piece}},
{\small\texttt{external\_\allowbreak{}owned}} & Each external zero-slice
torus is assigned to its actual connected cut component. \\
\end{longtable}

This definition has \textbf{no incompressibility hypothesis}. Its
auxiliary seams may be compressible; they need not be the final
essential tori. It supplies no irreducibility of a block and no
geometric metric. It supplies no circle-fiber compatibility between
opposite sides of an arbitrary torus pairing; all smooth
orientation-compatible torus matchings are allowed. The number of paired
or external tori can be zero. The blocks cannot form an empty list.

The external labels and collar parametrizations are genuinely retained
\textbf{inside the supplied presentation}. But
{\small\texttt{Raw\allowbreak{}Graph\allowbreak{}Presentation W}}
chooses them as part of its own data. A theorem merely concluding
{\small\texttt{Nonempty (Raw\allowbreak{}Graph\allowbreak{}Presentation W)}}
does not by that fact preserve an independently prescribed set of
external labels or parametrizations. A separate producer must state any
such comparison.

The code uses regular bundles, even when a later Seifert description has
exceptional fibers. To obtain a raw presentation from that later
description, one must actually remove exceptional-fiber neighborhoods
and use regular bundle pieces; an exceptional solid-torus carrier may
itself be given a different ordinary bundle structure over a disk. That
conversion is not a theorem in these five files.

\subsection{T5. A torus decomposition before geometry or
incompressibility}\label{t5}

\textbf{Declaration:} {\small\texttt{component\allowbreak{}Carrier}}
(TorusDecomposition:10), a definition.

Given the data \(C=\coprod_i C_i\) of T0 and an index \(i\), this is the
actual subset \(C_i\), equipped with its inherited smooth chart model
and the restriction of \(C\)'s orientation. Compactness comes from
\(C_i\) being closed in compact \(C\). No replacement by an abstract
diffeomorphic manifold occurs.

\textbf{Declaration:} {\small\texttt{Torus\allowbreak{}Decomposition M}}
(TorusDecomposition:17), a definition of data.

Here \(M\) is an actual connected, nonempty, closed, oriented smooth
3-manifold, \textbf{not assumed prime}. A torus decomposition supplies a
compact oriented cut carrier with its positive finite components, pairs
of torus ports exhausting its entire boundary, a smooth oriented
assembly of their literal quotient, and an orientation-preserving
diffeomorphism from the assembled quotient to \(M\). Every left and
right port is assigned to its actual component. The inherited assembly
includes the smooth quotient map, its orientation-preserving nonsingular
derivatives, an interior diffeomorphism, matching half-collars, and
two-sided seam charts with the exact positive/negative identities
described in T4.

Unlike T4, all boundary ports in this record are paired: the
reconstructed \(M\) is closed. Unlike
{\small\texttt{Raw\allowbreak{}Graph\allowbreak{}Presentation}}, this
definition asks for no circle fibrations on the pieces. Crucially,
\textbf{it also asks for neither incompressibility nor geometric
metrics}. Those are subsequent premises. There may be zero tori, and
there may be loop edges.

\textbf{Declaration:}
{\small\texttt{Torus\allowbreak{}Decomposition.\allowbreak{}component}}
(TorusDecomposition:30), a definition, is just the component carrier
above for this particular decomposition.

\textbf{Declaration:}
{\small\texttt{Torus\allowbreak{}Decomposition.\allowbreak{}to\allowbreak{}Geometric\allowbreak{}Decomposition}}
(TorusDecomposition:34), an actual data assembly without an admission.

If this same torus decomposition is supplied with (i) injection of every
reconstructed seam \(T^2\to M\) on \(\pi_1\), at every torus basepoint,
and (ii) a complete model geometry on each actual cut interior, then it
becomes a {\small\texttt{Geometric\allowbreak{}Decomposition M}} by
retaining every carrier, pairing, reconstruction map, and piece
assignment unchanged. This construction adds the supplied evidence; it
proves no new geometry or incompressibility. It is valid without
assuming \(M\) prime, because that assumption belongs to the surrounding
prime-decomposition certificate.

\subsection{T6. The first admitted topology producer: closed graph
carriers admit graph prime factors}\label{t6}

\textbf{Lean:}
{\small\texttt{GC.\allowbreak{}Graph\allowbreak{}Manifold.\allowbreak{}exists\_\allowbreak{}prime\_\allowbreak{}decomposition\_\allowbreak{}of\_\allowbreak{}raw\allowbreak{}Graph\allowbreak{}Presentation}},
\href{https://github.com/qinz1yang/differential-geometry-dev/blob/ea0fae60ee01ef8c8d9b57a51794c6538f679c5d/DifferentialGeometry/Topology/ThreeManifold/GraphManifold/Refinement.lean\#L23}{Refinement:23}.

\textbf{Status: admitted with {\small\texttt{sorry}}.}

For \textbf{every} connected closed oriented smooth 3-manifold \(M\),
and \textbf{every supplied} raw graph presentation of the actual \(M\),
there exists a {\small\texttt{Prime\allowbreak{}Decomposition M}} such
that \textbf{each of its chosen factors} also has a raw graph
presentation.

In symbols, with ``graph'' meaning existence of the full data in T4:

\[
\begin{gathered}
 \forall M\;\forall G\in\operatorname{RawGraphPresentation}(M),\\
 \exists\,(P_1,\ldots,P_m,\Phi),\quad
 m\geq1,\quad P_i\text{ prime and graph},\\
 \Phi:\#_iP_i\cong_+M.
\end{gathered}
\]

{\small\texttt{No\allowbreak{}Cuts.\allowbreak{}carrier M}} in the Lean
signature means \(M\) itself regarded as a compact carrier with no
boundary; it is not an extra topological assumption.

This is a \textbf{closed consequence} of the blueprint's relative
graph-prime construction. The conclusion preserves the actual selected
factor identities: it asserts graph presentations on precisely the
\(P_i\) in its reconstruction. But it gives no relation between the
original raw seams and the factor seams, no embedded sphere system in
\(M\), no punctured factors or cap-ball assignments, and no relative
boundary data. The theorem does not accept a preselected prime
decomposition which it must use.

It does not assume irreducibility of \(M\), nonempty torus cuts,
incompressibility of raw seams, or uniqueness of prime decomposition. It
permits \(S^3\) identity factors and \(S^2\times S^1\) factors under the
definition in T0.

\subsection{T7. The second admitted topology producer: a prime graph
carrier admits the endpoint geometry}\label{t7}

\textbf{Lean:}
{\small\texttt{GC.\allowbreak{}Graph\allowbreak{}Manifold.\allowbreak{}exists\_\allowbreak{}geometric\_\allowbreak{}decomposition\_\allowbreak{}of\_\allowbreak{}prime\_\allowbreak{}raw\allowbreak{}Graph\allowbreak{}Presentation}},
\href{https://github.com/qinz1yang/differential-geometry-dev/blob/ea0fae60ee01ef8c8d9b57a51794c6538f679c5d/DifferentialGeometry/Topology/ThreeManifold/GraphManifold/Refinement.lean\#L31}{Refinement:31}.

\textbf{Status: admitted with {\small\texttt{sorry}}.}

For every connected closed oriented smooth 3-manifold \(P\), if \(P\) is
prime and an actual raw graph presentation of \(P\) is supplied, then
\(P\) has a {\small\texttt{Geometric\allowbreak{}Decomposition}} as
defined in T0 and T5, including essential torus cuts and complete
metrics on all the actual cut interiors.

The final torus family and geometric models are chosen by the
conclusion. No statement identifies them with the raw auxiliary torus
family. Their actual reconstructed inclusions into \(P\) must be
injective on \(\pi_1\). The code imposes no minimality, canonical JSJ
property, nonparallel condition, prescribed fiber slope, or prohibition
on extra product cuts. A torus bundle that could support an uncut Sol
metric may therefore instead be represented by retained essential
product cuts with complete flat or Seifert-type interiors.

The output's model field may be any of the eight models. This theorem
does \textbf{not} additionally assert that hyperbolic pieces cannot
occur in its chosen output, even though a sharper graph-manifold
classification would provide such a statement. Whatever model is chosen
must have its actual model atlas and complete metric; a model tag alone
is insufficient. Only hyperbolic outputs have the additional
finite-volume requirement.

It also does not say that the complete metrics are restrictions of a
metric on \(P\), agree across a torus, or retain a prescribed boundary
metric. Completeness is required on the whole cut interior, and no
complete restriction is claimed after sphere puncturing.

\subsection{T8. The actual closed graph assembly}\label{t8}

\textbf{Lean:}
{\small\texttt{GC.\allowbreak{}Graph\allowbreak{}Manifold.\allowbreak{}geometrizes\_\allowbreak{}of\_\allowbreak{}raw\allowbreak{}Graph\allowbreak{}Presentation}},
\href{https://github.com/qinz1yang/differential-geometry-dev/blob/ea0fae60ee01ef8c8d9b57a51794c6538f679c5d/DifferentialGeometry/Topology/ThreeManifold/GraphManifold/Refinement.lean\#L37}{Refinement:37}.

\textbf{Status: actual proof, depending on the two admissions T6 and
T7.}

Given a raw graph presentation of a connected closed oriented smooth
\(M\), the proof obtains the selected prime decomposition from T6,
obtains geometric decomposition data on each of those same factors from
T7, and puts them into
{\small\texttt{Geometrization\allowbreak{}Certificate M}}. It uses
choice to select the finitely indexed witnesses from existence
statements. It does not prove the two producers.

This is a real, compiling connection between supplier and consumer: the
graph presentation returned for \(P_i\) is applied to that exact
\(P_i\), using the primality proof of that same factor. It does not
substitute merely an isomorphic fundamental group or an unmarked list of
prime types.

\subsection{T9. Hyperbolic-or-graph pieces and the mixed admitted
topology producer}\label{t9}

\textbf{Declarations:}
{\small\texttt{is\allowbreak{}Hyperbolic\allowbreak{}Interior\allowbreak{}Geometry}}
(Refinement:11) and
{\small\texttt{Hyperbolic\allowbreak{}Or\allowbreak{}Graph}} (:17), both
definitions.

{\small\texttt{is\allowbreak{}Hyperbolic\allowbreak{}Interior\allowbreak{}Geometry g}}
means that the {\small\texttt{Interior\allowbreak{}Geometry}} \(g\) of
T0 has its model field equal to the fixed hyperbolic model. The rest of
the mathematical force lies in the already supplied
{\small\texttt{Interior\allowbreak{}Geometry}}: completeness, an actual
hyperbolic model atlas, and finite volume.

For an actual cut component \(C_i\subset C\),
{\small\texttt{Hyperbolic\allowbreak{}Or\allowbreak{}Graph C components i}}
is data of \textbf{one of} the following kinds:

\begin{enumerate}
\def\labelenumi{\arabic{enumi}.}
\tightlist
\item
  A complete finite-volume hyperbolic structure on the actual whole
  interior \(C_i\cap C^\circ\).
\item
  A full raw graph presentation of the actual compact carrier \(C_i\).
\end{enumerate}

No metric is required in the graph alternative. The alternatives are
constructors of a data type, not a theorem that every component has one
of these forms. No relation to an ambient Ricci-flow metric is imposed
by this particular type; those metric relations belong to the separate
collapse/flow input layer.

\textbf{Lean:}
{\small\texttt{GC.\allowbreak{}Graph\allowbreak{}Manifold.\allowbreak{}exists\_\allowbreak{}prime\_\allowbreak{}geometric\_\allowbreak{}decomposition\_\allowbreak{}of\_\allowbreak{}hyperbolic\allowbreak{}Or\allowbreak{}Graph}},
\href{https://github.com/qinz1yang/differential-geometry-dev/blob/ea0fae60ee01ef8c8d9b57a51794c6538f679c5d/DifferentialGeometry/Topology/ThreeManifold/GraphManifold/Refinement.lean\#L46}{Refinement:46}.

\textbf{Status: admitted with {\small\texttt{sorry}}.}

For every connected closed oriented smooth 3-manifold \(M\), every
supplied actual torus decomposition \(D\) of \(M\), and every choice of
hyperbolic-or-graph data for each of its actual cut components, assume
additionally that \textbf{each of the input seam maps into this actual
\(M\)} induces an injection on fundamental groups, at every basepoint.
Then there exists a prime decomposition of \(M\), together with a
geometric decomposition of each of its selected prime factors.

The incompressibility assumption concerns the actual reconstructed seam
\(T^2\to M\). Injection into just a hyperbolic model or a single cut
piece is not a substitute. The theorem does not assume \(M\) already
prime or irreducible, and does not assume that every raw internal graph
seam is incompressible. It permits zero ambient seams, no hyperbolic
components, and self-pairings between ports of one component.

The conclusion is a new existentially chosen prime decomposition and new
geometric cut data. It contains \textbf{no comparison identifying the
final seams with the input seams, no requirement to retain the supplied
hyperbolic metric, and no requirement to preserve the input graph labels
or external collars in those final data}. The new output nevertheless
reconstructs the same actual \(M\) by the maps required in T0. These
omissions make this an endpoint consequence of the fuller relative
blueprint, not an implementation of all the relative data in AT13--AT18.

In particular, this admitted theorem has \textbf{not} been proved in
Lean by combining T6 and T7. Those are statements for closed graph
carriers; an argument for graph carriers with boundary, their relative
prime/cap construction, and the hyperbolic reattachment is still inside
this single mixed admission.

\subsection{T10. The actual mixed endpoint assembly}\label{t10}

\textbf{Lean:}
{\small\texttt{GC.\allowbreak{}Graph\allowbreak{}Manifold.\allowbreak{}geometrizes\_\allowbreak{}of\_\allowbreak{}hyperbolic\allowbreak{}Or\allowbreak{}Graph}},
\href{https://github.com/qinz1yang/differential-geometry-dev/blob/ea0fae60ee01ef8c8d9b57a51794c6538f679c5d/DifferentialGeometry/Topology/ThreeManifold/GraphManifold/Refinement.lean\#L54}{Refinement:54}.

\textbf{Status: actual proof, depending on the admission T9.}

Under exactly the hypotheses of T9, \(M\) geometrizes in the existing
certificate sense. The proof takes the prime decomposition and
per-factor geometric decompositions supplied by T9 and packages them
into the certificate. No further topological input is used in this
packaging. It is a short conditional assembly, not a proof of relative
graph refinement.

The two graph routes in these files are therefore:

\[
\begin{aligned}
  &\text{raw closed graph presentation}\\
  &\quad\xrightarrow{\text{T6, admitted}}\text{graph prime factors}\\
  &\quad\xrightarrow{\text{T7, admitted}}\text{factor geometries}\\
  &\quad\xrightarrow{\text{T8, proved assembly}}\text{certificate}.
\end{aligned}
\]

and separately

\[
\begin{aligned}
  &\text{actual mixed torus decomposition}\\
  &\qquad{}+\text{ ambient injection}\\
  &\quad\xrightarrow{\text{T9, admitted}}\text{prime factors + geometries}\\
  &\quad\xrightarrow{\text{T10, proved assembly}}\text{certificate}.
\end{aligned}
\]

\subsection{FR1. Finite-order compactness: the exact admitted local
statement}\label{fr1}

\textbf{Lean:}
{\small\texttt{Differential\allowbreak{}Geometry.\allowbreak{}Cheeger\allowbreak{}Gromov\allowbreak{}Compactness.\allowbreak{}exists\_\allowbreak{}bilinear\_\allowbreak{}form\_\allowbreak{}limit\_\allowbreak{}subsequence\_\allowbreak{}of\_\allowbreak{}bounded\_\allowbreak{}derivatives}},
\href{https://github.com/qinz1yang/differential-geometry-dev/blob/ea0fae60ee01ef8c8d9b57a51794c6538f679c5d/DifferentialGeometry/Analysis/Calculus/Compactness/FiniteOrder.lean\#L11}{Compactness:11}.

\textbf{Status: admitted with {\small\texttt{sorry}}.}

Let \(E\) be a finite-dimensional complete real normed vector space,
\(U\subset E\) an open set, and \(K\geq1\) an integer. Let \(g_i(x)\) be
continuous bilinear forms on \(E\), with \(g_i\) defined as a function
on all of \(E\) and of class \(C^K\) on \(U\). Symmetry and metric
inequalities below are required only on \(U\).

Assume:

\begin{enumerate}
\def\labelenumi{\arabic{enumi}.}
\tightlist
\item
  For \textbf{each} integer \(0\leq q\leq K\) and \textbf{each} compact
  set \(S\subset U\), there is a real number \(C_{q,S}\) such that, for
  all sufficiently large \(i\), \[
       \sup_{x\in S}\|D^q g_i(x)\|\leq C_{q,S}.
  \] Both \(C_{q,S}\) and the starting index may depend on \(q,S\).
  There is no single global derivative bound on \(U\), and finitely many
  early sequence terms need not satisfy a chosen bound.
\item
  There are two \textbf{fixed} real constants \(a,b\), with \(a>0\),
  such that for \textbf{every} index \(i\), every \(x\in U\), and every
  \(v,w\in E\), \[
    g_i(x)(v,w)=g_i(x)(w,v),\qquad
    a\|v\|^2\leq g_i(x)(v,v)\leq b\|v\|^2.
  \] The same \(a,b\) apply on all of \(U\) and to all terms, including
  the early ones. The code does not separately assume \(b>0\); if
  \(U\ne\varnothing\) and \(\dim E>0\), the displayed inequalities force
  \(b\geq a>0\).
\end{enumerate}

Then there exist a \textbf{single strictly increasing} subsequence map
\(\phi:\mathbb N\to\mathbb N\) and a bilinear-form-valued function
\(g_\infty\), defined on all of \(E\), such that:

\begin{enumerate}
\def\labelenumi{\arabic{enumi}.}
\tightlist
\item
  \(g_\infty\) is \(C^{K-1}\) on \(U\).
\item
  On \textbf{every} compact \(S\subset U\), \(g_{\phi(i)}\to g_\infty\)
  in \(C^{K-1}\).
\item
  On \(U\), \(g_\infty\) is symmetric and satisfies the same two-sided
  bounds with exactly the original \(a,b\).
\end{enumerate}

For item 2, the code's precise convergence condition is

\[
 \forall\varepsilon>0\;\exists i_0\;\forall i\geq i_0\;
 \forall q\leq K-1\;\forall x\in S,
 \quad \|D^q(g_{\phi(i)}-g_\infty)(x)\|\leq\varepsilon.
\]

This unfolds
{\small\texttt{Map\allowbreak{}CPConvergence\allowbreak{}On}} in
\href{https://github.com/qinz1yang/differential-geometry-dev/blob/ea0fae60ee01ef8c8d9b57a51794c6538f679c5d/DifferentialGeometry/Analysis/Calculus/MapConvergence/Basic.lean\#L25}{MapConvergence/Basic.lean:25}.
Its derivatives are ordinary ambient Fréchet derivatives; openness of
\(U\) ensures that at its points these depend only on the local
restrictions to \(U\). Values outside \(U\) are not constrained by the
conclusion.

There is no assumption that \(U\) is bounded, connected, relatively
compact, convex, or nonempty. At \(K=1\), the conclusion is only a
continuous positive-definite limit and uniform convergence on compact
subsets. At \(U=\varnothing\), the regularity, convergence and
positivity-on-\(U\) assertions are vacuous; a subsequence and arbitrary
globally defined bilinear forms still exist. Dimension zero is allowed.
These cases do not give a counterexample to the statement.

This is a \textbf{local coefficient compactness statement on a fixed
Euclidean domain}. It has no sequence of manifolds, basepoints,
injectivity-radius hypothesis, volume lower bound, complete limit
metric, global atlas, overlap transitions, pointed convergence,
comparison embeddings, or sectional-curvature conclusion. Those are
substantial additional obligations in the full LFR14/LC81 blueprint.
Uniform coordinate derivative bounds are premises here; this theorem
does not derive them from curvature derivative bounds.

\subsection{FR2. Pulling back a finite-regularity metric: the exact
admitted statement}\label{fr2}

\textbf{Lean:}
{\small\texttt{Differential\allowbreak{}Geometry.\allowbreak{}Geometry.\allowbreak{}exists\_\allowbreak{}pullback\_\allowbreak{}metric\_\allowbreak{}of\_\allowbreak{}finite\_\allowbreak{}diffeomorph}},
\href{https://github.com/qinz1yang/differential-geometry-dev/blob/ea0fae60ee01ef8c8d9b57a51794c6538f679c5d/DifferentialGeometry/Geometry/Metric/Pullback/FiniteRegularity.lean\#L16}{Pullback:16}.

\textbf{Status: admitted with {\small\texttt{sorry}}.}

Let \(M,N\) already be smooth manifolds, with general real
finite-dimensional models with corners \(I,J\), respectively. The
statement does not assume compactness, connectedness, nonemptiness,
completeness, boundarylessness, Hausdorffness, or second countability.
It uses the manifold structures and finite-dimensional normed tangent
models needed for the derivative and metric definitions. The source and
target model vector spaces need not be declared to have equal dimension
in advance; existence of the diffeomorphism is the relevant premise.

Let \(K,s,r\in\mathbb N\) satisfy

\[
                 r\leq K,\qquad r+1\leq s.
\]

Let \(g\) be a \(C^K\) Riemannian metric on \(N\), and let \(f:M\to N\)
be a \(C^s\) diffeomorphism, meaning an actual bijection whose forward
\textbf{and inverse} maps are \(C^s\). Then there exists a \(C^r\)
Riemannian metric \(h\) on \(M\) satisfying the \textbf{exact tensor
identity}

\[
          h_x(v,w)=g_{f(x)}(df_xv,df_xw)
          \quad\text{for every }x\in M,\;v,w\in T_xM.
\]

Here a \(C^r\) Riemannian metric means a continuously varying symmetric
positive-definite bilinear form of the asserted regularity on the
tangent bundle, compatible with its existing fiber topology. It is not
an arbitrary semidefinite tensor. The code's
{\small\texttt{Diffeomorph}} definition explicitly contains \(C^s\)
regularity of the inverse; since \(s\geq1\), the derivative has the
isomorphism property needed for positivity.

The order bound is the familiar one-derivative cost in pulling back a
metric:

\[
                h\in C^{\min(K,s-1)},
\]

or any weaker finite order \(r\) allowed by the inequalities. The code
states existence at each such specified \(r\), rather than a
maximum-order formula. Natural-number subtraction and the lower bound
\(s\geq1\) therefore cause no exceptional \(s=0\) case.

The conclusion contains no explicit length, distance, completeness,
volume, or curvature-preservation theorem. It contains the exact tensor
formula from which appropriate preservation results can subsequently be
proved. It asserts neither smoothing of \(g\) nor a curvature-preserving
perturbation. It also does \textbf{not} construct a compatible smooth
structure on a manifold initially given only a finite-regularity atlas:
both carriers are already smooth at the start of this theorem.

\subsection{FR3. The two proved regularity-budget
corollaries}\label{fr3}

\textbf{Lean:}
{\small\texttt{exists\_\allowbreak{}pullback\_\allowbreak{}metric\_\allowbreak{}of\_\allowbreak{}diffeomorph\_\allowbreak{}one\_\allowbreak{}order\_\allowbreak{}higher}},
\href{https://github.com/qinz1yang/differential-geometry-dev/blob/ea0fae60ee01ef8c8d9b57a51794c6538f679c5d/DifferentialGeometry/Geometry/Metric/Pullback/FiniteRegularity.lean\#L26}{Pullback:26}.

\textbf{Status: actual proof by specialization of the admission FR2.}
For any natural \(K\), a \(C^{K+1}\) diffeomorphism pulls a \(C^K\)
metric back to a \(C^K\) metric with the exact displayed tensor
identity. This includes \(K=0\): a \(C^1\) diffeomorphism pulls a
continuous metric back to a continuous metric. The proof substitutes
\(s=K+1\) and \(r=K\) into FR2.

\textbf{Lean:}
{\small\texttt{exists\_\allowbreak{}pullback\_\allowbreak{}metric\_\allowbreak{}of\_\allowbreak{}diffeomorph\_\allowbreak{}three\_\allowbreak{}orders\_\allowbreak{}lower}},
\href{https://github.com/qinz1yang/differential-geometry-dev/blob/ea0fae60ee01ef8c8d9b57a51794c6538f679c5d/DifferentialGeometry/Geometry/Metric/Pullback/FiniteRegularity.lean\#L35}{Pullback:35}.

\textbf{Status: actual proof by specialization of the admission FR2.}
For an integer \(m\geq4\), a \(C^{m-3}\) diffeomorphism pulls a \(C^m\)
metric back to a \(C^{m-4}\) metric with exactly the same tensor
formula. The hypothesis \(m\geq4\) controls truncated natural-number
subtraction. At \(m=4\) the map is \(C^1\) and the output metric is
\(C^0\). The proof supplies \(r=m-4\leq m\) and \(r+1\leq m-3\) to FR2.

For the blueprint's common substitution \(m=K-1\), this second statement
gives order \(K-5\), provided \(K\geq5\). It does not claim that every
model map in the blueprint is automatically \(C^{m-3}\); a consumer must
supply such a map. It also gives no quantitative norm bound on the
pullback or its difference from another metric.

\subsection{R1. What is and is not yet connected in
Lean}\label{topology-r1}

The statement network in this part has two actual topology assemblies
(T8 and T10), and two actual regularity specializations (FR3). The mixed
topology producer T9 is one large admission rather than a proof
assembled from the relative constructions of the blueprint.

The local compactness theorem FR1 and finite pullback theorem FR2 are
separate prepared statements. At this commit, searches of the Lean
sources show no mathematical application of FR1 beyond its own
declaration, and no mathematical application of FR2 beyond its two FR3
corollaries. In particular, these files do not yet demonstrate that
their outputs satisfy the hypotheses of the late-flow or static-collapse
producers. Importing a file or listing it in the declaration audit does
not establish such an application.

Similarly, the two closed graph admissions T6 and T7 do feed T8, but
they are not called in the proof body of the mixed admission T9. A
future relative-refinement proof needs additional intermediate
statements to expose that connection.

\subsection{R2. Approval points and weaknesses identified by this
rereading}\label{topology-r2}

I found \textbf{no specific counterexample or internal contradiction} in
these five admitted statements during this rereading. That is a bounded
review judgment, not a mathematical proof or a claim that all stronger
blueprint arguments have been checked. The most consequential decisions
for Bennett Chow and Peng Lu are these:

\begin{enumerate}
\def\labelenumi{\arabic{enumi}.}
\tightlist
\item
  \textbf{Approve the regular-bundle meaning of raw graph presentation.}
  Its fibers are all ordinary circles; exceptional Seifert fibers must
  be removed or represented by further regular blocks before supplying
  it. The bases may be nonorientable. There is no global circle action
  in the definition.
\item
  \textbf{Approve the endpoint consequence rather than a full relative
  export.} T6, T7 and especially T9 omit the sphere/cap ledger and
  comparisons preserving the given decorations. This suffices for their
  chosen final certificate conclusion. It does not yet provide the finer
  relative interfaces needed to implement the blueprint proof in
  parallel without further design.
\item
  \textbf{Check the existential freedom in T9 deliberately.} The final
  chosen metrics and cuts need not retain the input hyperbolic metrics
  or seams. This is a weaker conclusion, not an inconsistency; any task
  expecting to reuse those exact markings cannot treat T9 as already
  exporting them.
\item
  \textbf{Keep nonempty components separate from possibly empty seam
  families.} The underlying {\small\texttt{Components}} record excludes
  an empty carrier. Empty late-flow slices must be handled elsewhere.
  No-cut geometric manifolds remain legal because the torus index may be
  empty.
\item
  \textbf{Do not call FR1 the global compactness theorem.} It is only
  the fixed-chart coefficient subsequence lemma. Producing the chart
  bounds, a global limit and actual comparison embeddings remains
  unrepresented here.
\item
  \textbf{Do not call FR2 a smooth-metric category bridge.} Its tensor
  remains only finitely regular. Choosing the initial compatible smooth
  carrier and deriving metric/curvature invariance are additional steps.
  FR2's actual content is the precise finite pullback regularity
  statement.
\end{enumerate}

The existing complete-interior metric convention is also essential: only
hyperbolic pieces require finite volume, and extra essential product
cuts are permitted. Requiring all geometries to have finite volume or
demanding an uncut model for every torus bundle would change these
statements substantially.

\subsection{R3. Evidence, source comparison, and review
limits}\label{topology-r3}

This review read all five actual Lean files, then expanded their
dependencies for compact carriers and components, manifold-with-boundary
interiors, finite torus gluing relations, smooth oriented
reconstruction, prime decomposition, endpoint geometry, metric
regularity, diffeomorphism regularity, and the exact \(C^p\)-convergence
predicate. Source locations and authored-declaration coverage are
recorded in
{\small\texttt{topology\_\allowbreak{}finite\_\allowbreak{}regularity\_\allowbreak{}coverage.\allowbreak{}json}}
next to this report.

The corresponding natural-language blueprint bodies were reread at:

\begin{itemize}
\tightlist
\item
  {\small\texttt{master207A.\allowbreak{}tex}}, lines 8526--8588:
  G01--G05, especially ordinary circle bundles, nonorientable bases, raw
  versus essential cuts, and actual ambient injection.
\item
  {\small\texttt{master207A.\allowbreak{}tex}}, lines 9467--9579:
  RG05--RG06, actual relative graph-prime/good-block output and its
  terminal block list.
\item
  {\small\texttt{master207A.\allowbreak{}tex}}, lines 10378--10431:
  GM06, full marked graph refinement and complete E1 witnesses.
\item
  {\small\texttt{master207B.\allowbreak{}tex}}, lines 11895--12036:
  AT13--AT18, relative prime/cap data, good-block refinement, actual
  capped ambient injection and prime reconstruction.
\item
  {\small\texttt{master207A.\allowbreak{}tex}}, lines 24826--24860:
  LFR01, compatible smooth structures and the finite derivative ledger.
\item
  {\small\texttt{master207A.\allowbreak{}tex}}, lines 25872--25908 and
  25970--26030: LFR14, finite-model construction and actual comparison
  maps.
\item
  {\small\texttt{master207A.\allowbreak{}tex}}, lines 26128--26144: the
  \(m=K-1\) consumer regularity distinction.
\end{itemize}

Those comparisons identify which conclusions are intentionally omitted;
they do not silently import them into the Lean signatures. No new
theorem contract or source classification was changed. Consequently this
report does not claim a new primary-source literature or erratum audit;
the previously recorded source checks for these unchanged mathematical
routes retain their historical scope.

Proof status was cross-checked against the recorded elaborated
declaration audit at this commit: the five admitted theorem names and
four admission-dependent corollary/assembly theorem names have
{\small\texttt{sorry\allowbreak{}Ax}}; the three collar proofs do not.
This report did not rebuild Lean or reprove any admission. The
previously recorded build checks elaboration and dependencies; it does
not establish the mathematical truth of a {\small\texttt{sorry}}
statement.

\section{Suggested joint review and the limits of this
report}\label{joint-review}

\subsection{A practical order for the mathematical
review}\label{review-order}

Read E1--E2 first and decide whether that is the desired theorem. Then
read F3--F8 together: these show the exact hypotheses that supply the
endpoint, the order in which choices are made, and what the written
assembly actually uses. C2, C4--C10, HA-2--HA-4, and T9 contain the
principal remaining mathematical assertions. Finally, read the prepared
branches, including Mostow--Prasad, the two closed graph refinements,
and FR1--FR3. They are useful future work, but are not already proofs
inside the main composite admissions.

The following questions can serve as a review worksheet. A request to
change one of these statements is a specification correction, not a
request to begin filling its proof immediately.

\begin{longtable}[]{@{}
  >{\raggedright\arraybackslash}p{(\linewidth - 2\tabcolsep) * \real{0.5000}}
  >{\raggedright\arraybackslash}p{(\linewidth - 2\tabcolsep) * \real{0.5000}}@{}}
\toprule\noalign{}
\begin{minipage}[b]{\linewidth}\raggedright
Review question
\end{minipage} & \begin{minipage}[b]{\linewidth}\raggedright
Where to look
\end{minipage} \\
\midrule\noalign{}
\endhead
\bottomrule\noalign{}
\endlastfoot
Is the smooth, connected, closed, oriented endpoint, with a noncanonical
prime list and collared fundamental-group-injective cuts, exactly the
desired public statement? & E1, T0, T5 \\
Are the eight fixed model metrics and the finite-volume requirement only
for hyperbolic interiors correct? & E2, HA-1.3--HA-1.4 \\
Do the static hypotheses suffice with intrinsic balls, whole-ball
derivative estimates, uncapped curvature radius, and arbitrary positive
derivative function? & C2, C6, C8--C10 \\
Is the boundary statement correct with diameter measured in the carrier,
no pairwise collar-disjointness field, and only boundary-image/index
retention? & C4--C5, C8--C9 \\
Is the selected-flow assertion true with precisely its
sequence-dependent derivative function and its
fixed-flow-before-sequence quantifiers? & F4, F6--F8 \\
Is the stated disk class the one for which the intended existence and
comparison arguments work, including smooth immersion at the boundary
and exact parametrized boundary equality? & HA-3.1--HA-3.4 \\
Can the actual physical area have the stipulated continuity and strict
upper barriers at every surgery time on the same future half-line? & F5,
HA-2, HA-5 \\
Are the weaker existential outputs of mixed refinement sufficient, even
though they do not retain prescribed incoming seams, metrics or
markings? & T6--T10 \\
Are the finite-order statements correct at their lowest orders, and
understood as local coefficient compactness and finite pullback rather
than the full category bridge? & FR1--FR3 \\
Which of the large admitted producers must be subdivided before the team
can safely assign independent proof tasks? & F7, C11, HA-4.4, T9 \\
\end{longtable}

\subsection{What the rereading found}\label{review-findings}

The direct and independent rereadings found no concrete counterexample
or explicit contradiction in the new admitted statements. This is a
bounded review judgment, not a proof of those statements or
certification that all source-to-interface comparisons are complete. The
mathematical obligations are exactly the assertions marked admitted,
together with any inherited foundation obligations outside the scope of
this review.

The most substantial finding concerns \textbf{how much mathematics
remains inside the largest admissions}. The selected-flow theorem
contains the production of the actual late geometric pieces and the
incompressibility obstruction. The mixed-refinement theorem contains the
relative topology needed to reach the endpoint. Their short consumers
compile, but this does not make those large producers small tasks.

There is also a distinction between \textbf{a faithful endpoint
consequence and a full translation of a blueprint node}. Several
statements deliberately export less data than the corresponding
blueprint development: relative marking comparisons, persistent
families, global parameter coherence beyond the common delta function,
and the global finite-regularity construction remain outside the
displayed interfaces. These omissions have been listed where they occur
rather than supplied implicitly by the prose.

The observation that the area-obstruction premise is equivalent to
injection once the scalar contradiction is granted is a logical analysis
in this report. It is not an additional equivalence theorem claimed to
exist in Lean. Likewise, explanatory proof sketches for admitted
elementary lemmas are identified as intended arguments, not as completed
proof bodies.

\subsection{Evidence and the meaning of complete
coverage}\label{report-evidence}

The declaration register below accounts for \textbf{all 121 authored
declarations in the 22 new mathematical files}, including the private
collar helper. Generated constructors and projections are covered by
their parent data definitions and their field descriptions; they are not
presented as hundreds of separate mathematical claims. The recorded
elaborated types and axiom closures were used to distinguish direct
admissions from actual proofs with or without admitted dependencies.

For this report, four family translations were read against the actual
Lean files, and the endpoint/flow translation received independent
cross-checks from the other reviewers. That process corrected a prose
overstatement: the generic area-obstruction predicate does not itself
tie its torus carrier to a flow component; the surrounding late-sequence
predicate makes that identification. The report also makes explicit that
the topological assembly discards the supplied induced-metric
equalities, although those equalities constrain the geometric producer's
input to that assembly.

Fresh mechanical checks for the report verify the frozen source hashes,
complete declaration coverage, the 17 direct admission names, the six
endpoint-reachable admissions, and the absence of Lean-source changes.
The PDF is rendered and inspected for legible equations and tables.
These are document and consistency checks. The earlier successful Lean
build is reported from its pinned verification receipt; it is not
described as a new build or as a proof of the admitted mathematics.

The companion files include the editable combined Markdown and
standalone LaTeX, the exact text of the 22 skeleton modules in
{\small\texttt{exact\_\allowbreak{}sources.\allowbreak{}txt}}, coverage
records, and a source-hash inventory. The source text is frozen at the
commit named on the cover. References to inherited definitions have
their own pinned source links and hashes; including a source in that
inventory does not assert that every theorem in it received a fresh
proof audit. Primary-source and blueprint passages actually reopened are
listed in the family sections. Older source and errata checks retain
their stated historical scope.

Your approval can therefore be recorded at two levels: whether the
displayed mathematical statements are correct and appropriate, and
whether a particular producer has been decomposed sufficiently for
assignment. The present report does not presume either approval and does
not modify the Lean skeleton.

\section{Complete declaration register}\label{declaration-register}

This register is exhaustive for the 121 authored declarations. Each
entry gives its status, a link to the mathematical explanation, and the
exact source line at the frozen commit. All file paths are under
{\small\texttt{Differential\allowbreak{}Geometry/\allowbreak{}}}. A
proof without a skeleton admission still uses the inherited library;
this label is not a new audit of that whole library.

\Needspace{6\baselineskip}

\subsection{M01. FiniteOrder.lean}\label{m01.-finiteorder.lean}

\textbf{File:}
{\small\texttt{Analysis/\allowbreak{}Calculus/\allowbreak{}Compactness/\allowbreak{}Finite\allowbreak{}Order.\allowbreak{}lean}}

\begin{itemize}
\tightlist
\item
  {\small\texttt{Differential\allowbreak{}Geometry.\allowbreak{}Cheeger\allowbreak{}Gromov\allowbreak{}Compactness.\allowbreak{}exists\_\allowbreak{}bilinear\_\allowbreak{}form\_\allowbreak{}limit\_\allowbreak{}subsequence\_\allowbreak{}of\_\allowbreak{}bounded\_\allowbreak{}derivatives}}\\
  Admitted. \hyperref[fr1]{Explanation: FR1};
  \href{https://github.com/qinz1yang/differential-geometry-dev/blob/ea0fae60ee01ef8c8d9b57a51794c6538f679c5d/DifferentialGeometry/Analysis/Calculus/Compactness/FiniteOrder.lean\#L11}{source
  line 11}.
\end{itemize}

\Needspace{6\baselineskip}

\subsection{M02.
AreaUpperBarrier.lean}\label{m02.-areaupperbarrier.lean}

\textbf{File:}
{\small\texttt{Analysis/\allowbreak{}ODE/\allowbreak{}Area\allowbreak{}Upper\allowbreak{}Barrier.\allowbreak{}lean}}

\begin{itemize}
\item
  {\small\texttt{Differential\allowbreak{}Geometry.\allowbreak{}Analysis.\allowbreak{}has\allowbreak{}Local\allowbreak{}Smooth\allowbreak{}Upper\allowbreak{}Barrier}}\\
  Definition/construction. \hyperref[ha-2-1]{Explanation: HA-2.1};
  \href{https://github.com/qinz1yang/differential-geometry-dev/blob/ea0fae60ee01ef8c8d9b57a51794c6538f679c5d/DifferentialGeometry/Analysis/ODE/AreaUpperBarrier.lean\#L13}{source
  line 13}.
\item
  {\small\texttt{Differential\allowbreak{}Geometry.\allowbreak{}Analysis.\allowbreak{}not\_\allowbreak{}nonnegative\_\allowbreak{}area\_\allowbreak{}upper\_\allowbreak{}barriers\_\allowbreak{}shift}}\\
  Admitted. \hyperref[ha-2-2]{Explanation: HA-2.2};
  \href{https://github.com/qinz1yang/differential-geometry-dev/blob/ea0fae60ee01ef8c8d9b57a51794c6538f679c5d/DifferentialGeometry/Analysis/ODE/AreaUpperBarrier.lean\#L18}{source
  line 18}.
\item
  {\small\texttt{Differential\allowbreak{}Geometry.\allowbreak{}Analysis.\allowbreak{}not\_\allowbreak{}nonnegative\_\allowbreak{}area\_\allowbreak{}upper\_\allowbreak{}barriers}}\\
  Written proof; admitted dependence. \hyperref[ha-2-3]{Explanation:
  HA-2.3};
  \href{https://github.com/qinz1yang/differential-geometry-dev/blob/ea0fae60ee01ef8c8d9b57a51794c6538f679c5d/DifferentialGeometry/Analysis/ODE/AreaUpperBarrier.lean\#L28}{source
  line 28}.
\item
  {\small\texttt{Differential\allowbreak{}Geometry.\allowbreak{}Analysis.\allowbreak{}not\_\allowbreak{}nonnegative\_\allowbreak{}area\_\allowbreak{}upper\_\allowbreak{}barriers\_\allowbreak{}pi\_\allowbreak{}shift}}\\
  Written proof; admitted dependence. \hyperref[ha-2-3]{Explanation:
  HA-2.3};
  \href{https://github.com/qinz1yang/differential-geometry-dev/blob/ea0fae60ee01ef8c8d9b57a51794c6538f679c5d/DifferentialGeometry/Analysis/ODE/AreaUpperBarrier.lean\#L39}{source
  line 39}.
\item
  {\small\texttt{Differential\allowbreak{}Geometry.\allowbreak{}Analysis.\allowbreak{}not\_\allowbreak{}nonnegative\_\allowbreak{}area\_\allowbreak{}upper\_\allowbreak{}barriers\_\allowbreak{}pi}}\\
  Written proof; admitted dependence. \hyperref[ha-2-3]{Explanation:
  HA-2.3};
  \href{https://github.com/qinz1yang/differential-geometry-dev/blob/ea0fae60ee01ef8c8d9b57a51794c6538f679c5d/DifferentialGeometry/Analysis/ODE/AreaUpperBarrier.lean\#L52}{source
  line 52}.
\end{itemize}

\Needspace{6\baselineskip}

\subsection{M03. CurvatureScale.lean}\label{m03.-curvaturescale.lean}

\textbf{File:}
{\small\texttt{Geometry/\allowbreak{}Collapse/\allowbreak{}Curvature\allowbreak{}Scale.\allowbreak{}lean}}

\begin{itemize}
\item
  {\small\texttt{Differential\allowbreak{}Geometry.\allowbreak{}Geometry.\allowbreak{}Collapse.\allowbreak{}curvature\allowbreak{}Radius}}\\
  Definition/construction. \hyperref[collapse-radius]{Explanation: C2};
  \href{https://github.com/qinz1yang/differential-geometry-dev/blob/ea0fae60ee01ef8c8d9b57a51794c6538f679c5d/DifferentialGeometry/Geometry/Collapse/CurvatureScale.lean\#L23}{source
  line 23}.
\item
  {\small\texttt{Differential\allowbreak{}Geometry.\allowbreak{}Geometry.\allowbreak{}Collapse.\allowbreak{}ball\allowbreak{}Volume}}\\
  Definition/construction. \hyperref[collapse-radius]{Explanation: C2};
  \href{https://github.com/qinz1yang/differential-geometry-dev/blob/ea0fae60ee01ef8c8d9b57a51794c6538f679c5d/DifferentialGeometry/Geometry/Collapse/CurvatureScale.lean\#L28}{source
  line 28}.
\item
  {\small\texttt{Differential\allowbreak{}Geometry.\allowbreak{}Geometry.\allowbreak{}Collapse.\allowbreak{}euclidean\allowbreak{}Unit\allowbreak{}Ball\allowbreak{}Volume}}\\
  Definition/construction. \hyperref[collapse-radius]{Explanation: C2};
  \href{https://github.com/qinz1yang/differential-geometry-dev/blob/ea0fae60ee01ef8c8d9b57a51794c6538f679c5d/DifferentialGeometry/Geometry/Collapse/CurvatureScale.lean\#L31}{source
  line 31}.
\item
  {\small\texttt{Differential\allowbreak{}Geometry.\allowbreak{}Geometry.\allowbreak{}Collapse.\allowbreak{}volume\allowbreak{}Collapsed\allowbreak{}At\allowbreak{}Curvature\allowbreak{}Scale}}\\
  Definition/construction. \hyperref[collapse-radius]{Explanation: C2};
  \href{https://github.com/qinz1yang/differential-geometry-dev/blob/ea0fae60ee01ef8c8d9b57a51794c6538f679c5d/DifferentialGeometry/Geometry/Collapse/CurvatureScale.lean\#L33}{source
  line 33}.
\item
  {\small\texttt{Differential\allowbreak{}Geometry.\allowbreak{}Geometry.\allowbreak{}Collapse.\allowbreak{}curvature\allowbreak{}Derivatives\allowbreak{}Controlled}}\\
  Definition/construction. \hyperref[collapse-radius]{Explanation: C2};
  \href{https://github.com/qinz1yang/differential-geometry-dev/blob/ea0fae60ee01ef8c8d9b57a51794c6538f679c5d/DifferentialGeometry/Geometry/Collapse/CurvatureScale.lean\#L38}{source
  line 38}.
\item
  {\small\texttt{Differential\allowbreak{}Geometry.\allowbreak{}Geometry.\allowbreak{}Collapse.\allowbreak{}curvature\allowbreak{}Radius\_\allowbreak{}pos}}\\
  Admitted. \hyperref[collapse-radius]{Explanation: C2};
  \href{https://github.com/qinz1yang/differential-geometry-dev/blob/ea0fae60ee01ef8c8d9b57a51794c6538f679c5d/DifferentialGeometry/Geometry/Collapse/CurvatureScale.lean\#L46}{source
  line 46}.
\item
  {\small\texttt{Differential\allowbreak{}Geometry.\allowbreak{}Geometry.\allowbreak{}Collapse.\allowbreak{}curvature\allowbreak{}Radius\_\allowbreak{}eq\_\allowbreak{}top\_\allowbreak{}iff}}\\
  Admitted. \hyperref[collapse-radius]{Explanation: C2};
  \href{https://github.com/qinz1yang/differential-geometry-dev/blob/ea0fae60ee01ef8c8d9b57a51794c6538f679c5d/DifferentialGeometry/Geometry/Collapse/CurvatureScale.lean\#L50}{source
  line 50}.
\item
  {\small\texttt{Differential\allowbreak{}Geometry.\allowbreak{}Geometry.\allowbreak{}Collapse.\allowbreak{}curvature\allowbreak{}Derivatives\allowbreak{}Controlled\_\allowbreak{}mono\_\allowbreak{}threshold}}\\
  Written proof; no skeleton admission.
  \hyperref[collapse-radius]{Explanation: C2};
  \href{https://github.com/qinz1yang/differential-geometry-dev/blob/ea0fae60ee01ef8c8d9b57a51794c6538f679c5d/DifferentialGeometry/Geometry/Collapse/CurvatureScale.lean\#L55}{source
  line 55}.
\end{itemize}

\Needspace{6\baselineskip}

\subsection{M04. CuspBoundary.lean}\label{m04.-cuspboundary.lean}

\textbf{File:}
{\small\texttt{Geometry/\allowbreak{}Collapse/\allowbreak{}Cusp\allowbreak{}Boundary.\allowbreak{}lean}}

\begin{itemize}
\item
  {\small\texttt{Differential\allowbreak{}Geometry.\allowbreak{}Geometry.\allowbreak{}Collapse.\allowbreak{}cusp\allowbreak{}Metric\allowbreak{}Error}}\\
  Definition/construction.
  \hyperref[collapse-cuspembedding]{Explanation: C4};
  \href{https://github.com/qinz1yang/differential-geometry-dev/blob/ea0fae60ee01ef8c8d9b57a51794c6538f679c5d/DifferentialGeometry/Geometry/Collapse/CuspBoundary.lean\#L15}{source
  line 15}.
\item
  {\small\texttt{Differential\allowbreak{}Geometry.\allowbreak{}Geometry.\allowbreak{}Collapse.\allowbreak{}cusp\allowbreak{}Metric\allowbreak{}Error\allowbreak{}Bound}}\\
  Definition/construction.
  \hyperref[collapse-cuspembedding]{Explanation: C4};
  \href{https://github.com/qinz1yang/differential-geometry-dev/blob/ea0fae60ee01ef8c8d9b57a51794c6538f679c5d/DifferentialGeometry/Geometry/Collapse/CuspBoundary.lean\#L22}{source
  line 22}.
\item
  {\small\texttt{Differential\allowbreak{}Geometry.\allowbreak{}Geometry.\allowbreak{}Collapse.\allowbreak{}Cusp\allowbreak{}Embedding}}\\
  Definition/construction.
  \hyperref[collapse-cuspembedding]{Explanation: C4};
  \href{https://github.com/qinz1yang/differential-geometry-dev/blob/ea0fae60ee01ef8c8d9b57a51794c6538f679c5d/DifferentialGeometry/Geometry/Collapse/CuspBoundary.lean\#L29}{source
  line 29}.
\item
  {\small\texttt{Differential\allowbreak{}Geometry.\allowbreak{}Geometry.\allowbreak{}Collapse.\allowbreak{}Nearly\allowbreak{}Cuspidal\allowbreak{}Boundary}}\\
  Definition/construction.
  \hyperref[collapse-cuspembedding]{Explanation: C4};
  \href{https://github.com/qinz1yang/differential-geometry-dev/blob/ea0fae60ee01ef8c8d9b57a51794c6538f679c5d/DifferentialGeometry/Geometry/Collapse/CuspBoundary.lean\#L43}{source
  line 43}.
\item
  {\small\texttt{Differential\allowbreak{}Geometry.\allowbreak{}Geometry.\allowbreak{}Collapse.\allowbreak{}distance\allowbreak{}To\allowbreak{}Boundary}}\\
  Definition/construction.
  \hyperref[collapse-boundarydistance]{Explanation: C5};
  \href{https://github.com/qinz1yang/differential-geometry-dev/blob/ea0fae60ee01ef8c8d9b57a51794c6538f679c5d/DifferentialGeometry/Geometry/Collapse/CuspBoundary.lean\#L56}{source
  line 56}.
\item
  {\small\texttt{Differential\allowbreak{}Geometry.\allowbreak{}Geometry.\allowbreak{}Collapse.\allowbreak{}boundary\allowbreak{}Volume\allowbreak{}Collapsed}}\\
  Definition/construction.
  \hyperref[collapse-boundarydistance]{Explanation: C5};
  \href{https://github.com/qinz1yang/differential-geometry-dev/blob/ea0fae60ee01ef8c8d9b57a51794c6538f679c5d/DifferentialGeometry/Geometry/Collapse/CuspBoundary.lean\#L60}{source
  line 60}.
\item
  {\small\texttt{Differential\allowbreak{}Geometry.\allowbreak{}Geometry.\allowbreak{}Collapse.\allowbreak{}Cusp\allowbreak{}Embedding.\allowbreak{}weaken}}\\
  Definition/construction.
  \hyperref[collapse-cuspembedding]{Explanation: C4};
  \href{https://github.com/qinz1yang/differential-geometry-dev/blob/ea0fae60ee01ef8c8d9b57a51794c6538f679c5d/DifferentialGeometry/Geometry/Collapse/CuspBoundary.lean\#L64}{source
  line 64}.
\item
  {\small\texttt{Differential\allowbreak{}Geometry.\allowbreak{}Geometry.\allowbreak{}Collapse.\allowbreak{}Nearly\allowbreak{}Cuspidal\allowbreak{}Boundary.\allowbreak{}weaken}}\\
  Definition/construction.
  \hyperref[collapse-cuspembedding]{Explanation: C4};
  \href{https://github.com/qinz1yang/differential-geometry-dev/blob/ea0fae60ee01ef8c8d9b57a51794c6538f679c5d/DifferentialGeometry/Geometry/Collapse/CuspBoundary.lean\#L70}{source
  line 70}.
\end{itemize}

\Needspace{6\baselineskip}

\subsection{M05. GraphManifold.lean}\label{m05.-graphmanifold.lean}

\textbf{File:}
{\small\texttt{Geometry/\allowbreak{}Collapse/\allowbreak{}Graph\allowbreak{}Manifold.\allowbreak{}lean}}

\begin{itemize}
\item
  {\small\texttt{Differential\allowbreak{}Geometry.\allowbreak{}Geometry.\allowbreak{}Collapse.\allowbreak{}closed\allowbreak{}Collapse\allowbreak{}Hypotheses}}\\
  Definition/construction. \hyperref[collapse-hypotheses]{Explanation:
  C6};
  \href{https://github.com/qinz1yang/differential-geometry-dev/blob/ea0fae60ee01ef8c8d9b57a51794c6538f679c5d/DifferentialGeometry/Geometry/Collapse/GraphManifold.lean\#L11}{source
  line 11}.
\item
  {\small\texttt{Differential\allowbreak{}Geometry.\allowbreak{}Geometry.\allowbreak{}Collapse.\allowbreak{}boundary\allowbreak{}Collapse\allowbreak{}Hypotheses}}\\
  Definition/construction. \hyperref[collapse-hypotheses]{Explanation:
  C6};
  \href{https://github.com/qinz1yang/differential-geometry-dev/blob/ea0fae60ee01ef8c8d9b57a51794c6538f679c5d/DifferentialGeometry/Geometry/Collapse/GraphManifold.lean\#L17}{source
  line 17}.
\item
  {\small\texttt{Differential\allowbreak{}Geometry.\allowbreak{}Geometry.\allowbreak{}Collapse.\allowbreak{}static\allowbreak{}Collapse\allowbreak{}Hypotheses}}\\
  Definition/construction. \hyperref[collapse-hypotheses]{Explanation:
  C6};
  \href{https://github.com/qinz1yang/differential-geometry-dev/blob/ea0fae60ee01ef8c8d9b57a51794c6538f679c5d/DifferentialGeometry/Geometry/Collapse/GraphManifold.lean\#L22}{source
  line 22}.
\item
  {\small\texttt{Differential\allowbreak{}Geometry.\allowbreak{}Geometry.\allowbreak{}Collapse.\allowbreak{}volume\allowbreak{}Collapsed\allowbreak{}At\allowbreak{}Curvature\allowbreak{}Scale\_\allowbreak{}mono}}\\
  Written proof; no skeleton admission.
  \hyperref[collapse-hypotheses]{Explanation: C6};
  \href{https://github.com/qinz1yang/differential-geometry-dev/blob/ea0fae60ee01ef8c8d9b57a51794c6538f679c5d/DifferentialGeometry/Geometry/Collapse/GraphManifold.lean\#L26}{source
  line 26}.
\item
  {\small\texttt{Differential\allowbreak{}Geometry.\allowbreak{}Geometry.\allowbreak{}Collapse.\allowbreak{}closed\allowbreak{}Collapse\allowbreak{}Hypotheses\_\allowbreak{}mono}}\\
  Written proof; no skeleton admission.
  \hyperref[collapse-hypotheses]{Explanation: C6};
  \href{https://github.com/qinz1yang/differential-geometry-dev/blob/ea0fae60ee01ef8c8d9b57a51794c6538f679c5d/DifferentialGeometry/Geometry/Collapse/GraphManifold.lean\#L35}{source
  line 35}.
\item
  {\small\texttt{Differential\allowbreak{}Geometry.\allowbreak{}Geometry.\allowbreak{}Collapse.\allowbreak{}boundary\allowbreak{}Collapse\allowbreak{}Hypotheses\_\allowbreak{}mono}}\\
  Written proof; no skeleton admission.
  \hyperref[collapse-hypotheses]{Explanation: C6};
  \href{https://github.com/qinz1yang/differential-geometry-dev/blob/ea0fae60ee01ef8c8d9b57a51794c6538f679c5d/DifferentialGeometry/Geometry/Collapse/GraphManifold.lean\#L43}{source
  line 43}.
\item
  {\small\texttt{Differential\allowbreak{}Geometry.\allowbreak{}Geometry.\allowbreak{}Collapse.\allowbreak{}exists\_\allowbreak{}raw\allowbreak{}Graph\allowbreak{}Presentation\_\allowbreak{}of\_\allowbreak{}nonnegative}}\\
  Admitted. \hyperref[collapse-admitted]{Explanation: C8};
  \href{https://github.com/qinz1yang/differential-geometry-dev/blob/ea0fae60ee01ef8c8d9b57a51794c6538f679c5d/DifferentialGeometry/Geometry/Collapse/GraphManifold.lean\#L52}{source
  line 52}.
\item
  {\small\texttt{Differential\allowbreak{}Geometry.\allowbreak{}Geometry.\allowbreak{}Collapse.\allowbreak{}exists\_\allowbreak{}closed\_\allowbreak{}graph\_\allowbreak{}threshold}}\\
  Admitted. \hyperref[collapse-admitted]{Explanation: C8};
  \href{https://github.com/qinz1yang/differential-geometry-dev/blob/ea0fae60ee01ef8c8d9b57a51794c6538f679c5d/DifferentialGeometry/Geometry/Collapse/GraphManifold.lean\#L60}{source
  line 60}.
\item
  {\small\texttt{Differential\allowbreak{}Geometry.\allowbreak{}Geometry.\allowbreak{}Collapse.\allowbreak{}exists\_\allowbreak{}boundary\_\allowbreak{}graph\_\allowbreak{}threshold}}\\
  Admitted. \hyperref[collapse-admitted]{Explanation: C8};
  \href{https://github.com/qinz1yang/differential-geometry-dev/blob/ea0fae60ee01ef8c8d9b57a51794c6538f679c5d/DifferentialGeometry/Geometry/Collapse/GraphManifold.lean\#L68}{source
  line 68}.
\item
  {\small\texttt{Differential\allowbreak{}Geometry.\allowbreak{}Geometry.\allowbreak{}Collapse.\allowbreak{}exists\_\allowbreak{}graph\_\allowbreak{}threshold}}\\
  Written proof; admitted dependence.
  \hyperref[collapse-composition]{Explanation: C9};
  \href{https://github.com/qinz1yang/differential-geometry-dev/blob/ea0fae60ee01ef8c8d9b57a51794c6538f679c5d/DifferentialGeometry/Geometry/Collapse/GraphManifold.lean\#L80}{source
  line 80}.
\item
  {\small\texttt{Differential\allowbreak{}Geometry.\allowbreak{}Geometry.\allowbreak{}Collapse.\allowbreak{}exists\_\allowbreak{}componentwise\_\allowbreak{}graph\_\allowbreak{}threshold}}\\
  Written proof; admitted dependence.
  \hyperref[collapse-composition]{Explanation: C9};
  \href{https://github.com/qinz1yang/differential-geometry-dev/blob/ea0fae60ee01ef8c8d9b57a51794c6538f679c5d/DifferentialGeometry/Geometry/Collapse/GraphManifold.lean\#L97}{source
  line 97}.
\end{itemize}

\Needspace{6\baselineskip}

\subsection{M06.
TorusDecomposition.lean}\label{m06.-torusdecomposition.lean}

\textbf{File:}
{\small\texttt{Geometry/\allowbreak{}Collapse/\allowbreak{}Torus\allowbreak{}Decomposition.\allowbreak{}lean}}

\begin{itemize}
\item
  {\small\texttt{Differential\allowbreak{}Geometry.\allowbreak{}Geometry.\allowbreak{}Collapse.\allowbreak{}cut\allowbreak{}Piece\allowbreak{}Map}}\\
  Definition/construction. \hyperref[f3-induced-cut-metric]{Explanation:
  F3};
  \href{https://github.com/qinz1yang/differential-geometry-dev/blob/ea0fae60ee01ef8c8d9b57a51794c6538f679c5d/DifferentialGeometry/Geometry/Collapse/TorusDecomposition.lean\#L11}{source
  line 11}.
\item
  {\small\texttt{Differential\allowbreak{}Geometry.\allowbreak{}Geometry.\allowbreak{}Collapse.\allowbreak{}is\allowbreak{}Induced\allowbreak{}Cut\allowbreak{}Metric}}\\
  Definition/construction. \hyperref[f3-induced-cut-metric]{Explanation:
  F3};
  \href{https://github.com/qinz1yang/differential-geometry-dev/blob/ea0fae60ee01ef8c8d9b57a51794c6538f679c5d/DifferentialGeometry/Geometry/Collapse/TorusDecomposition.lean\#L16}{source
  line 16}.
\item
  {\small\texttt{Differential\allowbreak{}Geometry.\allowbreak{}Geometry.\allowbreak{}Collapse.\allowbreak{}Hyperbolic\allowbreak{}Or\allowbreak{}Collapsed}}\\
  Definition/construction. \hyperref[f3-induced-cut-metric]{Explanation:
  F3};
  \href{https://github.com/qinz1yang/differential-geometry-dev/blob/ea0fae60ee01ef8c8d9b57a51794c6538f679c5d/DifferentialGeometry/Geometry/Collapse/TorusDecomposition.lean\#L25}{source
  line 25}.
\item
  {\small\texttt{Differential\allowbreak{}Geometry.\allowbreak{}Geometry.\allowbreak{}Collapse.\allowbreak{}geometrizes\_\allowbreak{}of\_\allowbreak{}hyperbolic\allowbreak{}Or\allowbreak{}Collapsed}}\\
  Written proof; admitted dependence.
  \hyperref[f3-induced-cut-metric]{Explanation: F3};
  \href{https://github.com/qinz1yang/differential-geometry-dev/blob/ea0fae60ee01ef8c8d9b57a51794c6538f679c5d/DifferentialGeometry/Geometry/Collapse/TorusDecomposition.lean\#L41}{source
  line 41}.
\end{itemize}

\Needspace{6\baselineskip}

\subsection{M07.
UniformDerivativeBounds.lean}\label{m07.-uniformderivativebounds.lean}

\textbf{File:}
{\small\texttt{Geometry/\allowbreak{}Collapse/\allowbreak{}Uniform\allowbreak{}Derivative\allowbreak{}Bounds.\allowbreak{}lean}}

\begin{itemize}
\tightlist
\item
  {\small\texttt{Differential\allowbreak{}Geometry.\allowbreak{}Geometry.\allowbreak{}Collapse.\allowbreak{}exists\_\allowbreak{}common\_\allowbreak{}curvature\_\allowbreak{}derivative\_\allowbreak{}bound}}\\
  Admitted. \hyperref[collapse-uniform]{Explanation: C10};
  \href{https://github.com/qinz1yang/differential-geometry-dev/blob/ea0fae60ee01ef8c8d9b57a51794c6538f679c5d/DifferentialGeometry/Geometry/Collapse/UniformDerivativeBounds.lean\#L10}{source
  line 10}.
\end{itemize}

\Needspace{6\baselineskip}

\subsection{M08. Metric.lean}\label{m08.-metric.lean}

\textbf{File:}
{\small\texttt{Geometry/\allowbreak{}Connection/\allowbreak{}Tensor\allowbreak{}Nabla/\allowbreak{}Iterated/\allowbreak{}Metric.\allowbreak{}lean}}

\begin{itemize}
\item
  {\small\texttt{Differential\allowbreak{}Geometry.\allowbreak{}Geometry.\allowbreak{}Connection.\allowbreak{}metric\allowbreak{}Covariant\allowbreak{}Derivative}}\\
  Definition/construction. \hyperref[collapse-tensor]{Explanation: C1};
  \href{https://github.com/qinz1yang/differential-geometry-dev/blob/ea0fae60ee01ef8c8d9b57a51794c6538f679c5d/DifferentialGeometry/Geometry/Connection/TensorNabla/Iterated/Metric.lean\#L17}{source
  line 17}.
\item
  {\small\texttt{Differential\allowbreak{}Geometry.\allowbreak{}Geometry.\allowbreak{}Connection.\allowbreak{}iterated\allowbreak{}Metric\allowbreak{}Covariant\allowbreak{}Derivative}}\\
  Definition/construction. \hyperref[collapse-tensor]{Explanation: C1};
  \href{https://github.com/qinz1yang/differential-geometry-dev/blob/ea0fae60ee01ef8c8d9b57a51794c6538f679c5d/DifferentialGeometry/Geometry/Connection/TensorNabla/Iterated/Metric.lean\#L30}{source
  line 30}.
\end{itemize}

\Needspace{6\baselineskip}

\subsection{M09. DerivativeNorm.lean}\label{m09.-derivativenorm.lean}

\textbf{File:}
{\small\texttt{Geometry/\allowbreak{}Curvature/\allowbreak{}Metric/\allowbreak{}Derivative\allowbreak{}Norm.\allowbreak{}lean}}

\begin{itemize}
\tightlist
\item
  {\small\texttt{Differential\allowbreak{}Geometry.\allowbreak{}Geometry.\allowbreak{}Curvature.\allowbreak{}curvature\allowbreak{}Derivative\allowbreak{}Norm}}\\
  Definition/construction. \hyperref[collapse-tensor]{Explanation: C1};
  \href{https://github.com/qinz1yang/differential-geometry-dev/blob/ea0fae60ee01ef8c8d9b57a51794c6538f679c5d/DifferentialGeometry/Geometry/Curvature/Metric/DerivativeNorm.lean\#L18}{source
  line 18}.
\end{itemize}

\Needspace{6\baselineskip}

\subsection{M10.
CuspIncompressibility.lean}\label{m10.-cuspincompressibility.lean}

\textbf{File:}
{\small\texttt{Geometry/\allowbreak{}Flow/\allowbreak{}Ricci\allowbreak{}Flow/\allowbreak{}Long\allowbreak{}Time/\allowbreak{}Cusp\allowbreak{}Incompressibility.\allowbreak{}lean}}

\begin{itemize}
\item
  {\small\texttt{Differential\allowbreak{}Geometry.\allowbreak{}Geometry.\allowbreak{}Ricci\allowbreak{}Flow.\allowbreak{}cusp\_\allowbreak{}injective\_\allowbreak{}of\_\allowbreak{}exterior\_\allowbreak{}area\_\allowbreak{}barriers}}\\
  Written proof; admitted dependence. \hyperref[ha-4-1]{Explanation:
  HA-4.1};
  \href{https://github.com/qinz1yang/differential-geometry-dev/blob/ea0fae60ee01ef8c8d9b57a51794c6538f679c5d/DifferentialGeometry/Geometry/Flow/RicciFlow/LongTime/CuspIncompressibility.lean\#L15}{source
  line 15}.
\item
  {\small\texttt{Differential\allowbreak{}Geometry.\allowbreak{}Geometry.\allowbreak{}Ricci\allowbreak{}Flow.\allowbreak{}cusp\_\allowbreak{}slice\_\allowbreak{}injective\_\allowbreak{}iff}}\\
  Written proof; no skeleton admission. \hyperref[ha-4-2]{Explanation:
  HA-4.2};
  \href{https://github.com/qinz1yang/differential-geometry-dev/blob/ea0fae60ee01ef8c8d9b57a51794c6538f679c5d/DifferentialGeometry/Geometry/Flow/RicciFlow/LongTime/CuspIncompressibility.lean\#L43}{source
  line 43}.
\item
  {\small\texttt{Differential\allowbreak{}Geometry.\allowbreak{}Geometry.\allowbreak{}Ricci\allowbreak{}Flow.\allowbreak{}incompressible\_\allowbreak{}of\_\allowbreak{}area\_\allowbreak{}barriers}}\\
  Written proof; admitted dependence. \hyperref[ha-4-3]{Explanation:
  HA-4.3};
  \href{https://github.com/qinz1yang/differential-geometry-dev/blob/ea0fae60ee01ef8c8d9b57a51794c6538f679c5d/DifferentialGeometry/Geometry/Flow/RicciFlow/LongTime/CuspIncompressibility.lean\#L51}{source
  line 51}.
\item
  {\small\texttt{Differential\allowbreak{}Geometry.\allowbreak{}Geometry.\allowbreak{}Ricci\allowbreak{}Flow.\allowbreak{}incompressible\_\allowbreak{}of\_\allowbreak{}shifted\_\allowbreak{}area\_\allowbreak{}barriers}}\\
  Written proof; admitted dependence. \hyperref[ha-4-3]{Explanation:
  HA-4.3};
  \href{https://github.com/qinz1yang/differential-geometry-dev/blob/ea0fae60ee01ef8c8d9b57a51794c6538f679c5d/DifferentialGeometry/Geometry/Flow/RicciFlow/LongTime/CuspIncompressibility.lean\#L67}{source
  line 67}.
\end{itemize}

\Needspace{6\baselineskip}

\subsection{M11.
ExteriorDiskFlow.lean}\label{m11.-exteriordiskflow.lean}

\textbf{File:}
{\small\texttt{Geometry/\allowbreak{}Flow/\allowbreak{}Ricci\allowbreak{}Flow/\allowbreak{}Long\allowbreak{}Time/\allowbreak{}Exterior\allowbreak{}Disk\allowbreak{}Flow.\allowbreak{}lean}}

\begin{itemize}
\item
  {\small\texttt{GC.\allowbreak{}Long\allowbreak{}Time.\allowbreak{}post\allowbreak{}Stage}}\\
  Definition/construction. \hyperref[ha-5-1]{Explanation: HA-5.1};
  \href{https://github.com/qinz1yang/differential-geometry-dev/blob/ea0fae60ee01ef8c8d9b57a51794c6538f679c5d/DifferentialGeometry/Geometry/Flow/RicciFlow/LongTime/ExteriorDiskFlow.lean\#L18}{source
  line 18}.
\item
  {\small\texttt{GC.\allowbreak{}Long\allowbreak{}Time.\allowbreak{}post\allowbreak{}Metric}}\\
  Definition/construction. \hyperref[ha-5-2]{Explanation: HA-5.2};
  \href{https://github.com/qinz1yang/differential-geometry-dev/blob/ea0fae60ee01ef8c8d9b57a51794c6538f679c5d/DifferentialGeometry/Geometry/Flow/RicciFlow/LongTime/ExteriorDiskFlow.lean\#L22}{source
  line 22}.
\item
  {\small\texttt{GC.\allowbreak{}Long\allowbreak{}Time.\allowbreak{}post\allowbreak{}Stage\_\allowbreak{}regular\allowbreak{}Slice}}\\
  Written proof; no skeleton admission. \hyperref[ha-5-3]{Explanation:
  HA-5.3};
  \href{https://github.com/qinz1yang/differential-geometry-dev/blob/ea0fae60ee01ef8c8d9b57a51794c6538f679c5d/DifferentialGeometry/Geometry/Flow/RicciFlow/LongTime/ExteriorDiskFlow.lean\#L26}{source
  line 26}.
\item
  {\small\texttt{GC.\allowbreak{}Long\allowbreak{}Time.\allowbreak{}post\allowbreak{}Metric\_\allowbreak{}regular\allowbreak{}Slice}}\\
  Written proof; no skeleton admission. \hyperref[ha-5-3]{Explanation:
  HA-5.3};
  \href{https://github.com/qinz1yang/differential-geometry-dev/blob/ea0fae60ee01ef8c8d9b57a51794c6538f679c5d/DifferentialGeometry/Geometry/Flow/RicciFlow/LongTime/ExteriorDiskFlow.lean\#L37}{source
  line 37}.
\item
  {\small\texttt{GC.\allowbreak{}Long\allowbreak{}Time.\allowbreak{}exterior\allowbreak{}Disk\allowbreak{}Area}}\\
  Definition/construction. \hyperref[ha-5-4]{Explanation: HA-5.4};
  \href{https://github.com/qinz1yang/differential-geometry-dev/blob/ea0fae60ee01ef8c8d9b57a51794c6538f679c5d/DifferentialGeometry/Geometry/Flow/RicciFlow/LongTime/ExteriorDiskFlow.lean\#L57}{source
  line 57}.
\item
  {\small\texttt{GC.\allowbreak{}Long\allowbreak{}Time.\allowbreak{}exterior\allowbreak{}Disk\allowbreak{}Area\_\allowbreak{}nonneg}}\\
  Written proof; no skeleton admission. \hyperref[ha-5-4]{Explanation:
  HA-5.4};
  \href{https://github.com/qinz1yang/differential-geometry-dev/blob/ea0fae60ee01ef8c8d9b57a51794c6538f679c5d/DifferentialGeometry/Geometry/Flow/RicciFlow/LongTime/ExteriorDiskFlow.lean\#L62}{source
  line 62}.
\item
  {\small\texttt{GC.\allowbreak{}Long\allowbreak{}Time.\allowbreak{}exterior\allowbreak{}Disk\allowbreak{}Area\_\allowbreak{}eq}}\\
  Written proof; no skeleton admission. \hyperref[ha-5-4]{Explanation:
  HA-5.4};
  \href{https://github.com/qinz1yang/differential-geometry-dev/blob/ea0fae60ee01ef8c8d9b57a51794c6538f679c5d/DifferentialGeometry/Geometry/Flow/RicciFlow/LongTime/ExteriorDiskFlow.lean\#L68}{source
  line 68}.
\end{itemize}

\Needspace{6\baselineskip}

\subsection{M12.
LateDecomposition.lean}\label{m12.-latedecomposition.lean}

\textbf{File:}
{\small\texttt{Geometry/\allowbreak{}Flow/\allowbreak{}Ricci\allowbreak{}Flow/\allowbreak{}Long\allowbreak{}Time/\allowbreak{}Late\allowbreak{}Decomposition.\allowbreak{}lean}}

\begin{itemize}
\item
  {\small\texttt{GC.\allowbreak{}Long\allowbreak{}Time.\allowbreak{}has\allowbreak{}Common\allowbreak{}Neck\allowbreak{}Accuracy}}\\
  Definition/construction. \hyperref[f4-common-accuracy]{Explanation:
  F4};
  \href{https://github.com/qinz1yang/differential-geometry-dev/blob/ea0fae60ee01ef8c8d9b57a51794c6538f679c5d/DifferentialGeometry/Geometry/Flow/RicciFlow/LongTime/LateDecomposition.lean\#L17}{source
  line 17}.
\item
  {\small\texttt{GC.\allowbreak{}Long\allowbreak{}Time.\allowbreak{}Regular\allowbreak{}Slice.\allowbreak{}component\allowbreak{}Metric}}\\
  Definition/construction. \hyperref[f2-slices]{Explanation: F2};
  \href{https://github.com/qinz1yang/differential-geometry-dev/blob/ea0fae60ee01ef8c8d9b57a51794c6538f679c5d/DifferentialGeometry/Geometry/Flow/RicciFlow/LongTime/LateDecomposition.lean\#L23}{source
  line 23}.
\item
  {\small\texttt{GC.\allowbreak{}Long\allowbreak{}Time.\allowbreak{}has\allowbreak{}Exterior\allowbreak{}Area\allowbreak{}Obstruction\allowbreak{}After}}\\
  Definition/construction. \hyperref[f5-area-obstruction]{Explanation:
  F5};
  \href{https://github.com/qinz1yang/differential-geometry-dev/blob/ea0fae60ee01ef8c8d9b57a51794c6538f679c5d/DifferentialGeometry/Geometry/Flow/RicciFlow/LongTime/LateDecomposition.lean\#L29}{source
  line 29}.
\item
  {\small\texttt{GC.\allowbreak{}Long\allowbreak{}Time.\allowbreak{}has\allowbreak{}Late\allowbreak{}Sequence\allowbreak{}Tests}}\\
  Definition/construction. \hyperref[f6-sequence-tests]{Explanation:
  F6};
  \href{https://github.com/qinz1yang/differential-geometry-dev/blob/ea0fae60ee01ef8c8d9b57a51794c6538f679c5d/DifferentialGeometry/Geometry/Flow/RicciFlow/LongTime/LateDecomposition.lean\#L47}{source
  line 47}.
\item
  {\small\texttt{GC.\allowbreak{}Long\allowbreak{}Time.\allowbreak{}exists\_\allowbreak{}surgery\_\allowbreak{}with\_\allowbreak{}late\_\allowbreak{}sequence\_\allowbreak{}tests}}\\
  Admitted. \hyperref[f7-selected-flow]{Explanation: F7};
  \href{https://github.com/qinz1yang/differential-geometry-dev/blob/ea0fae60ee01ef8c8d9b57a51794c6538f679c5d/DifferentialGeometry/Geometry/Flow/RicciFlow/LongTime/LateDecomposition.lean\#L62}{source
  line 62}.
\item
  {\small\texttt{GC.\allowbreak{}Long\allowbreak{}Time.\allowbreak{}components\_\allowbreak{}geometrize\_\allowbreak{}of\_\allowbreak{}late\_\allowbreak{}sequence\_\allowbreak{}tests}}\\
  Written proof; admitted dependence.
  \hyperref[f8-bad-sequence-proof]{Explanation: F8};
  \href{https://github.com/qinz1yang/differential-geometry-dev/blob/ea0fae60ee01ef8c8d9b57a51794c6538f679c5d/DifferentialGeometry/Geometry/Flow/RicciFlow/LongTime/LateDecomposition.lean\#L71}{source
  line 71}.
\item
  {\small\texttt{GC.\allowbreak{}Long\allowbreak{}Time.\allowbreak{}geometrizes\_\allowbreak{}of\_\allowbreak{}metric}}\\
  Written proof; admitted dependence.
  \hyperref[f9-capstone]{Explanation: F9};
  \href{https://github.com/qinz1yang/differential-geometry-dev/blob/ea0fae60ee01ef8c8d9b57a51794c6538f679c5d/DifferentialGeometry/Geometry/Flow/RicciFlow/LongTime/LateDecomposition.lean\#L99}{source
  line 99}.
\end{itemize}

\Needspace{6\baselineskip}

\subsection{M13. RegularSlice.lean}\label{m13.-regularslice.lean}

\textbf{File:}
{\small\texttt{Geometry/\allowbreak{}Flow/\allowbreak{}Ricci\allowbreak{}Flow/\allowbreak{}Long\allowbreak{}Time/\allowbreak{}Regular\allowbreak{}Slice.\allowbreak{}lean}}

\begin{itemize}
\item
  {\small\texttt{GC.\allowbreak{}Long\allowbreak{}Time.\allowbreak{}Regular\allowbreak{}Slice}}\\
  Definition/construction. \hyperref[f2-slices]{Explanation: F2};
  \href{https://github.com/qinz1yang/differential-geometry-dev/blob/ea0fae60ee01ef8c8d9b57a51794c6538f679c5d/DifferentialGeometry/Geometry/Flow/RicciFlow/LongTime/RegularSlice.lean\#L13}{source
  line 13}.
\item
  {\small\texttt{GC.\allowbreak{}Long\allowbreak{}Time.\allowbreak{}Regular\allowbreak{}Slice.\allowbreak{}history}}\\
  Definition/construction. \hyperref[f2-slices]{Explanation: F2};
  \href{https://github.com/qinz1yang/differential-geometry-dev/blob/ea0fae60ee01ef8c8d9b57a51794c6538f679c5d/DifferentialGeometry/Geometry/Flow/RicciFlow/LongTime/RegularSlice.lean\#L24}{source
  line 24}.
\item
  {\small\texttt{GC.\allowbreak{}Long\allowbreak{}Time.\allowbreak{}Regular\allowbreak{}Slice.\allowbreak{}stage}}\\
  Definition/construction. \hyperref[f2-slices]{Explanation: F2};
  \href{https://github.com/qinz1yang/differential-geometry-dev/blob/ea0fae60ee01ef8c8d9b57a51794c6538f679c5d/DifferentialGeometry/Geometry/Flow/RicciFlow/LongTime/RegularSlice.lean\#L26}{source
  line 26}.
\item
  {\small\texttt{GC.\allowbreak{}Long\allowbreak{}Time.\allowbreak{}Regular\allowbreak{}Slice.\allowbreak{}metric}}\\
  Definition/construction. \hyperref[f2-slices]{Explanation: F2};
  \href{https://github.com/qinz1yang/differential-geometry-dev/blob/ea0fae60ee01ef8c8d9b57a51794c6538f679c5d/DifferentialGeometry/Geometry/Flow/RicciFlow/LongTime/RegularSlice.lean\#L29}{source
  line 29}.
\item
  {\small\texttt{GC.\allowbreak{}Long\allowbreak{}Time.\allowbreak{}Regular\allowbreak{}Slice.\allowbreak{}normalized\allowbreak{}Metric}}\\
  Definition/construction. \hyperref[f2-slices]{Explanation: F2};
  \href{https://github.com/qinz1yang/differential-geometry-dev/blob/ea0fae60ee01ef8c8d9b57a51794c6538f679c5d/DifferentialGeometry/Geometry/Flow/RicciFlow/LongTime/RegularSlice.lean\#L32}{source
  line 32}.
\item
  {\small\texttt{GC.\allowbreak{}Long\allowbreak{}Time.\allowbreak{}Regular\allowbreak{}Slice.\allowbreak{}curvature\allowbreak{}One\allowbreak{}Metric}}\\
  Definition/construction. \hyperref[f2-slices]{Explanation: F2};
  \href{https://github.com/qinz1yang/differential-geometry-dev/blob/ea0fae60ee01ef8c8d9b57a51794c6538f679c5d/DifferentialGeometry/Geometry/Flow/RicciFlow/LongTime/RegularSlice.lean\#L35}{source
  line 35}.
\item
  {\small\texttt{GC.\allowbreak{}Long\allowbreak{}Time.\allowbreak{}Regular\allowbreak{}Slice.\allowbreak{}initial}}\\
  Definition/construction. \hyperref[f2-slices]{Explanation: F2};
  \href{https://github.com/qinz1yang/differential-geometry-dev/blob/ea0fae60ee01ef8c8d9b57a51794c6538f679c5d/DifferentialGeometry/Geometry/Flow/RicciFlow/LongTime/RegularSlice.lean\#L38}{source
  line 38}.
\item
  {\small\texttt{GC.\allowbreak{}Long\allowbreak{}Time.\allowbreak{}has\allowbreak{}Arbitrarily\allowbreak{}Late\allowbreak{}Nonempty\allowbreak{}Slices}}\\
  Definition/construction. \hyperref[f2-slices]{Explanation: F2};
  \href{https://github.com/qinz1yang/differential-geometry-dev/blob/ea0fae60ee01ef8c8d9b57a51794c6538f679c5d/DifferentialGeometry/Geometry/Flow/RicciFlow/LongTime/RegularSlice.lean\#L43}{source
  line 43}.
\item
  {\small\texttt{GC.\allowbreak{}Long\allowbreak{}Time.\allowbreak{}exists\_\allowbreak{}regular\_\allowbreak{}slice\_\allowbreak{}after}}\\
  Written proof; no skeleton admission.
  \hyperref[f2-slices]{Explanation: F2};
  \href{https://github.com/qinz1yang/differential-geometry-dev/blob/ea0fae60ee01ef8c8d9b57a51794c6538f679c5d/DifferentialGeometry/Geometry/Flow/RicciFlow/LongTime/RegularSlice.lean\#L47}{source
  line 47}.
\item
  {\small\texttt{GC.\allowbreak{}Long\allowbreak{}Time.\allowbreak{}empty\_\allowbreak{}slice\_\allowbreak{}or\_\allowbreak{}arbitrarily\_\allowbreak{}late}}\\
  Written proof; no skeleton admission.
  \hyperref[f2-slices]{Explanation: F2};
  \href{https://github.com/qinz1yang/differential-geometry-dev/blob/ea0fae60ee01ef8c8d9b57a51794c6538f679c5d/DifferentialGeometry/Geometry/Flow/RicciFlow/LongTime/RegularSlice.lean\#L53}{source
  line 53}.
\item
  {\small\texttt{GC.\allowbreak{}Long\allowbreak{}Time.\allowbreak{}exists\_\allowbreak{}common\_\allowbreak{}late\_\allowbreak{}slice}}\\
  Written proof; no skeleton admission.
  \hyperref[f2-slices]{Explanation: F2};
  \href{https://github.com/qinz1yang/differential-geometry-dev/blob/ea0fae60ee01ef8c8d9b57a51794c6538f679c5d/DifferentialGeometry/Geometry/Flow/RicciFlow/LongTime/RegularSlice.lean\#L67}{source
  line 67}.
\item
  {\small\texttt{GC.\allowbreak{}Long\allowbreak{}Time.\allowbreak{}geometrizes\_\allowbreak{}of\_\allowbreak{}late\_\allowbreak{}slice\_\allowbreak{}supply}}\\
  Written proof; no skeleton admission.
  \hyperref[f2-slices]{Explanation: F2};
  \href{https://github.com/qinz1yang/differential-geometry-dev/blob/ea0fae60ee01ef8c8d9b57a51794c6538f679c5d/DifferentialGeometry/Geometry/Flow/RicciFlow/LongTime/RegularSlice.lean\#L91}{source
  line 91}.
\end{itemize}

\Needspace{6\baselineskip}

\subsection{M14. Cusp.lean}\label{m14.-cusp.lean}

\textbf{File:}
{\small\texttt{Geometry/\allowbreak{}Hyperbolic/\allowbreak{}Cusp.\allowbreak{}lean}}

\begin{itemize}
\item
  {\small\texttt{Differential\allowbreak{}Geometry.\allowbreak{}Geometry.\allowbreak{}Hyperbolic.\allowbreak{}Cusp\allowbreak{}Half\allowbreak{}Space}}\\
  Definition/construction. \hyperref[collapse-cuspmodel]{Explanation:
  C3};
  \href{https://github.com/qinz1yang/differential-geometry-dev/blob/ea0fae60ee01ef8c8d9b57a51794c6538f679c5d/DifferentialGeometry/Geometry/Hyperbolic/Cusp.lean\#L10}{source
  line 10}.
\item
  {\small\texttt{Differential\allowbreak{}Geometry.\allowbreak{}Geometry.\allowbreak{}Hyperbolic.\allowbreak{}cusp\allowbreak{}Domain}}\\
  Definition/construction. \hyperref[collapse-cuspmodel]{Explanation:
  C3};
  \href{https://github.com/qinz1yang/differential-geometry-dev/blob/ea0fae60ee01ef8c8d9b57a51794c6538f679c5d/DifferentialGeometry/Geometry/Hyperbolic/Cusp.lean\#L12}{source
  line 12}.
\item
  {\small\texttt{Differential\allowbreak{}Geometry.\allowbreak{}Geometry.\allowbreak{}Hyperbolic.\allowbreak{}Hyperbolic\allowbreak{}Cusp}}\\
  Definition/construction. \hyperref[collapse-cuspmodel]{Explanation:
  C3};
  \href{https://github.com/qinz1yang/differential-geometry-dev/blob/ea0fae60ee01ef8c8d9b57a51794c6538f679c5d/DifferentialGeometry/Geometry/Hyperbolic/Cusp.lean\#L14}{source
  line 14}.
\item
  {\small\texttt{Differential\allowbreak{}Geometry.\allowbreak{}Geometry.\allowbreak{}Hyperbolic.\allowbreak{}cusp\_\allowbreak{}constant\_\allowbreak{}sectional\_\allowbreak{}curvature}}\\
  Admitted. \hyperref[collapse-cuspmodel]{Explanation: C3};
  \href{https://github.com/qinz1yang/differential-geometry-dev/blob/ea0fae60ee01ef8c8d9b57a51794c6538f679c5d/DifferentialGeometry/Geometry/Hyperbolic/Cusp.lean\#L23}{source
  line 23}.
\end{itemize}

\Needspace{6\baselineskip}

\subsection{M15. ModelAtlas.lean}\label{m15.-modelatlas.lean}

\textbf{File:}
{\small\texttt{Geometry/\allowbreak{}Hyperbolic/\allowbreak{}Model\allowbreak{}Atlas.\allowbreak{}lean}}

\begin{itemize}
\item
  {\small\texttt{Differential\allowbreak{}Geometry.\allowbreak{}Geometry.\allowbreak{}Hyperbolic.\allowbreak{}has\_\allowbreak{}hyperbolic\_\allowbreak{}atlas\_\allowbreak{}of\_\allowbreak{}curvature\_\allowbreak{}neg\_\allowbreak{}one}}\\
  Admitted. \hyperref[ha-1-3]{Explanation: HA-1.3};
  \href{https://github.com/qinz1yang/differential-geometry-dev/blob/ea0fae60ee01ef8c8d9b57a51794c6538f679c5d/DifferentialGeometry/Geometry/Hyperbolic/ModelAtlas.lean\#L17}{source
  line 17}.
\item
  {\small\texttt{Differential\allowbreak{}Geometry.\allowbreak{}Geometry.\allowbreak{}Hyperbolic.\allowbreak{}finite\allowbreak{}Volume\allowbreak{}Geometric\allowbreak{}Structure}}\\
  Definition/construction; admitted dependence.
  \hyperref[ha-1-4]{Explanation: HA-1.4};
  \href{https://github.com/qinz1yang/differential-geometry-dev/blob/ea0fae60ee01ef8c8d9b57a51794c6538f679c5d/DifferentialGeometry/Geometry/Hyperbolic/ModelAtlas.lean\#L23}{source
  line 23}.
\item
  {\small\texttt{Differential\allowbreak{}Geometry.\allowbreak{}Geometry.\allowbreak{}Hyperbolic.\allowbreak{}finite\allowbreak{}Volume\allowbreak{}Geometric\allowbreak{}Structure\_\allowbreak{}model}}\\
  Written proof; admitted dependence. \hyperref[ha-1-4]{Explanation:
  HA-1.4};
  \href{https://github.com/qinz1yang/differential-geometry-dev/blob/ea0fae60ee01ef8c8d9b57a51794c6538f679c5d/DifferentialGeometry/Geometry/Hyperbolic/ModelAtlas.lean\#L42}{source
  line 42}.
\end{itemize}

\Needspace{6\baselineskip}

\subsection{M16. Rigidity.lean}\label{m16.-rigidity.lean}

\textbf{File:}
{\small\texttt{Geometry/\allowbreak{}Hyperbolic/\allowbreak{}Rigidity.\allowbreak{}lean}}

\begin{itemize}
\item
  {\small\texttt{Differential\allowbreak{}Geometry.\allowbreak{}Geometry.\allowbreak{}Hyperbolic.\allowbreak{}has\allowbreak{}Constant\allowbreak{}Sectional\allowbreak{}Curvature}}\\
  Definition/construction. \hyperref[ha-1-1]{Explanation: HA-1.1};
  \href{https://github.com/qinz1yang/differential-geometry-dev/blob/ea0fae60ee01ef8c8d9b57a51794c6538f679c5d/DifferentialGeometry/Geometry/Hyperbolic/Rigidity.lean\#L20}{source
  line 20}.
\item
  {\small\texttt{Differential\allowbreak{}Geometry.\allowbreak{}Geometry.\allowbreak{}Hyperbolic.\allowbreak{}mostow\_\allowbreak{}prasad}}\\
  Admitted. \hyperref[ha-1-2]{Explanation: HA-1.2};
  \href{https://github.com/qinz1yang/differential-geometry-dev/blob/ea0fae60ee01ef8c8d9b57a51794c6538f679c5d/DifferentialGeometry/Geometry/Hyperbolic/Rigidity.lean\#L24}{source
  line 24}.
\end{itemize}

\Needspace{6\baselineskip}

\subsection{M17.
FiniteRegularity.lean}\label{m17.-finiteregularity.lean}

\textbf{File:}
{\small\texttt{Geometry/\allowbreak{}Metric/\allowbreak{}Pullback/\allowbreak{}Finite\allowbreak{}Regularity.\allowbreak{}lean}}

\begin{itemize}
\item
  {\small\texttt{Differential\allowbreak{}Geometry.\allowbreak{}Geometry.\allowbreak{}exists\_\allowbreak{}pullback\_\allowbreak{}metric\_\allowbreak{}of\_\allowbreak{}finite\_\allowbreak{}diffeomorph}}\\
  Admitted. \hyperref[fr2]{Explanation: FR2};
  \href{https://github.com/qinz1yang/differential-geometry-dev/blob/ea0fae60ee01ef8c8d9b57a51794c6538f679c5d/DifferentialGeometry/Geometry/Metric/Pullback/FiniteRegularity.lean\#L16}{source
  line 16}.
\item
  {\small\texttt{Differential\allowbreak{}Geometry.\allowbreak{}Geometry.\allowbreak{}exists\_\allowbreak{}pullback\_\allowbreak{}metric\_\allowbreak{}of\_\allowbreak{}diffeomorph\_\allowbreak{}one\_\allowbreak{}order\_\allowbreak{}higher}}\\
  Written proof; admitted dependence. \hyperref[fr3]{Explanation: FR3};
  \href{https://github.com/qinz1yang/differential-geometry-dev/blob/ea0fae60ee01ef8c8d9b57a51794c6538f679c5d/DifferentialGeometry/Geometry/Metric/Pullback/FiniteRegularity.lean\#L26}{source
  line 26}.
\item
  {\small\texttt{Differential\allowbreak{}Geometry.\allowbreak{}Geometry.\allowbreak{}exists\_\allowbreak{}pullback\_\allowbreak{}metric\_\allowbreak{}of\_\allowbreak{}diffeomorph\_\allowbreak{}three\_\allowbreak{}orders\_\allowbreak{}lower}}\\
  Written proof; admitted dependence. \hyperref[fr3]{Explanation: FR3};
  \href{https://github.com/qinz1yang/differential-geometry-dev/blob/ea0fae60ee01ef8c8d9b57a51794c6538f679c5d/DifferentialGeometry/Geometry/Metric/Pullback/FiniteRegularity.lean\#L35}{source
  line 35}.
\end{itemize}

\Needspace{6\baselineskip}

\subsection{M18.
ExteriorDiskArea.lean}\label{m18.-exteriordiskarea.lean}

\textbf{File:}
{\small\texttt{Geometry/\allowbreak{}Minimal\allowbreak{}Surface/\allowbreak{}Exterior\allowbreak{}Disk\allowbreak{}Area.\allowbreak{}lean}}

\begin{itemize}
\item
  {\small\texttt{Differential\allowbreak{}Geometry.\allowbreak{}Geometry.\allowbreak{}Minimal\allowbreak{}Surface.\allowbreak{}is\allowbreak{}Exterior\allowbreak{}Spanning\allowbreak{}Disk}}\\
  Definition/construction. \hyperref[ha-3-1]{Explanation: HA-3.1};
  \href{https://github.com/qinz1yang/differential-geometry-dev/blob/ea0fae60ee01ef8c8d9b57a51794c6538f679c5d/DifferentialGeometry/Geometry/MinimalSurface/ExteriorDiskArea.lean\#L15}{source
  line 15}.
\item
  {\small\texttt{Differential\allowbreak{}Geometry.\allowbreak{}Geometry.\allowbreak{}Minimal\allowbreak{}Surface.\allowbreak{}least\allowbreak{}Exterior\allowbreak{}Disk\allowbreak{}Area}}\\
  Definition/construction. \hyperref[ha-3-2]{Explanation: HA-3.2};
  \href{https://github.com/qinz1yang/differential-geometry-dev/blob/ea0fae60ee01ef8c8d9b57a51794c6538f679c5d/DifferentialGeometry/Geometry/MinimalSurface/ExteriorDiskArea.lean\#L22}{source
  line 22}.
\item
  {\small\texttt{Differential\allowbreak{}Geometry.\allowbreak{}Geometry.\allowbreak{}Minimal\allowbreak{}Surface.\allowbreak{}least\allowbreak{}Exterior\allowbreak{}Disk\allowbreak{}Area\_\allowbreak{}nonneg}}\\
  Written proof; no skeleton admission. \hyperref[ha-3-2]{Explanation:
  HA-3.2};
  \href{https://github.com/qinz1yang/differential-geometry-dev/blob/ea0fae60ee01ef8c8d9b57a51794c6538f679c5d/DifferentialGeometry/Geometry/MinimalSurface/ExteriorDiskArea.lean\#L27}{source
  line 27}.
\item
  {\small\texttt{Differential\allowbreak{}Geometry.\allowbreak{}Geometry.\allowbreak{}Minimal\allowbreak{}Surface.\allowbreak{}least\allowbreak{}Exterior\allowbreak{}Disk\allowbreak{}Area\_\allowbreak{}le}}\\
  Written proof; no skeleton admission. \hyperref[ha-3-2]{Explanation:
  HA-3.2};
  \href{https://github.com/qinz1yang/differential-geometry-dev/blob/ea0fae60ee01ef8c8d9b57a51794c6538f679c5d/DifferentialGeometry/Geometry/MinimalSurface/ExteriorDiskArea.lean\#L33}{source
  line 33}.
\item
  {\small\texttt{Differential\allowbreak{}Geometry.\allowbreak{}Geometry.\allowbreak{}Minimal\allowbreak{}Surface.\allowbreak{}exterior\allowbreak{}Disk\allowbreak{}Area\allowbreak{}On}}\\
  Definition/construction. \hyperref[ha-3-3]{Explanation: HA-3.3};
  \href{https://github.com/qinz1yang/differential-geometry-dev/blob/ea0fae60ee01ef8c8d9b57a51794c6538f679c5d/DifferentialGeometry/Geometry/MinimalSurface/ExteriorDiskArea.lean\#L43}{source
  line 43}.
\item
  {\small\texttt{Differential\allowbreak{}Geometry.\allowbreak{}Geometry.\allowbreak{}Minimal\allowbreak{}Surface.\allowbreak{}exterior\allowbreak{}Disk\allowbreak{}Area\allowbreak{}On\_\allowbreak{}nonneg}}\\
  Written proof; no skeleton admission. \hyperref[ha-3-3]{Explanation:
  HA-3.3};
  \href{https://github.com/qinz1yang/differential-geometry-dev/blob/ea0fae60ee01ef8c8d9b57a51794c6538f679c5d/DifferentialGeometry/Geometry/MinimalSurface/ExteriorDiskArea.lean\#L52}{source
  line 52}.
\item
  {\small\texttt{Differential\allowbreak{}Geometry.\allowbreak{}Geometry.\allowbreak{}Minimal\allowbreak{}Surface.\allowbreak{}continuous\allowbreak{}On\_\allowbreak{}least\allowbreak{}Exterior\allowbreak{}Disk\allowbreak{}Area\_\allowbreak{}of\_\allowbreak{}local\_\allowbreak{}disk\_\allowbreak{}comparisons}}\\
  Admitted. \hyperref[ha-3-4]{Explanation: HA-3.4};
  \href{https://github.com/qinz1yang/differential-geometry-dev/blob/ea0fae60ee01ef8c8d9b57a51794c6538f679c5d/DifferentialGeometry/Geometry/MinimalSurface/ExteriorDiskArea.lean\#L65}{source
  line 65}.
\end{itemize}

\Needspace{6\baselineskip}

\subsection{M19. Theorem.lean}\label{m19.-theorem.lean}

\textbf{File:}
{\small\texttt{Topology/\allowbreak{}Three\allowbreak{}Manifold/\allowbreak{}Geometrization/\allowbreak{}Theorem.\allowbreak{}lean}}

\begin{itemize}
\item
  {\small\texttt{GC.\allowbreak{}Endpoint.\allowbreak{}geometrization}}\\
  Written proof; admitted dependence.
  \hyperref[f9-capstone]{Explanation: F9};
  \href{https://github.com/qinz1yang/differential-geometry-dev/blob/ea0fae60ee01ef8c8d9b57a51794c6538f679c5d/DifferentialGeometry/Topology/ThreeManifold/Geometrization/Theorem.lean\#L10}{source
  line 10}.
\item
  {\small\texttt{GC.\allowbreak{}Endpoint.\allowbreak{}geometrization\_\allowbreak{}conjecture}}\\
  Written proof; admitted dependence.
  \hyperref[f9-capstone]{Explanation: F9};
  \href{https://github.com/qinz1yang/differential-geometry-dev/blob/ea0fae60ee01ef8c8d9b57a51794c6538f679c5d/DifferentialGeometry/Topology/ThreeManifold/Geometrization/Theorem.lean\#L14}{source
  line 14}.
\item
  {\small\texttt{GC.\allowbreak{}Endpoint.\allowbreak{}smooth\_\allowbreak{}geometrization\_\allowbreak{}conjecture}}\\
  Written proof; admitted dependence.
  \hyperref[f9-capstone]{Explanation: F9};
  \href{https://github.com/qinz1yang/differential-geometry-dev/blob/ea0fae60ee01ef8c8d9b57a51794c6538f679c5d/DifferentialGeometry/Topology/ThreeManifold/Geometrization/Theorem.lean\#L16}{source
  line 16}.
\end{itemize}

\Needspace{6\baselineskip}

\subsection{M20. Presentation.lean}\label{m20.-presentation.lean}

\textbf{File:}
{\small\texttt{Topology/\allowbreak{}Three\allowbreak{}Manifold/\allowbreak{}Graph\allowbreak{}Manifold/\allowbreak{}Presentation.\allowbreak{}lean}}

\begin{itemize}
\item
  {\small\texttt{GC.\allowbreak{}Graph\allowbreak{}Manifold.\allowbreak{}Surface\allowbreak{}Model}}\\
  Definition/construction. \hyperref[t1]{Explanation: T1};
  \href{https://github.com/qinz1yang/differential-geometry-dev/blob/ea0fae60ee01ef8c8d9b57a51794c6538f679c5d/DifferentialGeometry/Topology/ThreeManifold/GraphManifold/Presentation.lean\#L10}{source
  line 10}.
\item
  {\small\texttt{GC.\allowbreak{}Graph\allowbreak{}Manifold.\allowbreak{}Surface\allowbreak{}Model.\allowbreak{}Space}}\\
  Definition/construction. \hyperref[t1]{Explanation: T1};
  \href{https://github.com/qinz1yang/differential-geometry-dev/blob/ea0fae60ee01ef8c8d9b57a51794c6538f679c5d/DifferentialGeometry/Topology/ThreeManifold/GraphManifold/Presentation.lean\#L13}{source
  line 13}.
\item
  {\small\texttt{GC.\allowbreak{}Graph\allowbreak{}Manifold.\allowbreak{}Surface\allowbreak{}Model.\allowbreak{}model}}\\
  Definition/construction. \hyperref[t1]{Explanation: T1};
  \href{https://github.com/qinz1yang/differential-geometry-dev/blob/ea0fae60ee01ef8c8d9b57a51794c6538f679c5d/DifferentialGeometry/Topology/ThreeManifold/GraphManifold/Presentation.lean\#L22}{source
  line 22}.
\item
  {\small\texttt{GC.\allowbreak{}Graph\allowbreak{}Manifold.\allowbreak{}Compact\allowbreak{}Surface}}\\
  Definition/construction. \hyperref[t1]{Explanation: T1};
  \href{https://github.com/qinz1yang/differential-geometry-dev/blob/ea0fae60ee01ef8c8d9b57a51794c6538f679c5d/DifferentialGeometry/Topology/ThreeManifold/GraphManifold/Presentation.lean\#L28}{source
  line 28}.
\item
  {\small\texttt{GC.\allowbreak{}Graph\allowbreak{}Manifold.\allowbreak{}Circle\allowbreak{}Fibration}}\\
  Definition/construction. \hyperref[t1]{Explanation: T1};
  \href{https://github.com/qinz1yang/differential-geometry-dev/blob/ea0fae60ee01ef8c8d9b57a51794c6538f679c5d/DifferentialGeometry/Topology/ThreeManifold/GraphManifold/Presentation.lean\#L43}{source
  line 43}.
\item
  {\small\texttt{GC.\allowbreak{}Graph\allowbreak{}Manifold.\allowbreak{}Boundary\allowbreak{}Tori}}\\
  Definition/construction. \hyperref[t2]{Explanation: T2};
  \href{https://github.com/qinz1yang/differential-geometry-dev/blob/ea0fae60ee01ef8c8d9b57a51794c6538f679c5d/DifferentialGeometry/Topology/ThreeManifold/GraphManifold/Presentation.lean\#L55}{source
  line 55}.
\item
  {\small\texttt{GC.\allowbreak{}Graph\allowbreak{}Manifold.\allowbreak{}Boundary\allowbreak{}Tori.\allowbreak{}torus\allowbreak{}Map}}\\
  Definition/construction. \hyperref[t2]{Explanation: T2};
  \href{https://github.com/qinz1yang/differential-geometry-dev/blob/ea0fae60ee01ef8c8d9b57a51794c6538f679c5d/DifferentialGeometry/Topology/ThreeManifold/GraphManifold/Presentation.lean\#L64}{source
  line 64}.
\item
  {\small\texttt{GC.\allowbreak{}Graph\allowbreak{}Manifold.\allowbreak{}Boundary\allowbreak{}Tori.\allowbreak{}image}}\\
  Definition/construction. \hyperref[t2]{Explanation: T2};
  \href{https://github.com/qinz1yang/differential-geometry-dev/blob/ea0fae60ee01ef8c8d9b57a51794c6538f679c5d/DifferentialGeometry/Topology/ThreeManifold/GraphManifold/Presentation.lean\#L67}{source
  line 67}.
\item
  {\small\texttt{GC.\allowbreak{}Graph\allowbreak{}Manifold.\allowbreak{}Boundary\allowbreak{}Tori.\allowbreak{}zero\_\allowbreak{}mem\_\allowbreak{}source}}\\
  Written proof; no skeleton admission. \hyperref[t2]{Explanation: T2};
  \href{https://github.com/qinz1yang/differential-geometry-dev/blob/ea0fae60ee01ef8c8d9b57a51794c6538f679c5d/DifferentialGeometry/Topology/ThreeManifold/GraphManifold/Presentation.lean\#L70}{source
  line 70}.
\item
  {\small\texttt{GC.\allowbreak{}Graph\allowbreak{}Manifold.\allowbreak{}Boundary\allowbreak{}Tori.\allowbreak{}torus\allowbreak{}Map\_\allowbreak{}smooth}}\\
  Written proof; no skeleton admission. \hyperref[t2]{Explanation: T2};
  \href{https://github.com/qinz1yang/differential-geometry-dev/blob/ea0fae60ee01ef8c8d9b57a51794c6538f679c5d/DifferentialGeometry/Topology/ThreeManifold/GraphManifold/Presentation.lean\#L77}{source
  line 77}.
\item
  {\small\texttt{GC.\allowbreak{}Graph\allowbreak{}Manifold.\allowbreak{}Boundary\allowbreak{}Tori.\allowbreak{}boundary\allowbreak{}Map}}\\
  Definition/construction. \hyperref[t2]{Explanation: T2};
  \href{https://github.com/qinz1yang/differential-geometry-dev/blob/ea0fae60ee01ef8c8d9b57a51794c6538f679c5d/DifferentialGeometry/Topology/ThreeManifold/GraphManifold/Presentation.lean\#L82}{source
  line 82}.
\item
  {\small\texttt{GC.\allowbreak{}Graph\allowbreak{}Manifold.\allowbreak{}Boundary\allowbreak{}Tori.\allowbreak{}torus\allowbreak{}Map\_\allowbreak{}is\allowbreak{}Embedding}}\\
  Written proof; no skeleton admission. \hyperref[t2]{Explanation: T2};
  \href{https://github.com/qinz1yang/differential-geometry-dev/blob/ea0fae60ee01ef8c8d9b57a51794c6538f679c5d/DifferentialGeometry/Topology/ThreeManifold/GraphManifold/Presentation.lean\#L86}{source
  line 86}.
\item
  {\small\texttt{GC.\allowbreak{}Graph\allowbreak{}Manifold.\allowbreak{}Boundary\allowbreak{}Tori.\allowbreak{}incompressible}}\\
  Definition/construction. \hyperref[t2]{Explanation: T2};
  \href{https://github.com/qinz1yang/differential-geometry-dev/blob/ea0fae60ee01ef8c8d9b57a51794c6538f679c5d/DifferentialGeometry/Topology/ThreeManifold/GraphManifold/Presentation.lean\#L93}{source
  line 93}.
\item
  {\small\texttt{GC.\allowbreak{}Graph\allowbreak{}Manifold.\allowbreak{}Torus\allowbreak{}Pairing}}\\
  Definition/construction. \hyperref[t3]{Explanation: T3};
  \href{https://github.com/qinz1yang/differential-geometry-dev/blob/ea0fae60ee01ef8c8d9b57a51794c6538f679c5d/DifferentialGeometry/Topology/ThreeManifold/GraphManifold/Presentation.lean\#L98}{source
  line 98}.
\item
  {\small\texttt{GC.\allowbreak{}Graph\allowbreak{}Manifold.\allowbreak{}Torus\allowbreak{}Pairing.\allowbreak{}Quotient\allowbreak{}Space}}\\
  Definition/construction. \hyperref[t3]{Explanation: T3};
  \href{https://github.com/qinz1yang/differential-geometry-dev/blob/ea0fae60ee01ef8c8d9b57a51794c6538f679c5d/DifferentialGeometry/Topology/ThreeManifold/GraphManifold/Presentation.lean\#L118}{source
  line 118}.
\item
  {\small\texttt{GC.\allowbreak{}Graph\allowbreak{}Manifold.\allowbreak{}Torus\allowbreak{}Pairing.\allowbreak{}quotient\allowbreak{}Map}}\\
  Definition/construction. \hyperref[t3]{Explanation: T3};
  \href{https://github.com/qinz1yang/differential-geometry-dev/blob/ea0fae60ee01ef8c8d9b57a51794c6538f679c5d/DifferentialGeometry/Topology/ThreeManifold/GraphManifold/Presentation.lean\#L124}{source
  line 124}.
\item
  {\small\texttt{GC.\allowbreak{}Graph\allowbreak{}Manifold.\allowbreak{}Raw\allowbreak{}Graph\allowbreak{}Presentation}}\\
  Definition/construction. \hyperref[t4]{Explanation: T4};
  \href{https://github.com/qinz1yang/differential-geometry-dev/blob/ea0fae60ee01ef8c8d9b57a51794c6538f679c5d/DifferentialGeometry/Topology/ThreeManifold/GraphManifold/Presentation.lean\#L129}{source
  line 129}.
\end{itemize}

\Needspace{6\baselineskip}

\subsection{M21. Refinement.lean}\label{m21.-refinement.lean}

\textbf{File:}
{\small\texttt{Topology/\allowbreak{}Three\allowbreak{}Manifold/\allowbreak{}Graph\allowbreak{}Manifold/\allowbreak{}Refinement.\allowbreak{}lean}}

\begin{itemize}
\item
  {\small\texttt{GC.\allowbreak{}Graph\allowbreak{}Manifold.\allowbreak{}is\allowbreak{}Hyperbolic\allowbreak{}Interior\allowbreak{}Geometry}}\\
  Definition/construction. \hyperref[t9]{Explanation: T9};
  \href{https://github.com/qinz1yang/differential-geometry-dev/blob/ea0fae60ee01ef8c8d9b57a51794c6538f679c5d/DifferentialGeometry/Topology/ThreeManifold/GraphManifold/Refinement.lean\#L11}{source
  line 11}.
\item
  {\small\texttt{GC.\allowbreak{}Graph\allowbreak{}Manifold.\allowbreak{}Hyperbolic\allowbreak{}Or\allowbreak{}Graph}}\\
  Definition/construction. \hyperref[t9]{Explanation: T9};
  \href{https://github.com/qinz1yang/differential-geometry-dev/blob/ea0fae60ee01ef8c8d9b57a51794c6538f679c5d/DifferentialGeometry/Topology/ThreeManifold/GraphManifold/Refinement.lean\#L17}{source
  line 17}.
\item
  {\small\texttt{GC.\allowbreak{}Graph\allowbreak{}Manifold.\allowbreak{}exists\_\allowbreak{}prime\_\allowbreak{}decomposition\_\allowbreak{}of\_\allowbreak{}raw\allowbreak{}Graph\allowbreak{}Presentation}}\\
  Admitted. \hyperref[t6]{Explanation: T6};
  \href{https://github.com/qinz1yang/differential-geometry-dev/blob/ea0fae60ee01ef8c8d9b57a51794c6538f679c5d/DifferentialGeometry/Topology/ThreeManifold/GraphManifold/Refinement.lean\#L23}{source
  line 23}.
\item
  {\small\texttt{GC.\allowbreak{}Graph\allowbreak{}Manifold.\allowbreak{}exists\_\allowbreak{}geometric\_\allowbreak{}decomposition\_\allowbreak{}of\_\allowbreak{}prime\_\allowbreak{}raw\allowbreak{}Graph\allowbreak{}Presentation}}\\
  Admitted. \hyperref[t7]{Explanation: T7};
  \href{https://github.com/qinz1yang/differential-geometry-dev/blob/ea0fae60ee01ef8c8d9b57a51794c6538f679c5d/DifferentialGeometry/Topology/ThreeManifold/GraphManifold/Refinement.lean\#L31}{source
  line 31}.
\item
  {\small\texttt{GC.\allowbreak{}Graph\allowbreak{}Manifold.\allowbreak{}geometrizes\_\allowbreak{}of\_\allowbreak{}raw\allowbreak{}Graph\allowbreak{}Presentation}}\\
  Written proof; admitted dependence. \hyperref[t8]{Explanation: T8};
  \href{https://github.com/qinz1yang/differential-geometry-dev/blob/ea0fae60ee01ef8c8d9b57a51794c6538f679c5d/DifferentialGeometry/Topology/ThreeManifold/GraphManifold/Refinement.lean\#L37}{source
  line 37}.
\item
  {\small\texttt{GC.\allowbreak{}Graph\allowbreak{}Manifold.\allowbreak{}exists\_\allowbreak{}prime\_\allowbreak{}geometric\_\allowbreak{}decomposition\_\allowbreak{}of\_\allowbreak{}hyperbolic\allowbreak{}Or\allowbreak{}Graph}}\\
  Admitted. \hyperref[t9]{Explanation: T9};
  \href{https://github.com/qinz1yang/differential-geometry-dev/blob/ea0fae60ee01ef8c8d9b57a51794c6538f679c5d/DifferentialGeometry/Topology/ThreeManifold/GraphManifold/Refinement.lean\#L46}{source
  line 46}.
\item
  {\small\texttt{GC.\allowbreak{}Graph\allowbreak{}Manifold.\allowbreak{}geometrizes\_\allowbreak{}of\_\allowbreak{}hyperbolic\allowbreak{}Or\allowbreak{}Graph}}\\
  Written proof; admitted dependence. \hyperref[t10]{Explanation: T10};
  \href{https://github.com/qinz1yang/differential-geometry-dev/blob/ea0fae60ee01ef8c8d9b57a51794c6538f679c5d/DifferentialGeometry/Topology/ThreeManifold/GraphManifold/Refinement.lean\#L54}{source
  line 54}.
\end{itemize}

\Needspace{6\baselineskip}

\subsection{M22. Decomposition.lean}\label{m22.-decomposition.lean}

\textbf{File:}
{\small\texttt{Topology/\allowbreak{}Three\allowbreak{}Manifold/\allowbreak{}Torus\allowbreak{}Cut/\allowbreak{}Decomposition.\allowbreak{}lean}}

\begin{itemize}
\item
  {\small\texttt{GC.\allowbreak{}Graph\allowbreak{}Manifold.\allowbreak{}component\allowbreak{}Carrier}}\\
  Definition/construction. \hyperref[t5]{Explanation: T5};
  \href{https://github.com/qinz1yang/differential-geometry-dev/blob/ea0fae60ee01ef8c8d9b57a51794c6538f679c5d/DifferentialGeometry/Topology/ThreeManifold/TorusCut/Decomposition.lean\#L10}{source
  line 10}.
\item
  {\small\texttt{GC.\allowbreak{}Graph\allowbreak{}Manifold.\allowbreak{}Torus\allowbreak{}Decomposition}}\\
  Definition/construction. \hyperref[t5]{Explanation: T5};
  \href{https://github.com/qinz1yang/differential-geometry-dev/blob/ea0fae60ee01ef8c8d9b57a51794c6538f679c5d/DifferentialGeometry/Topology/ThreeManifold/TorusCut/Decomposition.lean\#L17}{source
  line 17}.
\item
  {\small\texttt{GC.\allowbreak{}Graph\allowbreak{}Manifold.\allowbreak{}Torus\allowbreak{}Decomposition.\allowbreak{}component}}\\
  Definition/construction. \hyperref[t5]{Explanation: T5};
  \href{https://github.com/qinz1yang/differential-geometry-dev/blob/ea0fae60ee01ef8c8d9b57a51794c6538f679c5d/DifferentialGeometry/Topology/ThreeManifold/TorusCut/Decomposition.lean\#L30}{source
  line 30}.
\item
  {\small\texttt{GC.\allowbreak{}Graph\allowbreak{}Manifold.\allowbreak{}Torus\allowbreak{}Decomposition.\allowbreak{}to\allowbreak{}Geometric\allowbreak{}Decomposition}}\\
  Definition/construction. \hyperref[t5]{Explanation: T5};
  \href{https://github.com/qinz1yang/differential-geometry-dev/blob/ea0fae60ee01ef8c8d9b57a51794c6538f679c5d/DifferentialGeometry/Topology/ThreeManifold/TorusCut/Decomposition.lean\#L34}{source
  line 34}.
\end{itemize}

\end{document}
